% concatenated from main.tex
\documentclass[11pt]{article}
\usepackage{amsmath}
\usepackage[T1]{fontenc}
\usepackage[latin9]{inputenc}
\usepackage[a4paper]{geometry}
\usepackage{setspace}
\geometry{verbose,tmargin=0.75in,bmargin=0.75in,lmargin=0.75in,rmargin=0.75in}
\usepackage[affil-it]{authblk}
\usepackage[english]{babel}
\usepackage{blindtext}
\onehalfspacing
\usepackage{graphicx,psfrag,epsf}
\usepackage{amsmath,amssymb,mathtools,amsthm}
\usepackage[authoryear]{natbib}
\usepackage{mathtools}
\usepackage[unicode=true,
 bookmarks=false,
 breaklinks=false,pdfborder={0 0 1},colorlinks=true,citecolor = blue]{hyperref}
\usepackage{caption}
\usepackage{subcaption}
\usepackage{xcolor}
\usepackage{bm}
\usepackage{color}
\usepackage{tikz}
\usepackage{stmaryrd}
\usepackage{url}
\usepackage{mathrsfs}  
\usepackage{enumerate}
\usepackage{colortbl}
\usepackage{comment}
\usepackage{gensymb}
\usepackage[tableposition=top]{caption}
\RequirePackage{booktabs}
\usepackage{threeparttable}
\usepackage{array}

\newtheorem{assumption}{Assumption}
\newtheorem{example}{Example}
\newtheorem{definition}{Definition}
\newtheorem{corollary}{Corollary}
\newtheorem{property}{Property}
\newtheorem{proposition}{Proposition}
\newtheorem{theorem}{Theorem}
\newtheorem{remark}{Remark}
\newtheorem{defi}{Definition}
\newtheorem{thm}{Theorem}
\newtheorem{lemma}{Lemma}
\newcommand{\M}{\mathbf}
\newcommand{\lbt}{\llbracket}
\newcommand{\rbt}{\rrbracket}
\newcommand{\B}{\boldsymbol}
% \newcommand{\qed}{\hfill \ensuremath{\Box}}
\newcommand{\ind}{\mbox{$\perp\!\!\!\perp$}}
\newcommand{\nind}{\mbox{$\not\perp\!\!\!\perp$}}
\DeclareMathOperator{\sgn}{sgn}
\DeclareMathOperator{\diag}{diag}
\newcommand{\dd}{{\B\cdot}}
\newcommand{\od}{{\scriptstyle\star}}
 \newcommand\scalemath[2]{\scalebox{#1}{\mbox{\ensuremath{\displaystyle #2}}}}
\newcommand{\blind}{0}
\newcommand*{\QEDB}{\hfill\ensuremath{\square}}
\newcommand{\bee}{\begin{equation}\begin{aligned}}
	\newcommand{\ee}{\end{aligned}\end{equation}}
\newcommand{\beee}{\begin{eqnarray}}
	\newcommand{\eee}{\end{eqnarray}}
\newcommand{\indep}{\perp\!\!\!\perp}
\newcommand{\bone}{\mathbf{1}}
\newcommand{\red}{\color{black}}
\newcommand{\black}{\color{black}}
\newcommand{\Beta}{\text{Beta}}
\newcommand{\Binomial}{\text{Binomial}}
\newcommand{\Bern}{\text{Bern}}
\newcommand{\Expo}{\text{Expo}}
\newcommand{\Pois}{\text{Pois}}
\newcommand{\Unif}{\text{Unif}}
\newcommand{\Gammad}{\text{Gamma}}
\newcommand{\logit}{\text{logit}}
\newcommand{\expit}{\text{expit}}
\newcommand{\mexpit}{\text{mexpit}}
\newcommand{\pr}{\mathbb{P}}
\newcommand{\var}{\textnormal{var}}
\newcommand{\cov}{\textnormal{cov}}
\newcommand{\sumN}{\sum_{i=1}^N}
\newcommand{\nb}{\nonumber}
\newcommand{\E}{\mathbb{E}}
\newcommand{\LOTE}{\textup{LOTE}}
\newcommand{\lemmasone}{\textup{Lemma~\ref{lem:1}}}
\newcommand{\nn}{\nonumber}
\newcommand{\bZ}{\bm{Z}}
\newcommand{\bD}{\bm{D}}
\newcommand{\bz}{\bm{z}}
\newcommand{\bd}{\bm{d}}
\newcommand{\oD}{\overline{D}}
\newcommand{\oY}{\overline{Y}}
\newcommand{\bS}{\bar{s}}
\newcommand{\bs}{\bar{s}}
\newcommand{\bY}{\bm{Y}}
\newcommand{\bI}{\mathbf{I}}
\newcommand{\bM}{\mathbf{M}}
\newcommand{\bP}{\mathbf{P}}
\newcommand{\bQ}{\mathbf{Q}}
\newcommand{\bV}{\mathbf{V}}
\newcommand{\bU}{\mathbf{U}}
\newcommand{\bW}{\mathbf{W}}
\newcommand{\bepsilon}{\bm{\epsilon}}
\newcommand{\boldeta}{\bm{\eta}}
\newcommand{\ipw}{\textnormal{tp-ps}}
\newcommand{\ipsw}{\textnormal{sw}}
\newcommand{\reg}{\mathrm{reg}}
\newcommand{\regta}{\textnormal{tp-om}}
\newcommand{\regps}{\textnormal{ps-om}}
\newcommand{\aipw}{\mathrm{aw}}
\newcommand{\rmcr}{\mathrm{CR}}
\newcommand{\rmr}{\mathrm{R}}
\newcommand{\T}{\top}
\newcommand{\N}{\mathcal{N}}
\newcommand{\R}{\mathbb{R}}
\newcommand{\wtp}{\widetilde{p}}
\newcommand{\PM}{\mathrm{PM}_{2.5}}
\newcommand*\bigcdot{\mathpalette\bigcdot@{.5}}
\newcommand{\Mom}{\mathcal{M}_{\textnormal{om}}}
\newcommand{\Mps}{\mathcal{M}_{\textnormal{ps}}}
\newcommand{\Mta}{\mathcal{M}_{\textnormal{tp}}}
\newcommand{\Momps}{\mathcal{M}_{\textnormal{ps}+\textnormal{om}}}
\newcommand{\Momta}{\mathcal{M}_{\textnormal{tp}+\textnormal{om}}}
\newcommand{\Mpsta}{\mathcal{M}_{\textnormal{tp}+\textnormal{ps}}}
\newcommand{\Mompsta}{\mathcal{M}_{\textnormal{tp}+\textnormal{ps}+\textnormal{om} }}
\newtheorem{condition}{condition}

\newcommand{\norm}[1]{\left\lVert #1 \right\rVert}
\newcommand{\argmax}{\operatornamewithlimits{argmax}}
\newcommand{\argmin}{\operatornamewithlimits{argmin}}
\newcommand{\floor}[1]{\left\lfloor #1 \right\rfloor}
\newcommand{\ceil}[1]{\left\lceil #1 \right\rceil}
\newcommand{\proofend}{\hfill$\blacksquare$}

\usepackage{xr}
\makeatletter
\newcommand*{\addFileDependency}[1]{
  \typeout{(#1)}
  \@addtofilelist{#1}
  \IfFileExists{#1}{}{\typeout{No file #1.}}
}
\makeatother
\newcommand*{\myexternaldocument}[1]{%
    \externaldocument{#1}%
    \addFileDependency{#1.tex}%
    \addFileDependency{#1.aux}%
}
% \myexternaldocument{supp}

\makeatletter
\providecommand{\tabularnewline}{\\}
\@ifundefined{date}{}{\date{}}
\bibpunct{(}{)}{,}{autheryear}{,}{,}
\allowdisplaybreaks


\usepackage{babel}

\usepackage{babel}

\usepackage{babel}

\makeatother

\usepackage{babel} 

 
% Using \doublespacing in the preamble 
% changes the text to double-line spacing
% \onehalfspacing
% \renewcommand{\baselinestretch}{1.8} 
\begin{document}
\title{{\textbf{\Large{}{}{}{}Adversarial Estimation of Assortment Probabilities under Independence Structure}}}
\author{Alexandre Belloni\thanks{Fuqua School of Business, Duke University. Email: abn5@duke.edu} \thanks{Corresponding author.}, $\ $ Yan Chen\thanks{Fuqua School of Business, Duke University. Email: yc555@duke.edu} $\ $ and $\ $ Matthew Harding\thanks{Department of Economics,
University of California, Irvine. Email: matthew.harding@uci.edu}}
\maketitle

\vspace{-1.0cm}
\onehalfspacing
\begin{abstract}
We consider the problem of estimating assortment probabilities, which is common in operations management applications, including product bundling, advertising, etc. Existing approaches typically model each assortment as a category and apply multinomial models to estimate the choice probabilities; while computationally convenient, these methods do not exploit independence structures in the joint distribution and may therefore be statistically inefficient when the total number of items is large. Using the representation from \citet{bahadur1959representation}, we relate the sparsity of the generalized correlation coefficients to the independence structure of the binary components. We formulate the problem as estimating a high-dimensional vector of generalized correlation coefficients, together with low or moderate-dimensional nuisance parameters corresponding to the marginal probabilities. We develop a regularized adversarial estimator that attains the optimal rate under standard regularity conditions while remaining computationally feasible. The framework naturally extends to settings with covariates. We apply the proposed estimators to causal inference with multiple binary treatments and show substantial finite-sample improvements over non-adaptive methods. Numerical studies corroborate the theoretical results. 
\end{abstract}
\noindent%
{\it Keywords:} High-dimensional parameter of interest; moderate/low-dimensional nuisance parameter; localized regularized adversarial estimators; many binary components. 
\vfill

\onehalfspacing
% Text of your paper here
\section{Introduction}\label{sec-intro}
Many empirical applications involve estimating assortment probabilities, such as product bundle purchases \citep[e.g.][]{chung2003general,gabel2022product,blanchet2016markov,talluri2004revenue}, clinical trials \citep[e.g.][]{gueorguieva2001correlated,agresti2000strategies}, advertising \citep[e.g.][]{aouad2021display,li2025online} and causal inference \citep[e.g.][]{imbens2015causal,arkhangelsky2021double,fitzmaurice2008longitudinal}. In many applications, the underlying assortment distribution exhibits (partial) independence across products or treatments. For example, in retail demand models, purchase decisions are sometimes approximately independent across product categories \citep[e.g.][]{donnelly2021counterfactual}. In marketing, only a limited set of interaction effects is typically modeled across media channels or competing ads, implying independence across groups of exposures \citep[e.g.][]{zantedeschi2017measuring, danaher2008effect}. In clinical trials, certain combinations of drugs or therapies are designed to act independently \citep[e.g.][]{pomeroy2022drug, plana2022independent}. To estimate assortment probabilities for $M$ items, most existing methods do not exploit (in)dependence structures among products or treatments and can be inefficient when $M$ is large. Methods that exploit such structure typically rely on strong modeling assumptions, such as factorization models \citep[e.g.][]{gopalan2014bayesian}, low-dimensional latent feature representation \citep[e.g.][]{ruiz2020shopper}, shared confounders \citep[e.g.][]{ranganath2018multiple} or independence of pre-specified categories \citep[e.g.][]{donnelly2021counterfactual}.

In contrast, our approach exploits sparsity in the dependence structure to construct estimators with improved finite-sample performance and stronger theoretical guarantees. We only assume that the underlying dependence structure is sparse, without requiring knowledge of its support. Without additional structural assumptions, the distribution of an $M$-dimensional binary vector involves $2^M - 1$ parameters \citep{bahadur1959representation}, consisting of $M$ marginal probabilities and $2^M - M - 1$ \emph{generalized correlation coefficients}. Independence among some components could reduce the parameter number substantially. So in our setup, sparsity naturally arises from the intrinsic sparse generalized correlation structure across binary components.

A natural starting point is the maximum likelihood estimator (MLE). However, the log-likelihood function is highly non-concave in the marginal probabilities, making optimization computationally infeasible, and the MLE tends to overfit when $2^M$ is comparable to or exceeds $N$. Moreover, the MLE fails to adapt to the underlying sparse generalized correlation structure among the binary components. Notably, this non-concavity concerns only the $M$ marginal probabilities, for which consistent estimators are readily available (e.g., sample averages). Conditional on these marginal probabilities, the remaining $2^M - M - 1$ generalized correlation coefficients yield a concave log-likelihood. This naturally motivates a two-stage plug-in approach: first estimate the marginal probabilities, then plug them in to estimate the generalized correlation coefficients via a convex program with an $\ell_1$ penalty to exploit sparsity \citep[e.g.][]{tibshirani1996regression,tibshirani1997lasso,zou2006adaptive,belloni2014high}. We treat the $M$ marginal probabilities as nuisance parameters and the $2^M - M - 1$ generalized correlation coefficients as the parameters of interest. This formulation leads to a high-dimensional estimand in the presence of low- to moderate-dimensional nuisance parameters. In contrast, the debiased machine learning literature \citep[e.g.,][]{chernozhukov2017double} typically considers low-dimensional target parameters with high-dimensional nuisance components. We show that a plug-in approach induces misspecification in the design matrix and fails to achieve rate optimality. While from a statistical performance we want to adapt to the intrinsic dimension, from a computational perspective, the dimension of the problem is $p \asymp 2^M$. 

To address these challenges, we propose two \textit{regularized adversarial estimators}. The first, an \textit{adversarial estimator}, retains the $\ell_1$ penalty and maximizes the worst-case log-likelihood over the marginal probabilities, achieving faster convergence than the plug-in approach. However, it is computationally intractable due to the inner minimization over the marginals. We therefore introduce a novel \textit{first-order estimator}, which approximates the log-likelihood via a first-order expansion in the marginal probabilities, yielding an objective that is concave in the marginals. We construct the adversarial set for the marginal probabilities as a hyper-rectangle, so that the inner minimization is attained at one of its vertices (requiring only $p$ evaluations). Since the objective remains concave in the generalized correlation coefficients, both the inner minimization and outer maximization are computationally tractable. This estimator matches the performance of the adversarial estimator and is rate-optimal under mild conditions, achieving the same rate as in the oracle case where the marginal probabilities are known. We summarize our contributions as follows: 

\paragraph{Contributions}
First, we propose a novel method for estimating assortment probabilities that adapts to the independence structure among $M$ items using the Bahadur representation \citep{bahadur1959representation} without explicit structural assumptions. Second, we develop an adversarial estimation approach for high-dimensional estimand in the presence of low-dimensional nuisance parameters. When the support of the generalized correlation coefficients contains at most $s$ nonzero terms, without assuming their locations, we show that the oracle estimator achieves a rate of $\sqrt{sM/N}$ (the standard Lasso rate) when the marginal probabilities are known. When the marginals are unknown, the plug-in estimator attains a rate of $M\sqrt{s/N}$, while the first-order estimator achieves $\sqrt{sM/N}+\sqrt{M\log M/N}$. Consequently, the first-order estimator is rate-optimal when $s\gtrsim\log M$. Third, we extend the framework to settings with low-dimensional covariates, which require estimating multiple nuisance functions. We establish a similar convergence rate for the \textit{localized first-order estimator}, which outperforms the \textit{localized plug-in estimator}, as supported by our simulation studies. Finally, we apply our framework to causal inference with many binary treatments, where estimating generalized propensity scores \citep[e.g.,][]{rosenbaum1983central,imbens2000role} is challenging. Exploiting independence structures, our estimators are consistent and efficient, with substantial finite-sample gains when the number of treatment combinations is large. Simulations show that our approach outperforms the multinomial logit model \citep[e.g.][]{imai2004causal,li2019propensity,yang2016propensity} commonly used for multi-level treatments.

\paragraph{Related Literature.} 
Some existing works also study the estimation of joint probabilities of binary components or assortments based on \citet{bahadur1959representation}. For example, \citet{sinha2010multivariate} fits a Bahadur representation with higher-order terms set to zero, while \citet{yuan2021community} approximates the joint distribution of within-community connectivities using a truncated Bahadur representation. Unlike these works, we consider the full Bahadur representation without assuming that higher-order terms vanish or that the support of the nonzero terms is known; instead, we impose only a sparsity assumption.
The additional structure of assortments or binary components has also motivated alternative modeling approaches \citep[e.g.][]{banerjee2008model,urschel2017learning,rachkovskij2012similarity,rachkovskij2001representation,rachkovskij2013building,cox1972analysis,rao2004efficiency,gyllenberg1997classification,freund1991unsupervised,ichimura1998maximum}. Specifically, our formulation complements the literature on the conditional independence for undirected graphical models \citep{loh2012structure,morrison2017beyond,van2014new,fan2017high}, including classical Ising model \citep{cipra1987introduction,newell1953theory}. The Ising model, restricted to pairwise dependencies, may still exhibit nonzero higher-order correlations, precluding sparsity in our representation. As discussed in Remark~\ref{rmk:ass1:comparison:ising:graphic}, conditional independence in binary graphical models and sparsity of the Bahadur generalized correlation coefficients \citep{bahadur1959representation} are distinct; Examples~\ref{ex1:sparsity} and~\ref{ex2:sparsity} further show that neither implies the other. Our localized estimator relates to a rich literature on the local maximum likelihood estimation \citep{fan1998local, newey1994kernel, staniswalis1989local, stone1977consistent, tibshirani1987local, lewbel2007local, hansen2008uniform}.  
Our method is also related to recent developments on adversarial approach, including regularized M-estimation \citep{hirshberg2021augmented}, adversarial Riesz representation \citep{chernozhukov2021adversarial}, distributionally robust optimization \citep{blanchet2019quantifying,blanchet2019robust}.

\section{Problem Formulation and Estimators}\label{Sec:Setup}
We have $N$ i.i.d. observations $\{(X_i,Y_i)\}_{i=1}^N$ where $X_i\in\mathcal{X}$ for some compact set $\mathcal{X}\subset\mathbb{R}^d$ and $Y_i \in\{0,1\}^M$. We begin with characterizing $\mathbb{P}(Y_i = y \mid X_i = x)$ by extending the unconditional representation of $\mathbb{P}(Y_i = y)$ by \citet{bahadur1959representation}. For a binary vector $y \in \{0,1\}^M$ and marginal probabilities $\alpha=(\alpha_1,\ldots,\alpha_M)\in (0,1)^M$, we define the $j$-th standardized component as 
\begin{equation}\label{eq:z}
\displaystyle z_j(y,\alpha):=(y_{j}-\alpha_j)/\sqrt{\alpha_j(1-\alpha_j)}.
\end{equation}
For $j\in[M]:=\{1,2,\ldots,M\}$, we denote the $j$-th conditional marginal probability as $\alpha_{0j}(x) = \mathbb{P}(Y_{ij}=1 \mid X_i=x)$, which is the $j$-th component of vector $\alpha_0(x)\in(0,1)^M$. Note $\mathbb{E}[ z_j(Y_i,\alpha_0(x))\mid X_i=x]=0$ and $\mathrm{Var}[ z_j(Y_i,\alpha_0(x))\mid X_i=x]=1$. A $k$-tuple $ \ell = \{j_1,\ldots,j_k\} \subseteq [M]$ such that $k=|\ell|\geq 2$ represents an assortment. Since there are in total of $p=2^M-M-1$ possible assortments, we index each component $l \in [p]$ by its corresponding assortment $\ell = \ell(l)$. Given covariate $x\in\mathcal X$, the \textit{generalized correlation coefficients} are $r_0(x):=\left(r_{01}(x),\ldots,r_{0p}(x)\right)$, where for any $l\in [p]$, 
\begin{equation}\label{def:r0}
r_{0l}(x):= \mathbb{E}\left[ \prod_{j\in\ell(l)}z_j(Y_{i},\alpha_0(X_i)) \bigg| X_i=x\right].
\end{equation}
So each entry $r_{0l}(x)$ corresponds to an assortment $\ell(l)$ such that $|\ell(l)|\geq2$.  \citet{bahadur1959representation} characterizes the joint distribution of a multivariate binary vector in the absence of covariates, and its extension to the conditional distribution of all binary components given covariates is straightforward, following Propositions 1 and 2 in \citet{bahadur1959representation}:
\begin{theorem}\label{Thm:Bahadur}
$\mathbb{P}(Y_i = y \mid X_i=x )= f(y,\alpha_0(x),r_0(x))\prod_{j=1}^M \alpha_{0j}(x)^{y_j}(1-\alpha_{0j}(x))^{1-y_j}$ for any $y \in \{0,1\}^M$ and $x \in \mathcal{X}$, where $\displaystyle f(y,\alpha_0(x),r_0(x)):= 1 + \sum_{\ell(l)\subseteq [M] : \left|\ell(l)\right|\geq2} r_{0l}(x)\prod_{j\in\ell(l)} z_{j}(y,\alpha_0(x))$. 
\end{theorem}
Theorem \ref{Thm:Bahadur} characterizes the dependence between components through the generalized correlation coefficients $r_0(x)=(r_{0l}(x))_{l\in[p]}$ within function $f$. By defining a $p$-dimensional vector $W(\alpha,y)$ where for any $l\in[p]$, its $l$-th component is equal to $W_{l}(\alpha,y):=\prod_{j\in \ell(l)}z_j(y,\alpha)$, we can write the function $f$ as an affine function of the generalized correlation coefficients: $f(y,\alpha_0(x),r_0(x))= 1+W(\alpha_0(x),y)'r_0(x)$.
If all the binary components are independent, $r_{0}(x)=0$ and $f(y,\alpha_0(x),0)=1$, which recovers the conditional likelihood of $y$ given $x$ as $\prod_{j=1}^M \alpha_{0j}(x)^{y_j}(1-\alpha_{0j}(x))^{1-y_j}$ under full independence. According to \eqref{def:r0}, for any $l\in[p]$, if within the assortment $\ell(l)$ there is at least one binary component independent from the others conditioning on the covariates, then $r_{0l}(x)=0$. This connects independence structure with zero value of a  generalized correlation coefficient. Thus in applications of interest where independence holds among some of the binary components, we can expect the vector $r_0(x) \in \mathbb{R}^p$ to be sparse. Theorem~\ref{Thm:Bahadur} implies that estimating the assortment probabilities requires estimating $\alpha_0(x)$ and $r_0(x)$. 

For the remainder of this section, we first focus on the \textit{no-covariate case}, i.e. the unconditional probabilities $\mathbb{P}(Y_i=y)$. We begin with the \textit{maximum likelihood estimator} (MLE) as a natural baseline and discuss its limitations. Building on the MLE, we introduce a \textit{plug-in estimator} that improves upon it both statistically and computationally. However, the plug-in estimator is not rate-optimal. To address this, we consider two \textit{regularized adversarial estimators}. The first, termed the \textit{adversarial estimator}, is a natural extension of the plug-in approach but is computationally infeasible. We therefore propose the \textit{first-order estimator} as a second regularized adversarial estimator, which is computationally feasible and rate-optimal in common parameter regimes under mild conditions. Building on the results established for the no-covariate case, we then extend the proposed estimators to the \textit{with-covariate case}, i.e. the conditional probabilities $\mathbb{P}(Y_i=y\mid X_i=x)$. In this case, we obtain localized counterparts of the plug-in, adversarial, and first-order estimators. Theoretical properties of these estimators parallel those in the no-covariate case, while the inclusion of covariates introduces additional technical challenges.

\paragraph{Maximum Likelihood Estimator} We first consider MLE for $(\alpha_0,r_0)$:
$$(\hat \alpha^{mle}, \hat r^{mle}) \in \argmax_{(\alpha,r)\in \mathbb{K}} \sum_{i=1}^N\left\{\log( 1+W(\alpha,Y_i)'r )+ \sum_{j=1}^M Y_{ij} \log(\alpha_j)+ (1-Y_{ij}) \log(1-\alpha_j)\right\}$$
where $\mathbb{K}= \left\{ (\alpha,r) :  1+W(\alpha,y)'r \geq 0, \forall y\in\{0,1\}^M\right\}$ includes $(\alpha,r)$ leading to non-negative probability estimates for all possible assortments. However, the MLE faces practical limitations: First, the objective is non-concave in $\alpha$ due to the standardized outcomes \eqref{eq:z} in the formulation of vector $W(\alpha,Y_i)$. Second, the number of parameters grows exponentially with $M$, leading to overfitting even for moderate values of $M$. Third, the MLE fails to exploit the underlying independence structure, which could otherwise improve estimation efficiency.

\paragraph{Plug-in Estimator} Although the objective function for MLE is non-concave in $\alpha$, it is well-behaved in $r$. Fixing $\alpha$, the objective function of MLE becomes a concave maximization problem in $r$. Since independence implies sparsity in the coefficient vector $r_0$, the $\ell_1$-regularized estimator is particularly suitable. Moreover, consistent estimators for $\alpha_0$ (e.g., sample averages) are readily available. A straightforward approach is therefore to first estimate $\alpha_0$ by $\hat{\alpha}=(\hat{\alpha}_1,\ldots,\hat{\alpha}_M)$, plug $\hat{\alpha}$ into the objective of MLE, and apply an $\ell_1$ penalty to $r$. The resulting estimator is referred to as the plug-in estimator $\hat r^{PI}$ for the no-covariate case: 
\begin{equation}\label{eq:PMLE-no-covariates}
     \hat r^{PI}\! \in \!\arg\max_{r}  \displaystyle \frac{1}{N}\sum_{i=1}^N  \log( 1+ W(\hat \alpha,Y_i)'r ) - \lambda \|r\|_{1,w},
\end{equation}
where $\|r\|_{1,w}:= \sum_{j=1}^p w_j |r_j|$ is a weighted $\ell_1$-norm and $\alpha$ are the nuisance parameters. Additionally, the weights are defined as $w_j :=\{\frac{1}{N}\sum_{i=1}^NW_j^2(\hat{\alpha},Y_i)/(1+W(\hat{\alpha},Y_i)'\hat{r}^{(0)})^2\}^{1/2}$ for any $j\in[p]$, where $\hat{r}^{(0)}$ denotes an initial estimator for $r_0$. Theoretically, the ideal implementation would set $w_j =\{\frac{1}{N}\sum_{i=1}^NW_j^2(\alpha_0,Y_i)/(1+W(\alpha_0,Y_i)'r_0)^2\}^{1/2}$, which balances the heterogeneity across components. However, since $\alpha_0$ and $r_0$ are unknown, we instead use their estimates $\hat{\alpha}$ and $\hat{r}^{(0)}$. Although the weights depend on unknown quantities, stable estimates can be obtained via an iterative procedure with mild over-penalization. Note that $\mathrm{dim}(\alpha_0)=M$ and $\mathrm{dim}(r_0)=2^M-M-1=p$, so $\mathrm{dim}(r_0)\gg\mathrm{dim}(\alpha_0)$ as $M$ grows. When $p\gg M$, we can obtain estimators for $\alpha_0$ with a faster rate of convergence than for $r_0$. Specifically, for the rest of the paper, we proceed with the consistent estimator 
\begin{equation}\label{eq:SAA}
    \hat\alpha_j = \frac{1}{N}\sum_{i=1}^N Y_{ij}
\end{equation}
for $\alpha_{0j}$, $\forall j\in [M]$, to be used in all no-covariate case estimators. Since $\alpha_{0j}\in(0,1)$, for sufficiently large $N$, $\hat{\alpha}_j(1-\hat{\alpha}_j)>0$ with high probability, so $z_j(y,\hat{\alpha}_j)$ is well defined. In practice if $\hat{\alpha}_j(1-\hat{\alpha}_j)$ becomes numerically close to zero, one can apply a mild truncation or add a small constant to the denominator to avoid division by zero. 
\begin{remark}\label{remark:non-negativity}
In the estimation of $r_0$ for the plug-in estimator, we do not impose the constraint on $\mathbb{K}$, because we are using $\hat\alpha$, and it is possible that $(\hat\alpha,r_0) \not\in \mathbb{K}$. By construction of the objective function (\ref{eq:PMLE-no-covariates}), since $r=0$ is always feasible, by optimality of $\hat{r}^{PI}$ we have $1+W(\hat\alpha,Y_i)'\hat r^{PI}> 0$ for every $Y_i$ in the data, so the probability estimates is strictly positive for any outcome vector in the data. We also show in the proofs that under mild conditions, $1+W(\hat\alpha,Y_i)'r_0 > 0$ for any $Y_i$ in the data. Similar arguments hold for the estimators $\hat{r}^A, \hat{r}^{FO}$ defined later. This result is sufficient in many applications, because practitioners typically only need probability estimates for assortments that are actually observed in the data. For example, in A/B testing or recommendation systems, a firm may only evaluate the probability of a limited set of product or advertisement assortments that are presented to users, rather than all possible combinations of items. Similarly, in retail assortment analytics, firms typically analyze purchase patterns among product bundles that have appeared in historical transactions, without requiring probability estimates for bundles that have never been offered. Related discussion on possible extensions is included in Section~\ref{sec:discussion}. 
\proofend
\end{remark}
Although $\hat{r}^{PI}$ is computationally tractable, plugging in $\hat{\alpha}$ induces mis-specification in the score function related to \eqref{eq:PMLE-no-covariates}. The resulting discrepancy is dominated by the $\ell_1$ estimation error $\|\hat{\alpha}-\alpha_0\|_1 = \mathrm{O}_{\mathrm{p}}(M/\sqrt{N})$, introducing additional sparsity-dependent factors and yielding a slower convergence rate than the typical lasso rate. Detailed discussions are provided in Section~\ref{sec:theoretical}. 

\paragraph{Regularized Adversarial Estimators} Faced with the slower rate of convergence for the plug-in estimator, we propose an adversarial estimator $\hat{r}^A$, defined as 
\begin{equation}\label{eq:PMLE-no-covariates-non-linear}
    \displaystyle \hat r^{A}\in \argmax_{r}\min_{\alpha \in \mathcal{A}(\hat\alpha)}\frac{1}{N}\sum_{i=1}^N \log(1+ W(\alpha,Y_i)'r ) - \lambda \|r\|_{1,w}
\end{equation}
where $\mathcal{A}(\hat{\alpha})$ is an adversarial set such that $\alpha_0 \in \mathcal{A}(\hat{\alpha})$ with high probability. From a statistical perspective, the minimization of  $\alpha \in \mathcal{A}(\hat\alpha)$ acts as a regularization allowing us to better control the mis-specification for not knowing $\alpha_0$. This allows us to use a smaller penalty parameter $\lambda$, mitigating the bias induced by the overly large $\lambda$ required by the plug-in estimator due to overfitting on $\hat{\alpha}$. Since the minimum of concave functions is still concave, the objective function \eqref{eq:PMLE-no-covariates-non-linear} is concave in $r$. However, the inner problem of \eqref{eq:PMLE-no-covariates-non-linear} requires the minimization over $\alpha$ through the highly non-linear and non-concave mapping $W(\alpha,\cdot)$. This computational challenge motivates us to further propose the first-order estimator $\hat{r}^{FO}$, defined as
\begin{equation}\label{eq:PMLE-no-covariates-bilinear}   
    \hat r^{FO}\in\argmax_{r}\min_{\alpha \in \mathcal{A}(\hat\alpha)}\frac{1}{N}\sum_{i=1}^N\log\left(1+W^{FO}(\alpha,Y_i)'r\right)-\lambda \|r\|_{1,w},
\end{equation}
where $W^{FO}(\alpha,Y_i):=W(\hat\alpha,Y_i)+\nabla_{\alpha} W(\hat\alpha,Y_i)(\alpha-\hat\alpha)$.
Estimator (\ref{eq:PMLE-no-covariates-bilinear}) also adopts an adversarial approach but uses a first-order approximation of $W(\alpha,Y_i)$ with respect to $\alpha$ around $\hat{\alpha}$. The estimator has a bilinear form and is concave in $\alpha$ for any fixed $r$. Hence, the minimization of this concave function is attained at the extreme points of $\mathcal{A}(\hat{\alpha})$. For appropriately structured sets $\mathcal{A}(\hat{\alpha})$, such as a hyper-rectangle $\mathcal{A}(\hat{\alpha}) = \prod_{j=1}^M [\hat{\alpha}_j - c_j, \hat{\alpha}_j + c_j]$, where $c_j$ denotes the confidence interval half-width and $\hat{\alpha}_j$ is defined by~\eqref{eq:SAA}, the problem remains computationally feasible. As a result, the discrepancy in the score function becomes negligible since it includes only the second-order bias in the marginal probabilities, in contrast to the plug-in estimator $\hat{r}^{PI}$. As shown in Theorem~\ref{thm:main:FO}, the $\lambda$ for $\hat{r}^{FO}$ essentially bounds the $\ell_\infty$-norm of the score as if the marginal probabilities were known.

\paragraph{Incorporating Covariates} The insights and estimators from the no-covariate case extend naturally to the with-covariates case. However, the nuisance parameters associated with the marginal probabilities are now functions, which complicates the construction of the adversarial set, as discussed later. Given initial estimates $\hat\alpha(\cdot)=\left(\hat\alpha_1(\cdot),\ldots,\hat\alpha_M(\cdot)\right)$, we propose the following localized versions of the estimators \eqref{eq:PMLE-no-covariates}, \eqref{eq:PMLE-no-covariates-non-linear}, \eqref{eq:PMLE-no-covariates-bilinear} given any covariate $x$: 
$$\hat r^{PI}(\cdot,x) \in \arg \max_{a, b}  \displaystyle  \frac{1}{N}\sum_{i=1}^N  K_h\left( \mbox{$X_i-x$}\right)  \log( 1+W(\hat\alpha (X_i),Y_i)'r_{a,b}(X_i,x)) - \lambda_x\|r_{a,b}\|_{w_x,h,1}.$$
$$\hat r^{A}(\cdot,x) \in \arg \max_{a, b} \min_{\alpha \in \mathcal{A}_x(\hat\alpha)} \displaystyle  \frac{1}{N}\sum_{i=1}^N  K_h\left( \mbox{$X_i-x$}\right)  \log( 1+W(\alpha(X_i),Y_i)'r_{a,b}(X_i,x) ) - \lambda_x \|r_{a,b}\|_{w_x,h,1}$$
$$\begin{array}{l}
\displaystyle \hat r^{FO}(\cdot,x) \in \arg \max_{a, b} \min_{\alpha \in \mathcal{A}_x(\hat\alpha)}  \frac{1}{N}\sum_{i=1}^N  K_h\left( \mbox{$X_i-x$}\right)  \log\left( 1+ W(\hat\alpha(X_i),Y_i)'r_{a,b}(X_i,x) \right.\\
\displaystyle \hspace{5.5cm} \left. +(\alpha(X_i)-\hat\alpha(X_i) )'\nabla_\alpha W(\hat\alpha(X_i),Y_i)'r_{a,b}(X_i,x) \right)  - \lambda_x \|r_{a,b}\|_{w_x,h,1},
\end{array}$$
where $K_h(X_i-x):=\frac{1}{h^d}K\left(\frac{X_i-x}{h}\right)$, $K:\mathbb{R}^d\rightarrow \mathbb{R}$ is the kernel function, and $h$ is the bandwidth. For $a\in \mathbb{R}^p$, $b \in \mathbb{R}^{p\times d}$, $r_{a,b}(X_i,x):= a + b(X_i-x)$, $\|r_{a,b}\|_{w_x,h,1} :=\sum_{j=1}^p w_{xj} |a_j| + h\sum_{j=1}^p\sum_{k=1}^d w_{xjk}|b_{jk}|$, and the weights are $w_{xj} :=\left\{\frac{1}{N}\sum_{i=1}^NK_h(X_i-x)W_j(\hat{\alpha}(X_i),Y_i)^2/(1+W(\hat{\alpha}(X_i),Y_i)'\hat{r}^{(0)}(X_i,x))^2\right\}^{1/2}$, $w_{xjk}:=\left\{\frac{1}{N}\sum_{i=1}^NK_h(X_i-x)W_j(\hat{\alpha}(X_i),Y_i)^2[(X_{ik}-x_k)/h]^2/(1+W(\hat{\alpha}(X_i),Y_i)'\hat{r}^{(0)}(X_i,x))^2\right\}^{1/2}$ for any $j\in[p]$, $k\in[d]$, where $\hat{r}^{(0)}(X_i,x)$ denotes an initial estimator for $r_0(X_i)$. The initial estimates for $\alpha_{0j}(x)$ can be obtained from different approaches. In this paper we focus on the localized estimator defined for any $j\in[M]$ as follows:
\begin{equation}\label{def:localized:alpha:estimator}
\left(\tilde{\alpha}_j(x), \tilde{\beta}_j(x)\right)\in \arg\!\min_{(\tilde a,\tilde b)}\frac{1}{N}\!\sum_{i=1}^N\!K_h(X_i-x)\!\left\{Y_{ij} - \tilde a - \tilde b(X_i-x)\right\}^2.
\end{equation} 
We construct the data-dependent adversarial set $\mathcal{A}_x(\hat{\alpha})$ based upon \eqref{def:localized:alpha:estimator}. All three objectives for the estimators in the covariate setting are concave functions in $r_{a,b}$. At covariate $X_i$, the generalized correlation coefficient vector is $r_0(X_i)$, where $r_0(X_i)=r_0(x)+\nabla r_0(x)(X_i-x)+\Delta$ and $\|\Delta\|_{\infty}=\mathrm{O}\left(\|X_i-x\|^2\right)$ under proper regularity conditions, where $\|\cdot\|, \|\cdot\|_\infty$ denote the Euclidean norm and the $\ell_\infty$ norm throughout the paper. Suppose $K(\cdot)$ is supported on $\{u\in\mathbb{R}^d:\|u\|\le1\}$. Then, by the definition of $K_h(\cdot)$, for any $X_i$ such that $K_h(X_i-x)\neq0$, we have $\mathrm{O}\left(\|X_i-x\|^2\right)=\mathrm{O}(h^2)$. Hence, $r_{a,b}(X_i,x)$ effectively approximates $r_0(X_i)$ when $X_i$ is close enough to $x$, where $a$ approximates $r_0(x)$ and $b$ approximates its gradient $\nabla r_0(x)$. The minimization over $\alpha(\cdot) \in \mathcal{A}_x(\hat{\alpha})$ is again attained at an extreme point of the nuisance function set $\mathcal{A}_x(\hat{\alpha})$. However, care must be taken in specifying this set, as it can easily become too large. For instance, since the nuisance functions depend on $X_i$, one might be tempted to minimize over each $\alpha(X_i)$ separately, but this would create an overly adversarial setup that ignores the dependence across covariates $\{X_i\}_{i\in[N]}$. We define the set of functions $\mathcal{A}_x(\hat\alpha)$ such that the function $\alpha_{*x}(X)=\alpha_0(x)+\nabla_x\alpha_0(x)(X-x) \in \mathcal{A}_x(\hat\alpha)$ with high probability. Additionally, the penalty for the gradient estimation term $b$ is scaled by the bandwidth parameter $h$. Intuitively, since the estimation target is $r_0(x)$, the factor $h$ attenuates the effect of the estimation error in $\nabla r_0(x)$ through $b$. Note that $r_{a,b}(X_i,x)$ provides a first-order approximation to $r_0(X_i)$, and higher-order refinements are possible in principle. Such extensions, however, would require additional regularity and smoothness assumptions. Below we include a remark on higher-order approximations:
\begin{remark}\label{rmk:higher:order:approx}
One could alternatively define the localized estimator with higher order terms. For example, for the following version of the plug-in estimator
$$\begin{array}{rl}
\displaystyle\max_{a,b,c}\displaystyle\frac{1}{N} \sum_{i=1}^N  K_h\left( \mbox{$X_i-x$}\right)  \log( 1+W(\hat\alpha (X_i),Y_i)'r_{a,b,c }(X_i,x)  ) - \lambda_x \|r_{a,b,c}\|_{w_x,h,1}, 
\end{array}$$
where $a\in \mathbb{R}^p, b \in \mathbb{R}^{p\times d}, c \in \mathbb{R}^{p\times d\times d}$, $r_{a,b,c}(X_i,x):= a + b(X_i-x) + c[X_i-x,X_i-x]$,
$$\|r_{a,b,c}\|_{w_x,h,1} :=\sum_{j=1}^p w_{xj} |a_j| + h\sum_{j=1}^p\sum_{k=1}^d w_{xjk} |b_{jk}| + h^2\sum_{j=1}^p\sum_{k=1}^d\sum_{m=1}^d w_{xjkm} |c_{jkm}|,$$ 
and $c[X_i-x,X_i-x]$ is a contraction of the $p\times d\times d$ tensor with two copies of $X_i-x$, producing a $p$-dimensional vector. This formulation increases the dimension of the parameter space. It would require higher-order smoothness conditions and needs to account for the sparsity in the second-order derivative of $r_0(x)$. Meanwhile it would also reduce the approximation errors by improving the dependence on the bandwidth $h$. \proofend
\end{remark}

\paragraph{Notations.} Throughout the paper, $\|\cdot\|_1$, $\norm{\cdot}_2$ (or $\norm{\cdot}$), $\norm{\cdot}_\infty$, and $\norm{\cdot}_0$ denote the $\ell_1$-, $\ell_2$-, $\ell_\infty$-, and $\ell_0$-norms, respectively. For a positive integer $n$, $[n]:=\{1,\ldots,n\}$.
For a matrix $A\in\mathbb{R}^{m\times n}$, $\norm{A}_{\infty,\infty}:=\max_{i\in[m],j\in[n]}|A_{ij}|$, $\norm{A}_{k,k}:=(\sum_{i=1}^m\sum_{j=1}^n|A_{ij}|^k)^{1/k}$, $\norm{A}_0$ denotes the number of nonzero entries. $\mathbf{1}\{\ \cdot\ \}$ denotes the indicator. $\propto$ denotes proportionality. For any vector or matrix $V$, $V'$ denotes the transpose of $V$. Given $\delta\in\mathbb{R}^p$ and $T\subseteq[p]$, define $\delta_T$ as $\delta$ restricted to $T$. For random variables $Z$ and $W$, $Z\indep W$ denotes independence. For any finite-dimensional real vector $\nu$, $\mathrm{dim}(\nu)$ denotes its dimension. $\mathcal{N}(\mu,\sigma^2)$ denotes the Gaussian distribution with mean $\mu$ and variance $\sigma^2$. For deterministic sequences ${a_n}$ and ${b_n}$, we write $a_n = \mathrm{O}(b_n)$ if $|a_n| \le C |b_n|$ for some constant $C>0$ and all sufficiently large $n$, and $a_n = \mathrm{o}(b_n)$ if $a_n/b_n \to 0$ as $n \to \infty$. For stochastic sequences, $X_n = \mathrm{O}_{\mathrm{p}}(a_n)$ means $X_n/a_n$ is bounded in probability, and $X_n = \mathrm{o}_{\mathrm{p}}(a_n)$ means $X_n/a_n \xrightarrow{p} 0$, $X_n=\Omega_p(a_n)$ means $X_n/a_n$ is bounded away from zero with non-vanishing probability. We write $a \gtrsim b$ if there exists a universal constant $C_1>0$ such that 
$a \ge C_1 b$ for all sufficiently large $N$.  


\section{Theoretical Results}\label{sec:theoretical}
In this section, we present the main theoretical results. We start with assumptions for the with-covariate case. The no-covariate case can be viewed as a special case where the covariate space $\mathcal{X}$ is a singleton. In this case, $\alpha_0(x)=\alpha_0$, $r_0(x)=r_0$, and $\nabla_x r_0(x)=0$, so that the sparsity condition reduces to $\|r_0\|_0\le s$, corresponding to condition (i) of Assumption~\ref{ass:iid-min-den}. We focus on estimating $r_0(x)$ and $\alpha_0(x)$ at a given covariate $x$ in the interior of $\mathcal{X}$. 

\begin{assumption}[Sparsity and regularity conditions]\label{ass:iid-min-den}
(i) $\sup_{x \in \mathcal{X}} \|r_0(x)\|_0+\|\nabla r_0(x)\|_0 \leq s$. (ii) For some positive constant $\bar{L}$, ${ \displaystyle \max_{x\in\mathcal{X}, y\in \{0,1\}^M}} 1+W(\alpha_0(x),y)^\prime r_0(x) \leq \bar L$ (iii) ${ \displaystyle \min_{i\in [N]}} 1+W(\alpha_0(X_i),Y_i)'r_0(X_i) > \xi$ with probability $1-\epsilon$. (iv) $\alpha_{0j}(x)$ is bounded away from $0$ and $1$ (i.e. $\exists\ \eta\in(0,1)$, such that $\eta<\alpha_{0j}(x)<1-\eta$) for all $j\in[M]$, $\forall x\in\mathcal{X}$. 
\end{assumption}
Condition (i) of Assumption~\ref{ass:iid-min-den} bounds the degree of sparsity on $r_0(x) \in \mathbb{R}^p$ and $\nabla r_0(x) \in \mathbb{R}^{p\times d}$. When all binary components are independent from each other (for all values of $x\in\mathcal{X}$), we have $s=0$. The analysis also directly extends to a local measure of sparsity (i.e., $s_x$). For instance, when only pairwise interactions are present and higher-order interactions are absent, we have $\|r_0(x)\|_0+\|\nabla r_0(x)\|_0\leq(d+1)\binom{M}{2}\ll 2^M$.
Conditions (ii) and (iii) impose upper and lower bounds on the deviations from full independence. For example, this condition holds in empirical settings where there exist $q_1, q_2>0$ such that the conditional probabilities are bounded within multiplicative factors $q_1$ and $q_2$ of the independent probabilities, i.e., $q_1\prod_{j=1}^M\alpha_{0j}(x)^{y_j}(1-\alpha_{0j}(x))^{1-y_j}\le \mathbb{P}(Y_i=y\mid X_i=x) \le q_2\prod_{j=1}^M\alpha_{0j}(x)^{y_j}(1-\alpha_{0j}(x))^{1-y_j}$. Assumption \ref{ass:iid-min-den} allows for situations where some combinations of treatments have zero probability (so no $Y_i$ equal to that value exists in the data). 
\begin{remark}\label{rmk:ass1:comparison:ising:graphic}
The sparsity assumption in Assumption~\ref{ass:iid-min-den} differs from that in sparse Ising or binary graphical models, which impose sparsity on canonical interaction parameters. The two notions are fundamentally distinct: Even in a tree-structured Ising model, most pairwise correlations are typically nonzero, and Bahadur generalized correlation coefficients can be dense despite graph sparsity. Conversely, sparsity in the Bahadur generalized correlation coefficients does not imply a sparse conditional independence graph either. In particular, a distribution with only a single nonzero generalized correlation coefficient can still induce dense graphical conditional dependencies.
\proofend
\end{remark}
Examples~\ref{ex1:sparsity} and~\ref{ex2:sparsity} below illustrate our argument. In both examples, $Y_i=(Y_{i1},Y_{i2},Y_{i3},Y_{i4})\in\{0,1\}^4$ denotes the purchase decisions of the $i$-th customer for an assortment of four products. Suppose $\alpha_{0j}=\frac{1}{2}$ for each $j\in[4]$, and $Z_{ij}:=z_{j}(Y_{ij},\alpha_{0j})=2Y_{ij}-1\in\{-1,+1\}$. We use $r_{0,\ell}$ to denote the generalized correlation coefficient associated with the coordinate corresponding to bundle $\ell\subseteq[4]$.
\begin{example}\label{ex1:sparsity}
 Suppose $\mathbb{P}(Z_i=(z_1,z_2,z_3,z_4))\propto\exp\{\theta(z_1z_2+z_2z_3+z_3z_4)\}$ with $\theta\neq0$. This corresponds to a pairwise markov random field of the four-node Ising chain: ``$1-2-3-4$'' with only three edges $(1,2), (2,3), (3,4)$, and each node depends only on its neighbors, so $Z_{i1}\indep(Z_{i3},Z_{i4})|Z_{i2}$, $Z_{i4}\indep (Z_{i1},Z_{i2})|Z_{i3}$, etc. 
This is a sparse conditional independence graph. However, all the pairwise generalized correlation coefficients are nonzero, since for some $\rho\neq0$, $r_{0,12}=r_{0,23}=r_{0,34}=\rho$, $r_{0,13}=r_{0,24}=\rho^2$, $r_{0,14}=\rho^3$. Additionally, $r_{0,1234}=\rho^2$. So graphical model sparsity doesn't imply the sparsity assumption in Assumption~\ref{ass:iid-min-den}.  \proofend
\end{example}
\begin{example}\label{ex2:sparsity}
Suppose $\mathbb{P}(Z_i=(z_1,z_2,z_3,z_4))\propto(1+\gamma z_1z_2z_3z_4)/16$ where $0<|\gamma|<1$. Then $r_{0,jk}=0$ for any pair $\{j,k\}$, $r_{0,jkl}=0$ for any triple $\{j,k,l\}$, but $r_{0,1234}=\gamma\neq0$. So $r_0$ is sparse with the only nonzero element $r_{0,1234}$. Since $\mathbb{P}(Z_{i1}=z_1,Z_{i2}=z_2|Z_{i3}=z_3,Z_{i4}=z_4)\propto 1+\gamma z_1z_2z_3z_4$, it can be checked by direct calculation that $\mathbb{P}(Z_{i1}=z_1,Z_{i2}=z_2|Z_{i3}=z_3,Z_{i4}=z_4)\neq\mathbb{P}(Z_{i1}=z_1|Z_{i3}=z_3,Z_{i4}=z_4)\mathbb{P}(Z_{i2}=z_2|Z_{i3}=z_3,Z_{i4}=z_4)$ unless $\gamma=0$. Thus $\mathbb{E}[Z_{i1}Z_{i2}|Z_{i3},Z_{i4}]\neq0$. By symmetry, $\mathbb{E}[Z_{ik}Z_{ij}|Z_{il},Z_{im}]\neq0$ for all distinct $(k,j,l,m)$. So the binary graphical model has all the edges and is not sparse. \proofend
\end{example}
Next, Assumption~\ref{ass:moment-assumptions} imposes moment conditions and conditions on $M,N,h,d,s$:
\begin{assumption}[Moment conditions]\label{ass:moment-assumptions}
Let $\displaystyle\widetilde W(\alpha_0(x),Y)\!\!:=\!\!\frac{[\! W(\alpha_0(x),Y) ;\!\nabla_\alpha W(\alpha_0(x),Y)'r_0(x) ]}{1+W(\alpha_0(x),Y)'r_0(x)}$, $\bar \psi_2 := {\displaystyle \min_{x\in \mathcal{X}, l\in [p+M]}} \mathbb{E}[ \widetilde W_{l}^2(\alpha_0(x),Y)\mid X=x]^{1/2},$ $\bar K_m = {\displaystyle\max_{x\in \mathcal{X}}\max_{l\in [p+M]}}\mathbb{E}[|\widetilde W_{l}(\alpha_0(X),Y)|^m\mid X=x]^{1/m}$, and $\bar K_{\max} = \mathbb{E}\big[ {\displaystyle\max_{l\in [p+M], i \in [N]}}\widetilde W_{l}^2(\alpha_0(X_i),Y_i)\big]^{1/2}$. $\sqrt{\frac{M^3}{Nh^d}}\leq\delta_N$, $sdh^2\leq\delta_N\xi$, $Nh^{d+4}\geq\frac{\log(NMd)}{d^2}$, $d^2h<1$, $M^3\leq N\delta_N^{2+d/2}\left(\frac{\xi}{sd}\right)^{d/2}$, $\delta_N\leq\min\left\{\frac{1}{2},\frac{s}{d^3\xi},\frac{d^{d/(d+4)}M^{3/2}}{N^{\frac{2}{d+4}}[\log(NMd)]^{\frac{d}{2d+8}}}\right\}$, $\delta_N\rightarrow0$, $h\rightarrow0$, $Nh^d\rightarrow\infty$.
\end{assumption}
In the no-covariate case Assumption~\ref{ass:moment-assumptions} implies $\sqrt{M^3/N}\leq\delta_N$. In the with-covariate case, by direct calculation, Assumption \ref{ass:moment-assumptions} implies that $h$ satisfies $\left(\frac{M^3}{N\delta_N^2}\right)^{1/d}\leq h\leq\sqrt{\frac{\delta_N\xi}{sd}}$, this interval is nonempty since we assume $M^3\leq N\delta_N^{2+d/2}\left(\frac{\xi}{sd}\right)^{d/2}$.  Assumption \ref{ass:smoothness-assumptions} in the below imposes smoothness conditions:
\begin{assumption}[Smoothness conditions]\label{ass:smoothness-assumptions}
With probability $1-\delta_\alpha$, conditions (1) - (6) hold as follows: (1) $\alpha_0\in\mathcal{A}(\hat{\alpha})$, (2) ${\displaystyle\max_{j\in [p], \tilde \alpha \in \mathcal{A}(\hat\alpha)}}|v'\nabla_\alpha W_j(\tilde \alpha,Y_i)|  +|v'\nabla_\alpha W(\tilde \alpha,Y_i)'r_0| \leq M^{1/2}K_{\alpha,1}\|v\|$, \\
(3) $|u'\nabla_{\alpha}W(\alpha_0,Y_i)'r_0|\leq\bar{K}_w\|u\|_1$, $|u'\nabla_{\alpha}W(\alpha_0,Y_i)'v|\leq \bar{K}_w\norm{u}_1\norm{v}_{1}$, \\
(4) $\displaystyle\max_{j \in [p], \tilde \alpha \in \mathcal{A}(\hat\alpha)}|u'\nabla_\alpha^2 W(\tilde\alpha,Y_i)'r_0v|+|u'\nabla_\alpha^2 W_j(\tilde\alpha,Y_i)v|\leq K_{\alpha,2}\|u\|\|v\|$, (5) ${\displaystyle \max_{\tilde\alpha \in \mathcal{A}(\hat\alpha)}}\|\tilde \alpha-\alpha_0\|\frac{M^{1/2}K_{\alpha,1}}{\xi}<\frac{1}{4}$, ${\displaystyle\max_{\tilde\alpha \in \mathcal{A}(\hat\alpha)}}\|\tilde\alpha-\alpha_0\|^2 \frac{K_{\alpha,2}}{\xi}\leq\frac{1}{16}$, (6) $|W(\alpha_0,Y_i)'v|\leq\bar{K}_w\|v\|_1$, $u'(\nabla_{\alpha}W(\alpha_0,Y_i)'r_0)(\nabla_{\alpha}W(\alpha_0,Y_i)'r_0)'u\leq K_{\alpha,2}\|u\|^2$.

\noindent With probability $1-\delta_\alpha$, for the with-covariate case, over all $i\in[N]$ where $\|X_i-x\|\leq h$, we have:
(i) {$\alpha_{*x}(X):=\alpha_0(x)+\nabla_x\alpha_0(x)(X-x) \in \mathcal{A}_x(\hat\alpha)$},\\
(ii) ${\displaystyle\max_{j\in [p], \tilde \alpha \in \mathcal{A}_x(\hat\alpha)}}|v'\nabla_\alpha W_j(\tilde \alpha(X_i),Y_i)|  +|v'\nabla_\alpha W(\tilde \alpha(X_i),Y_i)'r_0(X_i)| \leq M^{1/2}K_{\alpha,1}\|v\|$, \\
(iii)$|u'\nabla_{\alpha}W(\alpha_0(X_i),Y_i)'r_0(X_i)|\leq\bar{K}_w\|u\|_1$, $|u'\nabla_{\alpha}W(\alpha_0(X_i),Y_i)'v|\leq \bar{K}_w\norm{u}_1\norm{v}_{1}$,\\
(iv)$\displaystyle\max_{j \in [p], \tilde \alpha \in \mathcal{A}_x(\hat\alpha)}|u'\nabla_\alpha^2 W(\tilde\alpha(X_i),Y_i)'r_0(X_i)v|+|u'\nabla_\alpha^2 W_j(\tilde\alpha(X_i),Y_i)v|\leq K_{\alpha,2}\|u\|\|v\|$,\\
(v)${\displaystyle \max_{\tilde\alpha \in \mathcal{A}_x(\hat\alpha),\|\tilde x-x\|\leq h}}\|\tilde \alpha(\tilde x)-\alpha_0(\tilde x)\|\frac{M^{1/2}K_{\alpha,1}}{\xi}<1/4$,  ${\displaystyle \max_{\tilde\alpha \in \mathcal{A}_x(\hat\alpha),\|\tilde x-x\|\leq h}}\|\tilde \alpha(\tilde x)-\alpha_0(\tilde{x})\|^2 \frac{K_{\alpha,2}}{\xi}\leq1/16$,\\
(vi) $|W(\alpha_0(X_i),Y_i)'v|\leq\bar{K}_w\|v\|_1$, $u'[\nabla_{\alpha}W(\alpha_0(X_i),Y_i)'r_0(X_i)][\nabla_{\alpha}W(\alpha_0(X_i),Y_i)'r_0(X_i)]'u\leq K_{\alpha,2}\|u\|^2$, ${\displaystyle\max_{j\in [p], k\in[M]}}|v'\nabla_{x}^2 r_{0j}(x)v|+ |v'\nabla^2\alpha_{0k}(x)v|\leq \bar{K}_{x,2}\|v\|^2$.
\end{assumption}
The adversarial sets $\mathcal{A}(\hat{\alpha})$, $\mathcal{A}_x(\hat{\alpha})$ constructed in Section~\ref{sec:marginals} make conditions (1), (i) hold with high probability. Without loss of generality, we next focus on discussing the conditions (ii) -- (vi) for the with-covariate case above. Given any $i\in[N],j\in[p],k\in[M]$, $\nabla_{\alpha_k}W_j(\alpha(X_i),Y_i)=\mathbf{1}\{k\in\ell(j)\}W_j(\alpha(X_i),Y_i)\frac{\partial_{\alpha_k}z_k(Y_i,\alpha(X_i))}{z_k(Y_i,\alpha(X_i))}$, where $\frac{\partial_{\alpha_k}z_k(Y_i,\alpha(X_i))}{z_k(Y_i,\alpha(X_i))}=\frac{1-2Y_{ik}}{2\alpha_{k}(X_i)(1-\alpha_{k}(X_i))}$. Additionally, note that $\nabla_{\alpha}^2W_j(\tilde{\alpha}(X_i),Y_i)=W_j(\tilde{\alpha}(X_i),Y_i)[g_{ij}g_{ij}'+D_{ij}]$, where $g_{ij}\in\mathbb{R}^M$ and $g_{ijk}=\frac{\mathbf{1}\{k\in\ell(j)\}(1-2Y_{ik})}{2\tilde{\alpha}_k(X_i)(1-\tilde{\alpha}_k(X_i))}$, $D_{ij}\in\mathbb{R}^{M\times M}$ is a diagonal matrix with $k$-th diagonal entry $-\frac{\mathbf{1}\{k\in\ell(j)\}(1-2Y_{ik})(1-2\tilde{\alpha}_k(X_i))}{2\tilde{\alpha}_k(X_i)^2(1-\tilde{\alpha}_k(X_i))^2}$. Particularly, when $\alpha_{0k}(X_i)$ are bounded away from both $0$ and $1$, since $\tilde{\alpha}\in\mathcal{A}_x(\hat{\alpha})$ are close to $\alpha_0(\cdot)$ with high probability by the construction of $\mathcal{A}_x(\hat{\alpha})$, $\tilde{\alpha}\in\mathcal{A}_x(\hat{\alpha})$ are also bounded away from both $0$ and $1$ with high probability, then condition (vi) implies that conditions (ii) -- (iv) hold. The theoretical analysis later implies ${\displaystyle \max_{\|\tilde x-x\|\leq h, \tilde\alpha \in \mathcal{A}_x(\hat\alpha)}}M^{1/2}\|\tilde \alpha(\tilde x)-\alpha_0(\tilde x)\|$ goes to zero as $N\rightarrow\infty$. So (v) holds when $N$ is sufficiently large. The last condition also includes a general smoothness assumption on $r_0(\cdot)$ and $\alpha_0(\cdot)$. 

Next, we impose the restricted eigenvalue condition, a standard assumption in high-dimensional regularized estimation that ensures sufficient curvature on a restricted set induced by the regularization, even when the design matrix may be rank-deficient globally. For a constant $c>1$, let $\bar c:=2\frac{c+1}{c-1}$. The restricted eigenvalues for both no-covariate and with-covariate cases are
$$\begin{array}{rl}
\kappa_{\bar{c}}^2 &\displaystyle := \min_{\|v_{T^c}\|_{w,1} \leq \bar{c}\|v_T\|_{w,1}} \frac{1}{N}\sum_{i=1}^N \frac{\{W(\alpha_0,Y_i)'v\}^2}{\{1+W(\alpha_0,Y_i)'r_0\}^2\|v_T\|_{w,2}^2},\\ 
\kappa_{x,\bar c}^2 &\displaystyle := \min_{\|v_{T_x^c}\|_{w_x,h,1} \leq \bar{c}\|v_{T_x}\|_{w_x,h,1}} \frac{1}{N}\sum_{i=1}^N K_h(X_i-x) \frac{\{W(\alpha_0(X_i),Y_i)'(v_a+v_b(X_i-x))\}^2}{ \{1+W(\alpha_0(X_i),Y_i)'r_0(X_i)\}^2\|(v_a,v_b)_{T_x}\|_{w_x,h,2}^2},
\end{array}$$
where $T$ is the support of $r_0$ in the no-covariate case and $T_x$ is the support of $(r_0(x),\nabla r_0(x))$ in the covariate-case. In the definition for $\kappa_{\bar{c}}^2$, $v\in\mathbb{R}^p$, $\|v_T\|_{w,k}= (\sum_{j\in T} w_j^k|v_j|^k)^{1/k}$. For $v=(v_a',v_b')'$ in the definition for $\kappa_{x,\bar{c}}^2$, $v_a\in\mathbb{R}^p$, $v_b\in\mathbb{R}^{p\times d}$. $\|(v_a,v_b)_{T_x}\|_{w_x,h,k}=\left(\sum_{j\in T_x}w_{xj}^k|v_{aj}|^k+h^k\sum_{(j,l)\in T_x}w_{xjl}^k|v_{b,jl}|^k\right)^{1/k}$. The restricted eigenvalue condition is a standard assumption for Lasso-type estimators in high-dimensional statistics \citep{bickel2009simultaneous,buhlmann2011statistics}. However, new technical challenges arise in our analysis compared to the classical Lasso setting. For example, we need to account for the non-separability between marginal probabilities and the binary components, as manifested in the structure of the design matrix. 

\subsection{Rates of Convergence without Covariates}\label{Sec:Rates:no-covariates}
We now provide the main theoretical results for the no-covariate case. First, we establish the convergence rate of the plug-in estimator $\hat{r}^{PI}$. We first define $\displaystyle\beta_N := \max_{j\in[p]}\frac{1}{w_j}\frac{1}{N}\sum_{i=1}^N \frac{|W_j(\alpha_0,Y_i)|}{1+W(\alpha_0,Y_i)'r_0}$. 
\begin{proposition}[Plug-in estimator]\label{prop:plug-in:no-covariate}
Suppose that Assumptions \ref{ass:iid-min-den}, \ref{ass:moment-assumptions}, \ref{ass:smoothness-assumptions} hold. Suppose for some constant $C'$, we set $\lambda$ such that for $\delta>1/2^M$, with probability at least $1-\delta$, 
$$\frac{\lambda}{c}\geq\frac{C'\Phi^{-1}(1-\delta/2p)}{\sqrt{N}}\max_{j\in[p]} \frac{1}{w_j}\left\{\frac{1}{N}\sum_{i=1}^N \frac{W_j^2(\alpha_0,Y_i)}{(1+W(\alpha_0,Y_i)'r_0)^2}\right\}^{1/2}\!\!\!\!+\frac{C'K_{\alpha,1}(1+\beta_N)\sqrt{M}\|\hat\alpha-\alpha_0\|}{\xi},$$
and $sM^{1/2}K_{\alpha,1}\|\hat\alpha-\alpha_0\|/ \{\xi\kappa_{\bar c}^2\min_{j\in[p]}w_j^2\}\leq \delta_N$, where $\delta_N$ is defined in Assumption~\ref{ass:moment-assumptions}. Suppose $\kappa_{\bar{c}}^2$ is bounded away from zero and $\bar K_2 + K_{\alpha,1} + K_{\alpha,2}/\xi^2 \leq C^\prime$ for some constant $C'$. Then $$\displaystyle\sqrt{\frac{1}{N}\sum_{i=1}^N \frac{\{W(\alpha_0,Y_i)'(\hat r^{PI} - r_0 )\}^2}{(1+W(\alpha_0,Y_i)'r_0)^2}}\leq C\sqrt{\frac{s(M+\log(1/\delta))}{N}} + C\sqrt{\frac{sM^2}{N}}$$ 
holds with probability $1-C\delta-C\delta_\alpha-C\epsilon$ for $N$ sufficiently large, where constant $C$ is independent of $N,s,M$.
\end{proposition}
The choice of penalty parameters $\lambda$ for Proposition~\ref{prop:plug-in:no-covariate} and later in Theorem~\ref{thm:main:FO} follows from the results of self-normalized moderate deviation theory. Example~\ref{ex:plug-in} in Appendix~\ref{appendix:no-covariate} of the ``electronic companions'' demonstrates the sharpness of the rate established for $\hat{r}^{PI}$ in Proposition~\ref{prop:plug-in:no-covariate}.
The rate of convergence for the plug-in estimator immediately implies the convergence rate for the \textit{oracle (plug-in) estimator} with true marginal probabilities $\alpha_0$:
\begin{proposition}[Oracle (plug-in) estimator]\label{prop:PI:infeasible}
Suppose $\hat{\alpha}=\alpha_0$, then under the conditions of Proposition \ref{prop:plug-in:no-covariate}, $\displaystyle\sqrt{\frac{1}{N}\sum_{i=1}^N \frac{\{W(\alpha_0,Y_i)'(\hat r^{PI} - r_0 )\}^2}{(1+W(\alpha_0,Y_i)'r_0)^2}} \leq  C\sqrt{\frac{sM+s\log(1/\delta)}{N}}$ holds with probability $1-C\delta-C\delta_\alpha-C\epsilon$, where $C$ is a constant independent of $N,s,M$.
\end{proposition}
Proposition~\ref{prop:plug-in:no-covariate} implies that the rate of convergence for $\hat{r}^{PI}$ is $\mathrm{O}_{\mathrm{p}}(\sqrt{sM^2/N})=\mathrm{O}_{\mathrm{p}}(\log p\sqrt{s/N})$ if $s\geq1$ and $M^2>\log(1/\delta)$, which is slower than that of $\hat{r}^{FO}$ as shown by Theorem~\ref{thm:main:FO} in the below:
\begin{theorem}[First-order estimator]\label{thm:main:FO}
Suppose Assumptions \ref{ass:iid-min-den}, \ref{ass:moment-assumptions}, \ref{ass:smoothness-assumptions} hold. Suppose for some constant $C'>1$, with probability at least $1-\delta$ we have\\
(i) $\lambda/(2c) \geq C'N^{-1/2} \Phi^{-1}(1-\delta/2p)\max_{j\in[p]} \frac{1}{w_j}\left\{\frac{1}{N}\sum_{i=1}^N \frac{W_j^2(\alpha_0,Y_i)}{(1+W(\alpha_0,Y_i)'r_0)^2}\right\}^{1/2}$, \\
(ii) $\lambda/(2c) \geq C'(1+\beta_N)\frac{K_{\alpha,2}}{\xi}\norm{\hat{\alpha}-\alpha_0}^2$,\\
(iii) $s\|\hat\alpha-\alpha_0\|^2 K_{\alpha,2} \bar K_3 / \{\xi\kappa_{\bar c}^2 {\displaystyle \min_{j\in[p]}} w_j^2\}\leq \delta_N$. \\
Then for $N$ sufficiently large, with probability $1- C\delta - C\delta_\alpha-C\epsilon$, $$\displaystyle\sqrt{\frac{1}{N}\sum_{i=1}^N \frac{\{W(\alpha_0,Y_i)'(\hat r^{FO} - r_0 )\}^2}{(1+W(\alpha_0,Y_i)'r_0)^2}}\leq C\lambda\sqrt{s}\frac{(1+1/c)}{\kappa_{\bar{c}}}+C\widetilde{\mathcal{R}}^{1/2}(\hat\alpha),$$ 
where $\widetilde{\mathcal{R}}(\hat\alpha) \leq C\bar{K}_2\sqrt{\frac{\log (M/\delta)}{N}}{\displaystyle\max_{\alpha \in \mathcal{A}(\hat\alpha)}}\|\alpha-\alpha_0\|_1  + C\frac{K_{\alpha,2}}{\xi^2}{\displaystyle\max_{\alpha \in \mathcal{A}(\hat\alpha)}}\|\alpha - \alpha_0\|^2$ and $C$ is independent of $s,N,M$.
\end{theorem}
The self-normalized moderate deviation theory essentially requires $\lambda$ to dominate the maximum norm of the score. The maximum norm of the score for the plug-in estimator is $\left\|\frac{1}{N}\sum_{i=1}^N\frac{W(\hat{\alpha},Y_i)}{1+W(\hat\alpha,Y_i)'r_0}\right\|_\infty=\left\|\frac{1}{N}\sum_{i=1}^N\frac{W(\alpha_0,Y_i)}{1+W(\alpha_0,Y_i)'r_0}\right\|_\infty+\mathrm{O}_p(\|\hat{\alpha}-\alpha_0\|_1)$, where $\left\|\frac{1}{N}\sum_{i=1}^N\frac{W(\alpha_0,Y_i)}{1+W(\alpha_0,Y_i)'r_0}\right\|_\infty=\mathrm{O}_p(\sqrt{\frac{\log(p)}{N}})=\mathrm{O}_p(\sqrt{\frac{M}{N}})$ and $\mathrm{O}_p(\|\hat{\alpha}-\alpha_0\|_1)=\mathrm{O}_p(\frac{M}{\sqrt{N}})$, we need to choose $\lambda\asymp\mathrm{O}(M/\sqrt{N})$ for $\hat{r}^{PI}$. For the first-order estimator, $\left\|\frac{1}{N}\sum_{i=1}^N\frac{W^{FO}(\alpha_0,Y_i)}{1+W^{FO}(\alpha_0,Y_i)'r_0}\right\|_\infty=\left\|\frac{1}{N}\sum_{i=1}^N\frac{W(\alpha_0,Y_i)}{1+W(\alpha_0,Y_i)'r_0}\right\|_\infty+\mathrm{O}_p(\|\hat{\alpha}-\alpha_0\|^2)$, where $W^{FO}(\alpha_0,Y_i)=W(\hat\alpha,Y_i)+\nabla_{\alpha} W(\hat\alpha,Y_i)(\alpha_0-\hat\alpha)$, which is a first-order expansion of $W(\alpha,Y_i)$ around $\hat{\alpha}$, so it introduces only a second-order bias in the score. This bias is smaller than the bias induced by directly plugging $\hat{\alpha}$ into the score. The penalty choice for $\hat{r}^{FO}$ is then $\lambda\asymp\mathrm{O}(\sqrt{M/N})$. By incorporating adversarial selection of $\alpha$, the first-order estimator regularizes the dependence of the score on $\hat{\alpha}$ and mitigates overfitting. As a result, a smaller penalty level $\lambda$ suffices compared with the plug-in estimator, leading to a lower estimation error. Moreover, the formulation of $\hat{r}^{FO}$ circumvents the non-concavity of the maximum likelihood objective and other computational challenges inherent in adversarial estimation of the marginal probabilities. 

The following corollary specializes Theorem~\ref{thm:main:FO} to a representative case, providing a precise convergence rate under the specific choice of the set $\mathcal{A}(\hat{\alpha})$ defined in~\eqref{def:A:no-covariate}, which admits efficient computation. The details of the construction of $\mathcal{A}(\hat\alpha)$ are in Section \ref{sec:marginals}.
\begin{corollary}\label{cor:FO}
    Under the conditions of Theorem \ref{thm:main:FO}, suppose $\kappa_{\bar{c}}$ is bounded away from zero and $\bar K_2 + K_{\alpha,1} + K_{\alpha,2}/\xi^2 \leq C^\prime$ for some constant $C'$, and $\mathcal{A}(\hat{\alpha})$ is constructed as~\eqref{def:A:no-covariate}. Then $$\displaystyle\sqrt{\frac{1}{N}\sum_{i=1}^N \frac{\{W(\alpha_0,Y_i)'(\hat r^{FO} - r_0 )\}^2}{(1+W(\alpha_0,Y_i)'r_0)^2}} \leq  C\sqrt{\frac{s(M+\log(1/\delta))}{N}} + C\sqrt{\frac{M\log(M/\delta)}{N}}$$ 
    holds with probability $1-C\delta-C\delta_\alpha-C\epsilon$, where $\delta\leq\delta_{\alpha}$, and $C$ is independent of $s,N,M$. \proofend
\end{corollary}
Corollary \ref{cor:FO} shows that the impact of estimated marginal probabilities on the first-order estimator is captured by $\widetilde{\mathcal{R}}(\hat\alpha)=\mathrm{O}_{\mathrm{p}}\left(\frac{M\log M}{N}\right)$ defined as in Theorem~\ref{thm:main:FO}. In many regimes of interest (e.g., $s \gg \log M$), this term is expected to be small as the adversarial estimation is over a nuisance parameter with dimension $M \ll p$. Note that the $\log M$ factor arises from the $\ell_\infty$ construction of $\mathcal{A}(\hat{\alpha})$ in~\eqref{def:A:no-covariate}, which requires controlling the maximum deviation across the $M$ marginal probability estimators. This factor could be avoided by instead defining $\mathcal{A}(\hat{\alpha})$ using an $\ell_2$ bound on the marginal estimation error (i.e., Euclidean balls). However, from a computational standpoint, representing the adversarial set as a hyper-rectangle is considerably more tractable, because the minimum of the inner optimization problem in~\eqref{eq:PMLE-no-covariates-bilinear} is attained at one of its finitely many extreme points. Since $M \asymp \log p$, Corollary~\ref{cor:FO} implies that the first-order estimator attains a convergence rate of order $\sqrt{sM/N}+\sqrt{M\log(M)/N}$. The rate for the oracle estimator with $\alpha_0$ known is of order $\sqrt{sM/N}$. So the first-order estimator is rate optimal when $s\gtrsim\log(M)$. Table \ref{table:comparison} compares all the aforementioned estimators when $s\gtrsim\log M$, showing the trade off between rates of convergence and computational complexity. 
\begin{table}[htp!]
\centering
\caption{Comparison of estimators when $s\gtrsim\log(M)$.}
\begin{tabular}{p{2.0cm}p{2.0cm}p{2.cm}p{2.8cm}p{4.8cm}}
\toprule
Estimator & Order of $\lambda$ & Rate & Optimize over $r$ & Evaluation over $\alpha$ \\
\midrule
MLE & $0$ & $\sqrt{p/N}$ & concave & non-concave maximization \\
Plug-in ($\hat{\alpha}$) & $\log p/\sqrt{N}$ & $\sqrt{s/N}\log p$ & concave & direct \\
Adversarial & $\sqrt{\log p/N}$ & $\sqrt{s\log p/N}$ & concave & non-convex minimization \\
First-order & $\sqrt{\log p/N}$ & $\sqrt{s\log p/N}$ & concave & tractable minimization \\
Oracle ($\alpha_0$) & $\sqrt{\log p/N}$ & $\sqrt{s\log p/N}$ & concave & direct \\
\bottomrule
\end{tabular}
\begin{tablenotes}
\small
\item[a] All estimators maximize a concave function in $r$, but differ in how the nuisance parameters $\alpha$ affect the penalty choice and evaluation over $\alpha$.
\end{tablenotes}
\label{table:comparison}
\end{table}

\subsection{Rates of Convergence with Covariates}\label{Sec:Rates:covariates}
The intuition from the no-covariate case extends to the with-covariate case. We therefore focus on the localized first-order estimator $\hat{r}^{FO}(\cdot,x)$. The analysis builds upon the no-covariate case but must additionally control approximation errors from the local approximation of $r_0(X)$ and $\alpha_0(X)$. Constructing the adversarial set pointwise in $X_i$ would introduce non-negligible bias. To address these challenges, we focus on a class of local linear functions for estimating the marginal probabilities. The next result provides rates of convergence for the generalized correlation coefficients in the with-covariate case. Define $\beta_{N,x}:={\displaystyle\max_{j\in[p],k\in[d]}}\left\{\frac{1}{w_{xj}}\frac{1}{N}\sum_{i=1}^N\frac{K_h(X_i-x)|W_j(\alpha_0(X_i),Y_i)|}{1+W(\alpha_0(X_i),Y_i)^\prime r_0(X_i)},\frac{1}{w_{xjk}}\frac{1}{N}\sum_{i=1}^N\frac{K_h(X_i-x)|W_j(\alpha_0(X_i),Y_i)(X_{ik}-x_k)/h|}{1+W(\alpha_0(X_i),Y_i)^\prime r_0(X_i)}\right\}$.
\begin{theorem}\label{thm:main:covariate}
Suppose Assumptions \ref{ass:iid-min-den}, \ref{ass:moment-assumptions}, \ref{ass:smoothness-assumptions} hold. Moreover suppose that for some $C'>1$, the conditions below hold with probability at least $1-\delta$, where $\delta\geq\max\{\exp(-C'\{\sqrt{Nh^d\log M}+Nh^d\}),\frac{\log(M)}{\delta_NNh^d},\frac{1}{2^M}\}$,
$$\begin{array}{rl}
\lambda_x/(3c) &\displaystyle \geq \frac{C'\Phi^{-1}(1-\delta/2p(1+d))}{(Nh^d)^{1/2}}{\displaystyle \max_{j\in [p]}} \frac{1}{w_{xj}} \sqrt{\frac{1}{N}\sum_{i=1}^N\frac{K_h(X_i-x)W_j^2(\alpha_0(X_i),Y_i)}{(1+W(\alpha_0(X_i),Y_i)'r_0(X_i))^2}},\\
\lambda_x/(3c) &\displaystyle \geq \frac{C'\Phi^{-1}(1-\delta/2p(1+d))}{(Nh^d)^{1/2} }{\displaystyle \max_{j\in [p], k\in[d]}} \frac{1}{w_{xjk}} \sqrt{\frac{1}{N}\sum_{i=1}^N\frac{K_h(X_i-x)W_j(\alpha_0(X_i),Y_i)^2\big(\frac{X_{ik}-x_k}{h}\big)^2}{(1+W(\alpha_0(X_i),Y_i)'r_0(X_i))^2}},\\
\lambda_x/(3c) &\displaystyle \geq C'(1+\beta_{N,x})K_{\alpha,2}\max_{\tilde \alpha\in\mathcal{A}_x(\hat\alpha), \|X_i-x\|\leq h}\norm{\tilde{\alpha}(X_i)-\alpha_0(X_i)}^2/\xi+C'h^2,\\
\delta_N &\displaystyle \geq \max_{\tilde \alpha\in\mathcal{A}_x(\hat\alpha), \|X_i-x\|\leq h}\|\tilde\alpha(X_i)-\alpha_0(X_i)\|^2s/ \{\xi\kappa_{x,\bar c}^2 \min_{j\in[p], k\in[d]} \{w_{xjk}^2,w_{xj}^2\}\},\\
\delta_N&\displaystyle \geq \max_{\tilde \alpha\in\mathcal{A}_x(\hat\alpha), \|X_i-x\|\leq h}\|\tilde\alpha(X_i)-\alpha_0(X_i)\|/ \{\xi\min_{j\in[p], k\in[d]}\{w_{xjk},w_{xj}\}\}.
\end{array}$$ 
Then with probability $1-C\delta$ we have
{$$ \sqrt{\frac{1}{N}\sum_{i=1}^N \frac{K_h(X_i-x)\{W(\alpha_0(X_i),Y_i)'(\hat r^{FO}(X_i,x) - r_{*x}(X_i) )\}^2}{(1+W(\alpha_0(X_i),Y_i)'r_0(X_i))^2}} \leq C\frac{\lambda_x\sqrt{s}}{\kappa_{x,\bar{c}}} +  \widetilde{\mathcal{R}}^{1/2}(\hat\alpha),$$}
where with probability $1-C\delta-C\delta_\alpha-\frac{C(\bar{K}_{\max}/\bar{K}_2)^2}{Nh^d}-C\epsilon$,
$$\begin{array}{rl}
\widetilde{\mathcal{R}}(\hat\alpha) &\displaystyle\leq C\max_{\|X_i-x\|\leq h,\alpha \in \mathcal{A}_x(\hat\alpha)}\|\alpha(X_i)-\alpha_0(X_i)\|_1 \sqrt{\frac{\log(M/\delta)}{Nh^d}}+ Csdh^2\sqrt{\frac{\log(1/\delta)}{Nh^d}}+Cs^2d^2h^4\\
&\displaystyle\quad + C\max_{\|X_i-x\|\leq h,\alpha \in \mathcal{A}_x(\hat\alpha)} \|\alpha(X_i)-\alpha_0(X_i)\|^2+Csdh^2\left(\frac{Md}{\sqrt{Nh^{d}}}+Mdh^2\right),
\end{array}$$
and $C$ is a universal constant independent of $s,M,N,h,d$. 
\end{theorem}
Assumption \ref{ass:local-linear} in the below is standard in the literature of local-linear regression \citep[e.g.][]{ruppert1994multivariate,fan1993local,lu1996multivariate, fan2003nonlinear}:
\begin{assumption}\label{ass:local-linear}
(i)$\int uu^TK(u)du=\mu_2\mathbf{I}_{d}$, where $\mu_2\neq0$ is a scalar and $\mathbf{I}_{d}$ is the $d\times d$ identity matrix. (ii) $X_i$'s have common density $\rho$ with compact support $\mathcal{X}\subseteq\mathbb{R}^d$ such that $\rho(\cdot)\geq\rho_0>0$ on its support (where $\rho_0$ is a constant) and $\rho$ is continuously differentiable for any $x$ in the interior of $\mathcal X$. (iii) For any $j\in[M]$, $\alpha_j(x)$ is continuously differentiable, and its second-order derivatives are continuous. (vi)  $K(\cdot): \mathbb{R}^d\mapsto\mathbb{R}$ is supported on $\{u:\|u\|\leq1\}$. For any $\norm{z}_2\leq1$, $K(-z)=K(z)$ and $\underline{\kappa}\leq K(z)\leq\mathcal{K}$, where $\underline{\kappa}$, $\mathcal{K}$ are positive absolute constants.
\end{assumption}
Assumption~\ref{ass:local-linear} is satisfied by many commonly used kernels, including the uniform (box) kernel $K(u)=\mathbf{1}\{\|u\|\leq1\}$ and the quadratic-positive kernel $K(u)=\mu_0(1+\mu_1\|u\|^2)\mathbf{1}\{\|u\|\leq1\}$. Similar properties hold for other compact-support kernels commonly used in nonparametric estimation. 
\begin{corollary}\label{cor:FO:covariate}
Suppose the conditions of Theorem \ref{thm:main:covariate} and  Assumption \ref{ass:local-linear} hold. Suppose $\kappa_{x,\bar{c}}$ is bounded away from zero,  $\bar{K}_2 + K_{\alpha,1} + K_{\alpha,2}/\xi^2 \leq C'$ for a universal constant $C'$. Let the adversarial set $\mathcal{A}_x(\hat{\alpha})$ be constructed as in Section \ref{sec:marginals}. Define $\displaystyle\Delta_{x,i}^{FO}:=\left\{\frac{W(\alpha_0(X_i),Y_i)'(\hat r^{FO}(X_i,x) - r_{*x}(X_i))}{1+W(\alpha_0(X_i),Y_i)'r_0(X_i)}\right\}^2$. Then with probability $1-C\delta-C\delta_\alpha-C\epsilon-C(\bar{K}_{\max}/\bar{K}_2)^2/(Nh^d)$, where $\delta\leq\delta_{\alpha}$, we have
\begin{equation}\label{eq:cor:cov:bound-1}
\begin{array}{rl}
&\displaystyle\quad\sqrt{\frac{1}{N}\sum_{i=1}^N K_h(X_i-x)\Delta_{x,i}^{FO}}\\
&\displaystyle\leq C\sqrt{\frac{sM+s\log(1/\delta)}{Nh^d}}+C(s+\sqrt{M})dh^2,
\end{array}
\end{equation}
where $C$ is a universal constant independent of $s,N,M,h,d$. Henceforth, the optimal bandwidth is $h^*=\min\left\{\sqrt{\frac{\delta_N\xi}{sd}},\max\left\{\left(\frac{M^3}{N\delta_N^2}\right)^{1/d},4^{-\frac{2}{d+4}}\left(\frac{sM+s\log(1/\delta)}{N(s+\sqrt{M})^2}\right)^{\frac{1}{d+4}}\right\}\right\}$.
Particularly, when $h=h^*$, with probability $1-\delta$, we have 
\begin{equation}\label{eq:bound:optimal-choice-h}
\sqrt{\frac{1}{N}\sum_{i=1}^N K_h(X_i-x)\Delta_{x,i}^{FO}}\leq\bar{C}\min\left\{B_{1},B_2\right\},
\end{equation}
where $B_1:=\delta_N\sqrt{\frac{s(M+\log(1/\delta))}{M^3}}+(\sqrt{M}+s)d\left(\frac{M^3}{N\delta_N^2}\right)^{\frac{2}{d}}$,
$B_2:=\sqrt{\frac{s(M+\log(1/\delta))}{N}}\left(\frac{sd}{\delta_N\xi}\right)^{d/4}\!\!+\frac{(\sqrt{M}+s)\delta_N\xi}{s}$, and $\bar{C}$ is a constant independent of $s,N,M,h,d$.\proofend
\end{corollary}

\section{Construction of the Adversarial Sets}\label{sec:marginals}
 This section details the construction of the adversarial sets for the no-covariate and with-covariate cases. The key idea is to construct $\mathcal{A}(\hat{\alpha})$ and $\mathcal{A}_x(\hat{\alpha})$ as hyper-rectangles, so that the minimum of the inner optimization problem with respect to $\alpha$ in the first order estimators is attained at one of their finitely many extreme points, thereby enabling efficient computation. Recall that for the no-covariate case, we consider the sample average estimators \eqref{eq:SAA} as the initial estimator for the marginal probabilities. We construct the adversarial set in the no-covariate case as
\begin{equation}\label{def:A:no-covariate}
\mathcal{A}(\hat\alpha) :=\prod_{j=1}^M [\hat\alpha_j - cv_{\delta_\alpha} \hat s_j, \hat\alpha_j + cv_{\delta_\alpha} \hat s_j],  
\end{equation}
where $\hat s_j= \sqrt{\hat\alpha_j(1-\hat\alpha_j)}$ and $cv_{\delta_\alpha} := (1-\delta_\alpha)\mbox{-quantile of} \max_{j\in[M]} \left| \frac{1}{N}\sum_{i=1}^N \xi_i \frac{(Y_{ij} - \hat \alpha_j)}{\sqrt{\hat\alpha_j(1-\hat\alpha_j)}}\right|$ conditioning on $\{Y_i\}_{i\in[N]}$ for $\delta_\alpha \in (0,1)$, and $\{\xi_i\}_{i\in[N]}$ are i.i.d. $\mathcal{N}(0,1)$ random variables. For the with-covariate case, recall that we use local linear regression estimators for the marginal probabilities. Specifically, we compute $(\tilde{\alpha}_{j}(x), \tilde\beta_j(x)) \in \arg\min_{(\alpha_j,\beta_j)}\sum_{i=1}^N K_h(X_i-x)(Y_{ij} - \alpha_j - \beta_j(X_i-x) )^2$ for any $j\in[M]$. Since $\alpha_0(x)\in[\eta,1-\eta]$ by Assumption~\ref{ass:iid-min-den}, we truncate $\tilde{\alpha}_{j}(x)$ and set
\begin{equation}\label{eq:truncate-local-linear}
\hat{\alpha}_j(x):=\min\{\max\{\eta,  \tilde{\alpha}_j(x)\} , 1-\eta\},\ \ \forall j\in[M].
\end{equation}
We construct $\mathcal{A}_x(\hat\alpha)$ as a set of linear functions to contain the function $\alpha_{*x}(X_i):=\alpha_0(x)+\nabla_x \alpha_0(x)(X_i-x)$ with high probability. Specifically,
$$\mathcal{A}_x(\hat\alpha) :=\prod_{j=1}^M \left\{ \!\alpha_j+\beta_j(X-x) :\!\!\! \begin{array}{l}\alpha_j \in [\hat \alpha_j(x) - cv_x \hat s_{j,0}(x), \hat\alpha_j(x) + cv_x \hat s_{j,0}(x)],  \\ 
\beta_{jk} \in [\tilde \beta_{jk}(x) - cv_x \hat s_{j,k}(x), \tilde \beta_{jk}(x) + cv_x \hat s_{j,k}(x)], k \in [d] 
\end{array}\!\!\right\},$$
and conditional on $\{Y_i,X_i\}_{i=1\in[N]}$, $cv_x := (1-\delta_\alpha)\mbox{-quantile of } {\displaystyle\max_{ j \in  [M], \ell \in\{0\}\cup[d]}}\displaystyle\left|\sum_{i=1}^N \xi_i \frac{\hat \epsilon_{ij} \hat{\zeta}_{i\ell}}{\hat s_{j,\ell}(x)}\right|$, where $\{\xi_i\}_{i\in[N]}$ are i.i.d. $\mathcal{N}(0,1)$ random variables, $\hat\epsilon_{ij}:=Y_{ij} - \hat{\alpha}_j(x)-\tilde\beta_j(x)(X_i-x)$, $\hat{\zeta}_{i} := (Z'W_xZ)^{-1}Z_i K_h(X_i-x)$  with $Z_i=(1, (X_i-x)')'$ and $W_{x}$ is a diagonal $N\times N$ matrix with its $i$-th diagonal entry equal to $W_{x,ii}:=K_h(X_i-x)$, $\hat s_{j,\ell}(x) := \sqrt{\sum_{i=1}^N \hat \epsilon_{ij}^2 \hat{\zeta}_{i\ell}^2}$, and $\hat{\zeta}_{i\ell}$ is the $\ell$-th component of vector $\hat{\zeta}_i$, $\forall \ell\in\{0\}\cup[d]$.

\section{Application to Estimating the Average Treatment Effects under Multiple Binary Treatments}\label{section:causal}
In this section, we apply our theoretical results to the estimation of average treatment effects in the presence of $M$ binary treatments. We consider a standard potential outcome framework \citep{imbens2015causal}. For a treatment vector $t \in \{0,1\}^M$, the $j$th treatment is treated if and only if $t_j = 1$. We observe the triplets $Z_i=(O_i,T_i,X_i)$ across $i\in[N]$, where $O_i\in\mathbb{R}$ denotes the outcome of interest, $T_i \in \{0,1\}^M$ denotes the treatment assortments, and $X_i \in \mathbb{R}^d$ is the covariate. We are interested in estimating the \textit{average treatment effect} (ATE) given two treatment levels $t$ and $t'$, namely $\tau(t,t') = \mathbb{E}[ O(t) - O(t')]$, where $O(t)$ denotes the potential outcome of receiving the treatment $t\in \{0,1\}^M$. We consider the cross-fitted \textit{augmented inverse probability weighted} (AIPW) estimator \citep[e.g.][]{robins1994estimation,robins1995semiparametric,hahn1998role,scharfstein1999adjusting,kennedy2020towards,newey2018cross}. 
Specifically, we split the data into $K$ disjoint folds of size $N/K$ and we use $\mathcal{I}_k$ to denote the $k$-th fold, for $k\in[K]$. It is not required that all the folds have the same size, but we proceed with this merely for simplicity, without loss of generality. For any $t\in\{0,1\}^M$, we let $\hat{\mu}_{(t)}^{-k}(X)$ and $\hat e_t^{-k}(X)$ denote the estimators for $\mu_{(t)}(X)=\mathbb{E}[O|T=t,X]$ and $e_t(X)=\mathbb{P}(T=t|X)$ without using observations from fold $k$. Then we estimate the average treatment effect using
$$\hat \tau(t,t')=\!\frac{1}{K}\!\!\sum_{k=1}^K\frac{1}{|\mathcal{I}_k|}\!\sum_{i\in\mathcal{I}_k}\!\left(\hat\mu_{(t)}^{-k}(X_i)\!-\!\hat \mu_{(t')}^{-k}(X_i)\!+\!\frac{1\{T_i=t\}}{\hat e_t^{-k}(X_i)}\!\!\left(O_i - \hat\mu_{(t)}^{-k}(X_i)\right) - \frac{1\{T_i=t'\}}{\hat e_{t'}^{-k}(X_i)}\!\!\left(O_i - \hat\mu_{(t')}^{-k}(X_i)\right)\right).$$
The generalized propensity scores are estimated based on (\ref{def:localized:alpha:estimator}) and the first-order estimator $\hat{r}^{FO}(\cdot,x)$ at the given covariate of interest $x$:
\begin{equation}\label{eq:est:GPS}
\hat e_t(x)= \Big(1+W(\hat\alpha(x),t)'\hat r^{FO}(x,x)\Big)\prod_{j=1}^M\hat\alpha_j(x)^{t_j}(1-\hat\alpha_j(x))^{1-t_j},
\end{equation}
We assume that the sample contains observations with treatment $t$ as well as observations receiving treatment $t'$. We impose the following regularity conditions:
\begin{assumption}\label{ass:causal:extra}
Suppose (i) $\{O_i(t)\}_{t\in\{0,1\}^M} \indep T_i \mid X_i$; (ii)  there is constant $c \in (0,1)$ such that  $c\prod_{j=1}^M \alpha_{0j}^{t_j}(x)(1-\alpha_{0j}(x))^{1-t_j} < e_t(x) < 1- c, \forall x \in \mathcal{X}$, for any $t\in\{0,1\}^M$. (iii) For any $t\in\{0,1\}^M$, $\sup_{x\in\mathcal{X}}|\hat\mu_{(t)}(x)|\leq\bar{\mu},\max_{i\in[N]}|O_i(t)|\leq\bar{\mu}$ for some absolute constant $\bar{\mu}>0$, $\mathbb{E}[\{\hat \mu_{(t)}(X_i) - \mu_{(t)}(X_i)\}^2]\leq\epsilon_N^2N^{-1+2\gamma}$, $\sqrt{\frac{1}{N}\sum_{i=1}^N \{\hat \mu_{(t)}(X_i) - \mu_{(t)}(X_i) \}^2} \leq \epsilon_N N ^{-1/2+\gamma}$ with probability $1-\epsilon_N$ 
for some $0<\gamma\leq1/4$, where $\epsilon_N\to0$.
\end{assumption}

Assumption \ref{ass:causal:extra} imposes the standard unconfoundness assumption and a version of the overlapping condition. Since we allow $M$ to go to infinity, the traditional overlapping condition for generalized propensity scores (i.e. $\eta<e_t(x)<1-\eta$ for $\forall x\in\mathcal{X}, t\in\{0,1\}^M$, where $\eta\in(0,1)$ is a constant) doesn't hold. Instead, we impose condition (ii) of Assumption \ref{ass:causal:extra}, which allows for some assortment of treatments to be common but no assortment is exceedingly rare relative to the case where treatments are assigned independently. Condition (iii) on the conditional mean estimation ensures sufficient smoothness. 
\begin{theorem}\label{thm:causal-inference:AIPW}
Suppose Assumptions \ref{ass:iid-min-den}, \ref{ass:moment-assumptions}, \ref{ass:smoothness-assumptions}, \ref{ass:local-linear}, \ref{ass:causal:extra} hold. Suppose $\kappa_{x,\bar{c}}$ is bounded away from zero, $\bar{K}_2 + K_{\alpha,1} + K_{\alpha,2}/\xi^2\leq C'$ and $\mathcal{A}_x(\hat{\alpha})$ is constructed as before. Additionally, suppose that $0<\gamma\leq\frac{1}{4}$, $\psi_NN^\gamma<1$, $\psi_NN^\gamma\to0$, and for some constant $C$,
\begin{equation}\label{consistency:ATE:condition}
\begin{array}{rl}   
    &\displaystyle\quad\sqrt{s}\left(\sqrt{\frac{sM}{Nh^d}}\!+(s+\sqrt{M})dh^2\right)\\
    &\displaystyle\leq\frac{C\psi_N}{\sup_{x\in\mathcal{X}}\!\prod_{j=1}^M\!\alpha_{0j}(x)^{t_j}(1-\alpha_{0j}(x))^{1-t_j}},
\end{array}
\end{equation} 
\begin{equation}\label{consistency:ATE:2}
Mdh^2+\frac{M}{\sqrt{Nh^d}}\leq C\psi_N,
\end{equation}
where $M=\mathrm{o}(\min\{\log(1/\psi_N),\log(1/\epsilon_N)\})$, then we have $\sqrt{N}\hat \Sigma^{-1/2} (\hat \tau(t,t') - \tau(t,t')) \to\mathcal{N}(0,1)$, where $\hat{\Sigma}$ is the semiparametric efficiency variance bound. \proofend
\end{theorem}
Theorem \ref{thm:causal-inference:AIPW} shows that independence among some treatments induces sparsity in $r_0(x)$ and substantially improves estimation efficiency when $s\ll 2^M$. 
Remark~\ref{rmk:causal:1} and Remark~\ref{rmk:causal:2} show the compatibility of the conditions in Theorem~\ref{thm:causal-inference:AIPW} and provide two examples of $(N,M,d,s)$ satisfying them.
\begin{remark}\label{rmk:causal:1}
According to condition (iii) of Assumption \ref{ass:causal:extra} and the conditions of Theorem \ref{thm:causal-inference:AIPW}, we establish that the cross-term, given by the product of the outcome and propensity score estimation errors, is $\mathrm{o}_{\mathrm{P}}(N^{-1/2})$. We allow the rate of convergence for the propensity score estimator to be slower than the classical $\mathrm{o}_{\mathrm{P}}(N^{-1/4})$ to adapt to the constraints with the covariate dimension $d$ for the kernel estimator. As a result, we require the rate of convergence for the outcome estimator to be faster than $\mathrm{o}_{\mathrm{P}}(N^{-1/4})$ in condition (iii) of Assumption \ref{ass:causal:extra}. 
This assumption holds in some general setup. For example, suppose $\hat{\mu}_{(t)}(\cdot)$ is a parametric outcome model that can be estimated at $\sqrt{N}$ rate consistently, then if $\epsilon_N^2N^{2\gamma}\geq\mathbb{E}[N\{\hat \mu_{(t)}(X_i) - \mu_{(t)}(X_i)\}^2]$, and it is sufficient to let $\epsilon_NN^{\gamma}\rightarrow\infty$. Further, by leveraging moderate deviation theory in this case, under proper moment conditions, with probability $1-\delta$, $\max_{t\in\{0,1\}^M}\sqrt{\frac{1}{N}\sum_{i=1}^N\{\hat{\mu}_{(t)}(X_i)-\mu_{(t)}(X_i)\}^2}\leq \frac{c'\log(p/\delta)}{\sqrt{N}}$, where $p=2^M$, $c'$ is some absolute constant depending on the corresponding moment conditions. Hence if additionally we have $M\leq\epsilon_NN^{\gamma},\ \mbox{where} \ 0<\gamma\leq\frac{1}{4}$, condition (iii) of Assumption \ref{ass:causal:extra} holds.
\proofend 
\end{remark}
\begin{remark}\label{rmk:causal:2}
For Assumption~\ref{ass:moment-assumptions}, all conditions of Theorem~\ref{thm:causal-inference:AIPW} and the sufficient conditions implied by Remark~\ref{rmk:causal:1} for (iii) of Assumption~\ref{ass:causal:extra} (namely $\epsilon_NN^{\gamma}\rightarrow\infty$, $M\leq\epsilon_NN^{\gamma}$) to hold, 
by direct calculation, we need $d<\frac{2}{\gamma}-4$.
Particularly, if $\gamma=\frac{1}{4}$, then the classical $\mathrm{o}_{\mathrm{P}}(N^{-1/4})$-consistency condition is implied, and we can only let $d\leq3$ when $\gamma=1/4$. As $\gamma$ decreases we allow higher values of covariate dimension. Recall that we assume the covariate dimension $d$ is low so that the kernel method is effective. Particularly, two possible choices are to both take $h\asymp\left(\frac{\log(NMd)}{Nd^2}\right)^{1/(d+4)}$, and set $s=\log M, M=\mathrm{o}\left(N^{2/(d+4)}\psi_N[\log(NMd)]^{-\frac{2}{d+4}}d^{-\frac{d}{d+4}}\right)$ or $s=\sqrt{M}, M=\mathrm{o}\big(N^{\frac{8}{5(d+4)}}\psi_N^{\frac{4}{5}}[d\log(NMd)]^{-\frac{8}{5(d+4)}}\big)$, then the conditions in Theorem \ref{thm:causal-inference:AIPW} and Assumption~\ref{ass:moment-assumptions} are all satisfied. \proofend
\end{remark}
Our result is related to the doubly robust kernel estimators. The condition in standard semiparametric settings requires $\mathrm{o}_{\mathrm{P}}(N^{-1/4})$ consistency (e.g. \citep{chernozhukov2017double}). For non-parametric methods, \citet{kennedy2017non} estimates continuous treatment effect and derives $\sqrt{N}$-local rate of convergence with a kernel smoothing approach, by requiring that $h=\mathrm{O}(N^{-1/5})$, and the outcome estimation and propensity score estimation are both required to be of order $\mathrm{o}_{\mathrm{P}}(N^{-1/5})$. \citet{rothe2019properties} proposes a semiparametric two-step doubly robust estimators by requiring $\mathrm{o}_{\mathrm{P}}(N^{-1/6})$ consistency for both outcome and propensity score estimation, as long as their product is of order $\mathrm{o}_{\mathrm{P}}(N^{-1/2})$. 
Our estimator allows a slower rate of convergence for estimating the generalized propensity scores by imposing additional sparsity assumptions with a prudent choice of $N,M,d,h$. 

\section{Numerical Simulations}\label{sec:simulations}
For the no-covariate case, we compare $\hat{r}^{PI}$ (\ref{eq:PMLE-no-covariates}) and $\hat{r}^{FO}$ (\ref{eq:PMLE-no-covariates-bilinear}) by running simulations with $s=2,M=4,N=500$ and $s=5, M=4, N=500$. We focus on two types of penalty weighting by specifying different values for $\{w_k\}_{k\in[2^M-M-1]}$: (I) $w_{k}=\{\frac{1}{N}\sum_{i=1}^NW_{k}(\hat{\alpha},Y_i)^2\}^{1/2}$; (II) $w_{k}=\{\frac{1}{N}\sum_{i=1}^N W_{k}^2(\hat{\alpha},Y_i)/(1+W_i^\prime r_0)^2\}^{1/2}$, where in type (\textrm{II}), we use the true correlation coefficient $r_0$. We denote the estimators corresponding to these two weightings choices as $w_{\textrm{I}}$ and $w_{\textrm{II}}$ respectively. PI (oracle) refers to plug-in estimator with true marginal probabilities $\alpha_0$. For any $y\in\{0,1\}^M$, let $p_y$ be the true probability of vector $y$, and $\hat{p}_y=
(1+W(\hat{\alpha},y)^\prime \hat{r})\prod_{j=1}^M\hat{\alpha}_j^{y_j}(1-\hat{\alpha}_j)^{1-y_j}$. We compute three metrics, \textit{root mean squared error} (RMSE) ($\mathbb{E}[\norm{\hat{r}-r_0}_2^2]^{1/2}$), \textit{maximum absolute probability estimation error} ($\mathbb{E}[\max_{y\in\{0,1\}^M}|\hat{p}_y-p_y|]$), and \textit{mean expected probability estimation error} ($\mathbb{E}[\sum_{y\in\{0,1\}^M}p_y|\hat{p}_y-p_y|]$) for the estimation of all $2^M$ probabilities. When $\hat{p}_y \notin [0,1]$, we truncate it to $[0,1]$. The expectations in the metrics are computed by averaging over $200$ independent replications under random environments. $\lambda_{cv}$ refers to the penalty chosen via cross-validation, and $\lambda_{\mathrm{theory}}$ denotes the theoretically derived penalty parameter, chosen at the scale prescribed by Theorem~\ref{thm:main:FO} for the first-order estimators and Proposition~\ref{prop:plug-in:no-covariate} for the plug-in estimators. Specifically, we choose 
$\lambda_{\mathrm{theory}}^{\mathrm{PI}}=\frac{\Phi^{-1}(1-\delta/2p)}{\sqrt{N}}+\frac{M}{\sqrt{N}}$, $\lambda_{\mathrm{theory}}^{\mathrm{FO}}=\frac{\Phi^{-1}(1-\delta/2p)}{\sqrt{N}}+\frac{M}{N}$ and $\lambda_{\mathrm{theory}}^{\mathrm{PI}_{\mathrm{oracle}}}=\frac{\Phi^{-1}(1-\delta/2p)}{\sqrt{N}}$. For the with-covariate case, we also focus on two types of penalty weighting $\{w_k\}_{k \in [2^M - M - 1]}$: type (I) $w_{k}=\sqrt{\sum_{i=1}^N\{\frac{K_h(X_i-x)}{N}W_{k}(\hat{\alpha}(X_i),Y_i)^2}\}$ and $w_{k,1}=\sqrt{\sum_{i=1}^N\frac{K_h(X_i-x)}{N}W_{k}(\hat{\alpha}(X_i),Y_i)^2(X_{i}-x)^2}$; type (II) $w_{k}=\sqrt{\sum_{i=1}^N\frac{K_h(X_i-x)W_{k}(\hat{\alpha}(X_i),Y_i)^2}{N[1+W(\hat{\alpha}(X_i),Y_i)^\prime r_0(X_i)]^2}}$ and 
$w_{k,1}=\sqrt{\sum_{i=1}^N\frac{K_h(X_i-x)W_{k}(\hat{\alpha}(X_i),Y_i)^2(X_{i}-x)^2}{N[1+W(\hat{\alpha}(X_i),Y_i)^\prime r_0(X_i)]^2}}$, where in type (\textrm{II}), we use the true correlation coefficient $r_0(X_i)$.
PI (oracle) refers to plug-in estimator with true marginal probabilities $\alpha_0(X_i)$. Similar to the no-covariate case, $\lambda_{\mathrm{theory}}$'s are specified as
$\lambda_{\mathrm{theory}}^{PI} = \frac{\Phi^{-1}(1 - \delta / 2p(d + 1))}{\sqrt{Nh^d}} + \frac{M}{\sqrt{Nh^d}} + dMh^2$,
$\lambda_{\mathrm{theory}}^{FO} = \frac{\Phi^{-1}(1 - \delta / 2p(d + 1))}{\sqrt{Nh^d}} + \frac{M}{Nh^d} + Md^2h^4$, and
$\lambda_{\mathrm{theory}}^{PI_{\mathrm{oracle}}} = \frac{\Phi^{-1}(1 - \delta / 2p(d + 1))}{\sqrt{Nh^d}}$. For the no-covariate case, when $s=2,M=4,N=500$, the marginal probability vector is set as $(0.613,0.491,0.653,0.510)$, $r_0 = (0,0,0,0,0, 0,0,-0.339,0,0,0.249)$. For $s=5, M=4, N=500$, the marginal probability vector is set as $(0.767,0.360,0.389,0.535)$, and $r_0=(0,-0.009,0,0,0.160,-0.011,0,0,0.031,0,-0.066)$. For the with-covariate case, We sample $N$ i.i.d. $X_i\sim\mathrm{Uniform}[0,1]$. For any $j\in[M]$, set $\theta=(0.1,0.2,0.3,0.4)$ and $\alpha_{0j}(X_i) = \frac{1}{1+e^{\theta_jX_i}}$ as the $j$-th marginal probability for $\alpha_0(X_i)$. We first choose $s$ entries of $r_0(\cdot)$ as its support. 
For each $i\in[N]$, $r_{0j}(X_i)=\{(-1)^j0.03jX_i\}$ if $j$ is in the support. 
Here we provide the histogram of $1+W_i^\prime r_0$ and $1+W(\alpha_0(X_i),Y_i)^\prime r_0(X_i)$ over $1000$ replications for all numerical setups here. The values span a range not concentrated around $1$, indicating that the underlying probability distributions deviate substantially from the case of full independence, which highlights the generality of our setup. Table \ref{tab:unconditional} displays the results for $s=2,M=4,N=500$ and $s=5,M=4,N=500$ for the no-covariate cases. The detailed simulation results are included in the electronic companion. Table~\ref{tab:conditional} there reports the simulation results for $(s, M, N, x) = (2, 4, 100, 0.5)$ and $(5, 4, 100, 0.5)$ for the with-covariate cases. Each numerical result is based on 200 replications. The results indicate that the first-order estimators achieve lower RMSE than the plug-in estimator and provide more accurate probability estimates.
\begin{table}[htp!]
\centering
\begin{threeparttable}
\caption{Simulation results for the no-covariate case.}
\label{tab:unconditional}
\begin{tabular}{lcccccc}
\toprule
& \multicolumn{2}{c}{RMSE}
& \multicolumn{2}{c}{$\mathbb{E}\!\left[\max_{y\in\{0,1\}^M}\left|\hat{p}_y-p_y\right|\right]$}
& \multicolumn{2}{c}{$\mathbb{E}\!\left[\sum_{y\in\{0,1\}^M} p_y\left|\hat{p}_y-p_y\right|\right]$} \\
\cmidrule(lr){2-3}\cmidrule(lr){4-5}\cmidrule(lr){6-7}
Estimator & $\lambda_{\mathrm{cv}}$ & $\lambda_{\mathrm{theory}}$
& $\lambda_{\mathrm{cv}}$ & $\lambda_{\mathrm{theory}}$
& $\lambda_{\mathrm{cv}}$ & $\lambda_{\mathrm{theory}}$ \\
\midrule
\multicolumn{7}{l}{\textit{Panel A.} $s=2$, $M=4$, $N=500$} \\
PI ($w_{\mathrm{I}}$)              & 0.1203 & 0.3502 & 0.0138 & 0.0255 & 0.0089 & 0.0140 \\
PI ($w_{\mathrm{II}}$)             & 0.1304 & 0.3399 & 0.0171 & 0.0254 & 0.0061 & 0.0098 \\
FO ($w_{\mathrm{I}}$)              & 0.0913 & 0.1991 & 0.0136 & 0.0217 & 0.0070 & 0.0128 \\
FO ($w_{\mathrm{II}}$)             & 0.0902 & 0.2073 & 0.0141 & 0.0237 & 0.0057 & 0.0094 \\
PI (oracle $w_{\mathrm{I}}$)       & 0.0872 & 0.1079 & 0.0064 & 0.0205 & 0.0040 & 0.0045 \\
PI (oracle $w_{\mathrm{II}}$)      & 0.0858 & 0.1394 & 0.0063 & 0.0201 & 0.0037 & 0.0052 \\
\addlinespace
\multicolumn{7}{l}{\textit{Panel B.} $s=5$, $M=4$, $N=500$} \\
PI ($w_{\mathrm{I}}$)              & 0.0996 & 0.2809 & 0.0202 & 0.0320 & 0.0082 & 0.0143 \\
PI ($w_{\mathrm{II}}$)             & 0.0987 & 0.2840 & 0.0203 & 0.0292 & 0.0083 & 0.0125 \\
FO ($w_{\mathrm{I}}$)              & 0.0277 & 0.1551 & 0.0128 & 0.0271 & 0.0061 & 0.0113 \\
FO ($w_{\mathrm{II}}$)             & 0.0280 & 0.1548 & 0.0140 & 0.0247 & 0.0062 & 0.0116 \\
PI (oracle $w_{\mathrm{I}}$)       & 0.0182 & 0.0982 & 0.0102 & 0.0216 & 0.0054 & 0.0106 \\
PI (oracle $w_{\mathrm{II}}$)      & 0.0187 & 0.0995 & 0.0106 & 0.0212 & 0.0053 & 0.0102 \\
\bottomrule
\end{tabular}
\begin{tablenotes}
\footnotesize
\item[a] ``PI'' denotes the plug-in estimator and ``FO'' denotes the first-order estimator.
\item[b] $w_{\mathrm{I}}$ and $w_{\mathrm{II}}$ correspond to weighting schemes (I) and (II), respectively, for the no-covariate setting.
\item[c] All estimation errors are rounded to four decimal places.
\end{tablenotes}
\end{threeparttable}
\end{table}

\begin{table}[htp!]
\centering
\begin{threeparttable}
\caption{Simulation results for the with-covariate case.}
\label{tab:conditional}
\begin{tabular}{lcccccc}
\toprule
& \multicolumn{2}{c}{RMSE}
& \multicolumn{2}{c}{$\mathbb{E}\!\left[\max_{y\in\{0,1\}^M}\left|\hat{p}_y-p_y\right|\right]$}
& \multicolumn{2}{c}{$\mathbb{E}\!\left[\sum_{y\in\{0,1\}^M} p_y\left|\hat{p}_y-p_y\right|\right]$} \\
\cmidrule(lr){2-3}\cmidrule(lr){4-5}\cmidrule(lr){6-7}
Estimator & $\lambda_{\mathrm{cv}}$ & $\lambda_{\mathrm{theory}}$
& $\lambda_{\mathrm{cv}}$ & $\lambda_{\mathrm{theory}}$
& $\lambda_{\mathrm{cv}}$ & $\lambda_{\mathrm{theory}}$ \\
\midrule
\multicolumn{7}{l}{\textit{Panel A.} $s=2$, $M=4$, $N=100$, $x=0.5$, $h=0.1634$} \\
Plug-in ($w_{\mathrm{I}}$)         & 0.0915 & 0.9042 & 0.0498 & 0.0524 & 0.0194 & 0.0195 \\
Plug-in ($w_{\mathrm{II}}$)        & 0.0913 & 0.9058 & 0.0498 & 0.0524 & 0.0194 & 0.0195 \\
FO ($w_{\mathrm{I}}$)              & 0.0913 & 0.1681 & 0.0498 & 0.0498 & 0.0194 & 0.0194 \\
FO ($w_{\mathrm{II}}$)             & 0.0913 & 0.1682 & 0.0498 & 0.0498 & 0.0194 & 0.0194 \\
PI (oracle $w_{\mathrm{I}}$)       & 0.0912 & 0.1657 & 0.0070 & 0.0308 & 0.0056 & 0.0126 \\
PI (oracle $w_{\mathrm{II}}$)      & 0.0912 & 0.1657 & 0.0071 & 0.0308 & 0.0056 & 0.0126 \\
\addlinespace
\multicolumn{7}{l}{\textit{Panel B.} $s=5$, $M=4$, $N=100$, $x=0.5$, $h=0.1634$} \\
Plug-in ($w_{\mathrm{I}}$)         & 0.2552 & 0.2758 & 0.0894 & 0.0895 & 0.0327 & 0.0333 \\
Plug-in ($w_{\mathrm{II}}$)        & 0.2553 & 0.2749 & 0.0893 & 0.0895 & 0.0328 & 0.0332 \\
FO ($w_{\mathrm{I}}$)              & 0.2531 & 0.2539 & 0.0893 & 0.0892 & 0.0326 & 0.0332 \\
FO ($w_{\mathrm{II}}$)             & 0.2532 & 0.2538 & 0.0892 & 0.0890 & 0.0326 & 0.0332 \\
PI (oracle $w_{\mathrm{I}}$)       & 0.2432 & 0.2522 & 0.0327 & 0.0794 & 0.0134 & 0.0330 \\
PI (oracle $w_{\mathrm{II}}$)      & 0.2410 & 0.2429 & 0.0324 & 0.0773 & 0.0133 & 0.0332 \\
\bottomrule
\end{tabular}
\begin{tablenotes}
\footnotesize
\item[a] ``Plug-in'' and ``FO'' denote the plug-in and first-order estimators, respectively.
\item[b] $w_{\mathrm{I}}$ and $w_{\mathrm{II}}$ denote the choices of $w_k$ and $w_{k,1}$ under methods (I) and (II), respectively, for the with-covariate setting.
\item[c] All estimation errors are rounded to four decimal places.
\end{tablenotes}
\end{threeparttable}
\end{table}

\begin{figure}[htp!]
    \centering
    \captionsetup{font=small}
    \captionsetup[subfigure]{font=small}

    % Row 1
    \begin{subfigure}{0.3\linewidth}
        \centering
        \includegraphics[width=\linewidth]{new_simulation_results/uncond_s=2.png}
        \caption{$M=4, s=2$ (no covariate)}
    \end{subfigure}
    \hspace{0.02\linewidth}
    \begin{subfigure}{0.3\linewidth}
        \centering
        \includegraphics[width=\linewidth]{new_simulation_results/unconditional_s=5.png}
        \caption{$M=4, s=5$ (no covariate)}
    \end{subfigure}

    \vspace{0.2cm}

    % Row 2
    \begin{subfigure}{0.3\linewidth}
        \centering
        \includegraphics[width=\linewidth]{new_simulation_results/cond_s=2.png}
        \caption{$M=4, s=2$ (with covariate)}
    \end{subfigure}
    \hspace{0.02\linewidth}
    \begin{subfigure}{0.3\linewidth}
        \centering
        \includegraphics[width=\linewidth]{new_simulation_results/cond_s=5.png}
        \caption{$M=4, s=5$ (with covariate)}
    \end{subfigure}
    \caption{Histograms of $1+W_i'r_0$ (top row, no covariates) and $1+W(\alpha_0(X_i),Y_i)'r_0(X_i)$ (bottom row, with covariates), under sparsity levels $s \in \{2,5\}$ with $M=4$.}
    \label{fig:histogram}
\end{figure}

\subsection{Numerical Results for Estimating Average Treatment Effects}\label{appendix:ATE:numerical}
We numerically illustrate average treatment effect estimation with $M=4$, yielding $16$ (vector) treatment levels (binary tuples). Setting $(0,0,0,0)$ as the (vector) control level, we estimate average treatment effects for the remaining $15$ treatment levels versus $(0,0,0,0)$. The conditional outcome models are set as $O_i=\beta_t+\mu_t X_i+\epsilon_i$, where $\epsilon_i\indep X_i$, $\epsilon_i\overset{i.i.d.}{\sim}\mathcal{N}(0,1)$, $\beta_t=0.1H(t)$, $\mu_t=0.5+0.1H(t)$, where $H(t)=\sum_{j=1}^Mt_j2^{j-1}$. We generate $X_i$ uniformly from $\{0.1,0.2,0.3,0.4,0.5,0.6,0.7,0.8,0.9\}$.
Thus the average treatment effect of level $t$ is equal to $(\beta_t-\beta_0)+(\mu_t-\mu_0)\mathbb{E}[X_i]=0.15H(t)$, and the heterogeneous treatment effect of treatment level $t$ versus control $(0,0,0,0)$ is equal to $\tau(X_i)=(\beta_t-\beta_0)+(\mu_t-\mu_0)X_i$. The propensity scores are defined through Bahadur's representation: for any $t\in\{0,1\}^M$, $e_t(x)=(1+W(\alpha_0(x),t)'r_0(x))\prod_{j=1}^M\alpha_{0j}(x)^{t_j}(1-\alpha_{0j}(x))^{1-t_j}$, where $\alpha_{0}(x)$, $r_0(x)$ are computed in the same way as the with-covariate case simulation before. So we obtain the semiparametric efficiency variance for treatment level $t$ as $V^*(t)=\mathrm{Var}[\tau(X_i)]+\mathbb{E}\left[\frac{\sigma_0(X_i)^2}{e_{0}(X_i)}\right]+\mathbb{E}\left[\frac{\sigma_t(X_i)^2}{e_t(X_i)}\right]$, where $\mathrm{Var}[\tau(X_i)]=(\mu_t-\mu_0)^2\mathrm{Var}(X_i)=0.01/15H(t)^2$,
$\sigma_0(X_i)^2=\mathrm{Var}[O_i|X_i=x,T_i=(0,0,0,0)]=1$, and $\sigma_t(x)^2=\mathrm{Var}[O_i|X_i=x,T_i=t]=1$, thus $V^*(t)=0.01/15H(t)^2+\mathbb{E}[1/e_0(X_i)]+\mathbb{E}[1/e_t(X_i)]$ for any $t\in\{0,1\}^M\setminus\{(0,0,0,0)\}$. We compare the performance of the AIPW estimator using four generalized propensity score estimators. The first is the multinomial (MNL) estimator, which fits a multinomial logit model of $T_i$ on $X_i$ and is widely used in multi-level treatment settings. The remaining three estimators are computed as $\hat{e}_t(x)=(1+W(\hat{\alpha}(x),t)'\hat{r}(x,x))\prod_{j=1}^M\hat{\alpha}_j(x)^{t_j}(1-\hat{\alpha}_j(x))^{1-t_j}$, where we employ the Nadaraya-Watson (NW estimator), plug-in ($w_{\textrm{I}}$), and first-order ($w_{\textrm{I}}$) estimators for $\hat{r}(x,x)$ to estimate the generalized correlation coefficients $r_0(x)$. The kernel bandwidth is chosen as $h \propto (\log(p)/N)^{1/5}$, ranging from $0.3$ to $0.4$ depending on the sample size $N$. Specifically, we use $\displaystyle\hat{r}_{l}^{NW}(x):= \frac{\sum_{i=1}^NK_h(X_i-x)\prod_{j\in\ell(l)}z_{j}(T_{i},\hat{\alpha}(X_i))}{\sum_{i=1}^NK_h(X_i-x)}$ as the $l$-th component of the NW estimator $\hat{r}^{NW}(x)$ to approximate \eqref{def:r0}. We checked the probability estimates before plugging them into the AIPW estimator and found that they are all positive, so this is not an issue in our simulations. In practice, if they fall outside the feasible range, one can truncate them before plugging them into the AIPW estimator. We replicate the simulations 200 times and calculate the coverage ratios of the semiparametric efficient confidence intervals for the true ATEs at each treatment level. Simulations are conducted for $s=0,2,5,10$ with total sample size $N=200,300,400$ in each case. The effective sample size for each treatment combination is equal to $N/16$ on average. Figure \ref{fig:coverage} presents the coverage ratios of true ATE for all estimators across combinations of $s$ and $N$, where the vertical axis corresponds to treatment levels and the horizontal axis to estimators. 

Several key observations can be made from the numerical results:
First, the coverage ratios of the MNL model and the NW estimator remain stable across different values of $s$ and treatment levels. However, their coverage ratios increase only slightly as $N$ grows. Second, for the plug-in and first-order estimators, the coverage ratio generally increases as $s$ decreases and $N$ grows. In contrast, such pattern is not observed for the MNL model and the NW estimator, as both do not account for the underlying independence structure of the binary treatments. Third, the NW estimator performs significantly worse than both the plug-in and first-order estimators. 
\begin{figure}
    \centering
    \includegraphics[width=\linewidth]{coverage_plot.png}
    \caption{Coverage ratios of true ATEs under different propensity score estimators, across sparsity levels $s$ and sample sizes $N$, with $M=4$ (16 treatment combinations; $(0,0,0,0)$ as control).}
    \label{fig:coverage}
\end{figure}

\section{Discussion}\label{sec:discussion}
We develop a method for estimating assortment probabilities. Existing approaches often ignore independence structures in the joint distribution and can therefore be statistically inefficient in high dimensions. Our approach is adaptive to the independence structure by leveraging sparsity in the generalized correlation coefficients \citep{bahadur1959representation}, yielding improved finite-sample performance both theoretically and numerically. In our framework, the marginal distributions are treated as nuisance parameters, while the high-dimensional generalized correlation coefficients are the target of interest. 
Our adversarial approach also advances the study of high-dimensional estimand with low-dimensional nuisance parameters. This perspective contrasts with the de-biased machine learning literature \citep[e.g.,][]{chernozhukov2017double}, which typically studies low-dimensional targets with high-dimensional nuisance parameters instead.

We close the paper by a discussion on possible extensions. First, for the covariate-case, we establish a pointwise rate of convergence result for $r_0(x)$. So a natural extension is to establish uniform convergence results over all $x\in\mathcal{X}$. Second, it would also be interesting to incorporate prior information about independence directly into the generalized correlation coefficients and to develop localized sparsity constraints, such as sparsity levels $s_x$ that vary across covariate strata. We expect our method to remain valid provided that the covariate dimension is low. It is also interesting to incorporate high-dimensional covariates, which would require new modeling strategies and estimators. Further, our approach offers a fundamentally different way to model independence from the graphical conditional independence. \citet{yuan2021community} shows in simulations that their community detection method based on a truncated Bahadur representation outperforms approaches that assume conditional independence. This suggests that our method may also be useful in similar settings, such as community detection, to capture dependence among edges. Computationally, the adversarial set of the first-order estimator has at least $\mathrm{O}(p)$ extreme points with $p\asymp2^M$, manageable via parallelization for moderate $M$. For very large $M$, an approximate dual reformulation of the inner minimization may improve tractability, which we leave for future work. Lastly, recall from Remark~\ref{remark:non-negativity} that since $(\hat \alpha, r_0)$ may not lie in $\mathbb{K}$, we do not impose this constraint in the optimization. For the plug-in estimator, enforcing $(\hat{\alpha},r)\in\mathbb{K}$ requires adding $2^M$ linear constraints. For the first-order estimator, this requirement would be even more restrictive, since it would require $(\alpha,r)\in\mathbb{K}$ for every $\alpha\in\mathcal{A}(\hat{\alpha})$, which can be computationally prohibitive. Moreover, this introduces a tradeoff: the feasible set may become overly restrictive, and since it is possible that $(\hat \alpha, r_0)\notin\mathbb{K}$, imposing such constraints can lead to large estimation bias. The discussion in Remark~\ref{remark:non-negativity} indicates that our result is sufficient for a broad range of practical applications. For the causal inference application, $\min_{i\in[N]}1+W(\hat{\alpha}(X_i),t)'\hat{r}^{FO}(X_i,X_i)>0$ with high probability, so the AIPW estimator is generally not impacted either. In the rare cases where positivity of the probability estimates fails, a simple truncation step can be applied. There could exist alternative projection-based approaches that could enforce this constraint, but analyzing such methods would require a different type of analysis. As discussed above, the current approach is already sufficiently general for many practically relevant settings, so we leave such extensions to future work.


\bibliographystyle{abbrvnat}
\bibliography{bibliography}

\def\spacingset#1{\renewcommand{\baselinestretch}%
{#1}\small\normalsize} \spacingset{1}
\onehalfspacing
\begin{center}
\textbf{\Large{}{}Appendix for\\ ``Adversarial Estimation of Assortment Probabilities under Independence Structure''}
{\Large{} }{\Large\par}
\par\end{center}

\begin{center}
Alexandre Belloni\footnote{Fuqua School of Business, Duke University. Email: abn5@duke.edu}, Yan Chen\footnote{Fuqua School of Business, Duke University. Email: yc555@duke.edu} and Matthew Harding\footnote{Department of Economics,
University of California, Irvine. Email: matthew.harding@uci.edu}
{\Large{} }{\Large\par}
\par\end{center}

\onehalfspacing
\section{Proof of Theorem \ref{Thm:Bahadur}}
The proof follows similarly to Proposition 1 in  \cite{bahadur1959representation}. We include a proof for completeness.  Fix $X=x$. Since $p(y\mid x)=\mathbb{P}(Y=y\mid X=x)$ has a domain with $2^M$ points and $\sum_{y\in\mathcal{Y}}p(y\mid x) = 1$, there are $2^M-1$ parameters to characterize $p(\cdot\mid x)$. Let $p_{[I\mid x]}(y\mid x) = \prod_{j=1}^M \alpha_j(x)^{y_j}(1-\alpha_j(x))^{1-y_j}$ and note that $p_{[I]}(y,x)>0$ since $0<\alpha_j(x)<1$ for all $j\in[M]$. Consider the vector space $V:=\{ h: \mathcal{Y} \to \mathbb{R} \} $, with inner-product $\langle h,g\rangle = \mathbb{E}_{p_{[I\mid x]}}[f\cdot g] = \sum_{y \in \mathcal{Y}} h(y)g(y)p_{[I]}(y\mid x)$.
Under $p_{[I\mid x]}$ it follows that the set of functions $S=\left\{1; \prod_{k=1}^{|\ell|} z_{\ell_k}(\cdot,\alpha(x)),  \ell \subseteq [M] \right\}$
is an orthonormal basis for $V$. Indeed under $p_{[I\mid x]}$ the components are independence so that unit norm holds $$\mathbb{E}_{p_{[I\mid x]}}\left[ \left(\prod_{k=1}^{|\ell|} z_{\ell_k}(\cdot ,\alpha(x))\right)^2\right]=\mathbb{E}_{p_{[I\mid x]}}\left[ \prod_{k=1}^{|\ell|} z_{\ell_k}^2(\cdot ,\alpha(x))\right]=\prod_{k=1}^{|\ell|} \mathbb{E}_{p_{[I\mid x]}}[z_{\ell_k}^2(\cdot ,\alpha(x))]=1$$ 
and the orthogonality follows from noting that for any different $\ell,\tilde\ell \subseteq M$, letting $\tilde \ell' = \ell \cap  \tilde \ell$ and $\ell' = \ell \cup  \tilde \ell \setminus \ell \cap  \tilde \ell$ we have
{$$\begin{array}{rl} 
&\displaystyle\quad \mathbb{E}_{p_{[I\mid x]}}\left[\prod_{k=1}^{|\ell|} z_{\ell_k}(\cdot ,\alpha(x)) \prod_{k=1}^{|\tilde \ell|} z_{\tilde \ell_k}(\cdot ,\alpha(x))\right] \\ 
&\displaystyle=\mathbb{E}_{p_{[I\mid x]}}\left[\left(\prod_{k=1}^{|\ell'|} z_{\ell_k}(\cdot,\alpha(x))\right) \left(\prod_{k=1}^{|\tilde \ell'|} z_{\tilde \ell_k'}^2(\cdot,\alpha(x))\right)\right]\\
&\displaystyle = \left(\prod_{k=1}^{|\ell'|} \mathbb{E}_{p_{[I\mid x]}}\left[ z_{\ell_k}( \cdot,\alpha(x))\right]\right) \prod_{k=1}^{|\tilde \ell'|}\mathbb{E}_{p_{[I\mid x]}}[  z_{\tilde \ell_k'}^2( \cdot,\alpha(x))]=0
\end{array}$$}
Therefore, since $\frac{p(\cdot\mid x)}{p_{[I\mid x]}(\cdot\mid x)} \in V$, it has a linear representation with the basis $S$, namely $\frac{p(y\mid x)}{p_{[I\mid x]}(y\mid x)} = \sum_{h \in S} \langle \frac{p(\cdot\mid x)}{p_{[I\mid x]}(\cdot\mid x)}, h(\cdot)\rangle h(y)$, where for $h(y)= \prod_{k=1}^{|\ell|}z_{\ell_k}(y,\alpha(x))$ we have
$$\begin{array}{rl} 
\displaystyle\left\langle \frac{p(\cdot\mid x)}{p_{[I\mid x]}(\cdot\mid x)}, h(\cdot)\right\rangle &\displaystyle = \sum_{y\in\mathcal{Y}} \frac{p(y\mid x)}{p_{[I\mid x]}(y\mid x)}\prod_{k=1}^{|\ell|}z_{\ell_k}(y,\alpha(x)) \cdot p_{[I\mid x]}(y\mid x) \\
&\displaystyle = \sum_{y\in\mathcal{Y}} \prod_{k=1}^{|\ell|}z_{\ell_k}(y,\alpha(x)) \cdot p(y\mid x) = r_{0\ell}(x)\\
\end{array}$$
Thus the specific representation of 
$$f(y,\alpha(x),r_0) = \frac{p(\cdot\mid x)}{p_{[I\mid x]}(\cdot\mid x)} = 1 + \sum_{k=2}^M \sum_{\ell \subset [M], |\ell|=k} r_{0\ell}(x) \prod_{m=1}^k z_{\ell_m}(y,\alpha(x))$$  
follows from $\mathbb{E}_{p_{[I\mid x]}}\left[\frac{p(\cdot\mid x)}{p_{[I\mid x]}(\cdot\mid x)}\right]=\sum_{y\in\mathcal{Y}}p(y\mid x) = 1$ and $\mathbb{E}_{p_{[I\mid x]}}\left[\frac{p(\cdot \mid x)}{p_{[I\mid x]}(\cdot\mid x)}z_j(\cdot,\alpha(x))\right]=\mathbb{E}[z_j(Y,\alpha(x))]=0$.


\section{Proofs and Additional Results of Section \ref{Sec:Rates:no-covariates}}\label{sec:proof:no-cov}
\begin{proof}[Proof of Proposition~\ref{prop:plug-in:no-covariate}]
By Lemma \ref{lemma:Score:nocovariate}, with probability $1-\epsilon-\delta_{\alpha}-\delta(1+\mathrm{o}(1))$, we have 
$$\left\|\frac{{\rm diag}^{-1}(w)}{N}\!\!\sum_{i=1}^N\!\!\frac{W(\hat\alpha,Y_i)}{1+W(\hat\alpha,Y_i)'r_0}\right\|_\infty\!\!\!\!\!\leq\left\|\frac{{\rm diag}^{-1}(w)}{N}\!\!\sum_{i=1}^N\!\!\frac{W(\alpha_0,Y_i)}{1+W(\alpha_0,Y_i)'r_0}\right\|_\infty \!\!\!\!+C(1+\beta_N)\|\hat\alpha-\alpha_0\| \frac{M^{1/2}K_{\alpha,1}}{\xi}.$$ 
Let $\mathcal{W}:={\rm diag}(w)$. By H\"older's inequality, we have 
\begin{equation}\label{ineq:no-cov:holder}
\displaystyle v'\hat Uv \geq v'Uv - \|v\|_{w,1}^2\|\mathcal{W}^{-1}(U-\hat U)\mathcal{W}^{-1}\|_{\infty,\infty}.    
\end{equation}
We apply \eqref{ineq:no-cov:holder} for 
$U:=\frac{1}{N}\sum_{i=1}^N\frac{W(\alpha_0,Y_i)W(\alpha_0,Y_i)'}{(1+W(\alpha_0,Y_i)^\prime r_0)^2},\ \hat U :=\frac{1}{N}\sum_{i=1}^N\frac{W(\hat\alpha,Y_i)W(\hat\alpha,Y_i)'}{(1+W(\hat\alpha,Y_i)^\prime r_0)^2}$. Note that with probability $1-\delta_{\alpha}$, $1+W(\hat{\alpha},Y_i)'r_0\geq_{(1)}1+W(\alpha_0,Y_i)'r_0-|(\hat{\alpha}-\alpha_0)'\nabla_{\alpha}W(\tilde{\alpha},Y_i)'r_0|\geq_{(2)}\xi-M^{1/2}K_{\alpha,1}\|\hat{\alpha}-\alpha_0\|\geq_{(3)}\frac{3}{4}\xi>0$, where (1) follows from Taylor expansion and $\tilde{\alpha}$ is on the line segment between $\alpha_0, \hat{\alpha}$. So when $\alpha_0\in\mathcal{A}(\hat{\alpha})$, by the construction of $\mathcal{A}(\hat{\alpha})$, $\tilde\alpha\in\mathcal{A}(\hat{\alpha})$. Hence following from Assumption~\ref{ass:iid-min-den} and  Assumption~\ref{ass:smoothness-assumptions}, $\min_{i\in[N]}1+W(\hat{\alpha},Y_i)'r_0\geq\frac{3}{4}\xi>0$ with probability $1-\epsilon-3\delta_{\alpha}$. Hence with probability $1-C\epsilon-C\delta_{\alpha}-C\delta(1+\mathrm{o}(1))$, we have  
$$\begin{array}{rl}
&\displaystyle\quad\frac{c_0}{N}\sum_{i=1}^N\left\{\frac{W(\alpha_0,Y_i)^\prime(\hat{r}^{PI}-r_0)}{1+W(\alpha_0,Y_i)^\prime r_0}\right\}^2\\ 
&\displaystyle\leq_{(1)}\frac{c_0}{N}\sum_{i=1}^N\left\{\frac{W(\hat\alpha,Y_i)^\prime(\hat{r}^{PI}-r_0)}{1+W(\hat\alpha,Y_i)^\prime r_0}\right\}^2+c_0\|\hat r^{PI} - r_0\|_{w,1}^2 \|\mathcal{W}^{-1}(\hat U - U)\mathcal{W}^{-1} \|_{\infty,\infty}\\
&\displaystyle \leq_{(2)}2\lambda\frac{c+1}{c-1}\norm{\hat{r}_T^{PI}-r_0}_{w,1} + c_0\left(\frac{2c}{c-1}\right)^2\|\hat r_T^{PI} - r_0\|_{w,1}^2 \|\mathcal{W}^{-1}(\hat U - U)\mathcal{W}^{-1} \|_{\infty,\infty}\\
&\displaystyle \leq_{(3)} \frac{2\lambda\sqrt{s}(c+1)/(c-1)}{\kappa_{\bar c}\sqrt{c_0}}\sqrt{\frac{c_0}{N} \sum_{i=1}^N\left\{\frac{W(\alpha_0,Y_i)^\prime(\hat{r}^{PI}-r_0)}{1+W(\alpha_0,Y_i)^\prime r_0}\right\}^2}  \\
&\displaystyle\quad\quad + \frac{4c^2s\|\mathcal{W}^{-1}(\hat U - U)\mathcal{W}^{-1} \|_{\infty,\infty}}{(c-1)^2\kappa_{\bar c}^2}\frac{c_0}{N} \sum_{i=1}^N\left\{\frac{W(\alpha_0,Y_i)^\prime(\hat{r}^{PI}-r_0)}{1+W(\alpha_0,Y_i)^\prime r_0}\right\}^2,
\end{array}$$
where (1) follows by (\ref{ineq:no-cov:holder}). The first term of (2) follows since by Lemma \ref{lemma:fast-rate:estimated-alpha}, for some universal constant $c_0>0$, with probability $1-\epsilon-\delta_{\alpha}-\delta(1+\mathrm{o}(1))$, we have $\frac{c_0}{N}\sum_{i=1}^N\!\left\{\frac{W(\hat\alpha,Y_i)^\prime(\hat{r}^{PI}-r_0)}{(1+W(\hat\alpha,Y_i)^\prime r_0)}\right\}^2\!\!\leq2\lambda\frac{c+1}{c-1}\norm{\hat{r}_T^{PI}-r_0}_{w,1}$. The second term of (2) follows since $\|\hat{r}^{PI}-r_0\|_{w,1}=\|\hat{r}_T^{PI}-r_0\|_{w,1}+\|\hat{r}_{T^c}^{PI}-r_0\|_{w,1}\leq\frac{2c}{c-1}\|\hat{r}_T^{PI}-r_0\|_{w,1}$ according to Lemma~\ref{lem:diff-restricted}. (3) follows from the definition of restricted eigenvalues and $\|\hat{r}_T^{PI}-r_0\|_{w,1}\leq\sqrt{s}\|\hat{r}_T^{PI}-r_0\|_{w,2}$. By Taylor expansion, for each $j\in[p]$, $i\in[N]$ there exists $\tilde \alpha^{ji}$ on the line segment between $\hat \alpha$ and $\alpha_0$ such that
$W_j(\hat\alpha,Y_i) = W_j(\alpha_0,Y_i) + (\hat\alpha-\alpha_0)'\nabla_\alpha W_j(\tilde \alpha^{ji},Y_i)$. Simiarly, for some $\tilde\alpha$ on the line segment between $\hat{\alpha}$ and $\alpha_0$, for any $j\in[p]$, 
$$\begin{array}{rl}
\displaystyle\frac{W_j(\hat\alpha,Y_i)}{1+W(\hat\alpha,Y_i)'r_0}&\displaystyle= \frac{W_j(\alpha_0,Y_i) +   (\hat\alpha-\alpha_0)'\nabla_\alpha W_j(\tilde \alpha^{ji},Y_i)}{1+W(\alpha_0,Y_i)'r_0 + (\hat\alpha-\alpha_0)'\nabla_\alpha W(\tilde \alpha,Y_i)'r_0} \\
&\displaystyle = \frac{W_j(\alpha_0,Y_i) }{1+W(\alpha_0,Y_i)'r_0} \left( 1 -  \frac{(\hat\alpha-\alpha_0)'\nabla_\alpha W(\tilde \alpha,Y_i)'r_0}{1+W(\alpha_0,Y_i)'r_0 + (\hat\alpha-\alpha_0)'\nabla_\alpha W(\tilde \alpha,Y_i)'r_0}\right)\\
&\displaystyle\quad + \frac{(\hat\alpha-\alpha_0)'\nabla_\alpha W_j(\tilde \alpha^{ji},Y_i)}{1+W(\alpha_0,Y_i)'r_0 + (\hat\alpha-\alpha_0)'\nabla_\alpha W(\tilde \alpha,Y_i)'r_0} = \frac{W_j(\alpha_0,Y_i) }{1+W(\alpha_0,Y_i)'r_0} + \widetilde R_j(\hat\alpha,Y_i).
\end{array}$$ 
By the construction of $\mathcal{A}(\hat{\alpha})$ as shown later in Section~\ref{sec:marginals}, $\hat{\alpha}\in\mathcal{A}(\hat{\alpha})$. Particularly, 
$$|(\hat\alpha-\alpha_0)'\nabla_\alpha W(\tilde \alpha,Y_i)'r_0|\leq_{(a)} K_{\alpha,1}\sqrt{M}\|\hat\alpha-\alpha_0\|\leq_{(b)}K_{\alpha,1}c_1\bigg(\sqrt{\frac{M}{N}}+\sqrt{\frac{\log(1/\delta)}{N}}\bigg)<_{(c)}\frac{1}{2}\xi,$$ 
where (a) holds with probability $1-\delta_{\alpha}$ from Assumption~\ref{ass:smoothness-assumptions}, (b) holds with probability $1-\delta$ according to Proposition~\ref{prop:concentration:no covariate}, (c) holds when $N$ is sufficiently large. Since $1+W(\alpha_0,Y_i)'r_0>\xi$ with probability $1-\epsilon$, $|\widetilde R_j(\hat\alpha,Y_i)|\leq \frac{CK_{\alpha,1}}{\xi}\left(1+\frac{|W_j(\alpha_0,Y_i)|}{1+W(\alpha_0,Y_i)'r_0}\right)M^{1/2}\|\hat\alpha-\alpha_0\|$ with probability $1-\delta_\alpha-\delta-\epsilon$ when $N$ is sufficiently large, where $C$ is an absolute constant.
Define $W_i := \frac{W(\alpha_0,Y_i)}{1+W(\alpha_0,Y_i)'r_0}$, $\hat W_i := \frac{W(\hat \alpha,Y_i)}{1+W(\hat \alpha,Y_i)'r_0}$. For any $l\in[p]$, define $\tilde R_{il}:=\tilde{R}_l(\hat{\alpha},Y_i)$. By Assumption \ref{ass:smoothness-assumptions}, with probability $1-\delta_\alpha$, $\left(1+\frac{|W_l(\alpha_0,Y_i)| }{1+W(\alpha_0,Y_i)'r_0}\right)\|\hat\alpha - \alpha_0\|\frac{M^{1/2}K_{\alpha,1}}{\xi} \leq\frac{1}{4}\left(1+\frac{|W_l(\alpha_0,Y_i)|}{1+W(\alpha_0,Y_i)'r_0}\right)$, and $\frac{1}{N}\sum_{i=1}^N |\tilde R_{il}\tilde R_{ik}|\leq\frac{C}{4N}\sum_{i=1}^N (1+|W_{il}|)|\tilde R_{ik}|$. Note $|\hat W_{il}\hat W_{ik} -W_{il}W_{ik}| \leq \left| W_{il}\tilde R_{ik} + \tilde R_{il} W_{ik} + \tilde R_{il}\tilde R_{ik}\right|$. Thus with probability $1-C'\delta_{\alpha}-C'\delta-C'\epsilon$, $\displaystyle\left| U_{lk} - \hat U_{lk} \right| \leq C'\frac{M^{1/2}K_{\alpha,1}}{\xi}\|\hat\alpha - \alpha_0\|\max_{l,k\in[p]}\frac{1}{N}\sum_{i=1}^N  \left(1+| W_{il}|\right)\left(1+| W_{ik}|\right)$ for some constant $C'$ independent of $s,N,M$, so the given conditions further imply that for some constant $\bar{C}$ independent of $s,N,M$, we have $s\|\mathcal{W}^{-1}(\hat U - U)\mathcal{W}^{-1} \|_{\infty,\infty}/\kappa_{\bar{c}}^2\leq \bar{C}sM^{1/2}K_{\alpha,1}\|\hat\alpha-\alpha_0\|/\{\xi\kappa_{\bar c}^2 \min_{j\in[p]} w_j^2\}\leq\delta_N\bar{C}$.
Recall that $\delta_N\rightarrow0$. So for sufficiently large $N$, $4c^2s\|\mathcal{W}^{-1}(\hat U - U)\mathcal{W}^{-1} \|_{\infty,\infty}/\{\kappa_{\bar{c}}^2(c-1)^2\}\leq1/2$. Let $x=\sqrt{\frac{c_0}{N}\sum_{i=1}^N\left\{\frac{W(\alpha_0,Y_i)^\prime(\hat{r}^{PI}-r_0)}{1+W(\alpha_0,Y_i)^\prime r_0}\right\}^2}$. Note that $\frac{x^2}{2} \leq \frac{2\lambda\sqrt{s}}{\sqrt{c_0}\kappa_{\bar{c}}}\frac{c+1}{c-1} x \Rightarrow x \leq \frac{4\lambda\sqrt{s}}{\sqrt{c_0}\kappa_{\bar{c}}}\frac{c+1}{c-1}$. This implies 
$\sqrt{\frac{c_0}{N} \sum_{i=1}^N\left\{\frac{W(\alpha_0,Y_i)^\prime(\hat{r}^{PI}-r_0)}{1+W(\alpha_0,Y_i)^\prime r_0}\right\}^2}\leq\frac{4\lambda\sqrt{s}}{\sqrt{c_0}\kappa_{\bar{c}}}\frac{c+1}{c-1}$. By Proposition \ref{prop:concentration:no covariate}, with probability $1-\delta$, $\|\hat{\alpha}-\alpha_0\|\leq c_1\big(\sqrt{\frac{M}{N}}+\sqrt{\frac{\log(1/\delta)}{N}}\big)$ for some $c_1>0$. Further note $\delta>1/2^M$. So the result follows. 
\end{proof}
\medskip
\begin{proof}[Proof of Theorem~\ref{thm:main:FO}]
Firstly, define 
$Q^{FO}(\alpha,r):= \frac{1}{N}\sum_{i=1}^N \log(1+ W^{FO}(\alpha,Y_i)'r)$, where $W^{FO}(\alpha,y):=W(\hat \alpha,y)+\nabla_\alpha W(\hat\alpha,y)(\alpha-\hat\alpha)$. $\mathcal{E}_0:= \{ \alpha_0 \in \mathcal{A}(\hat\alpha)\}$, $\mathcal{E}_1:= \{ \min_{\alpha \in \mathcal{A}(\hat\alpha)} Q^{FO}(\alpha,r_0)-Q^{FO}(\alpha_0,r_0) \geq - \widetilde{\mathcal{R}}(\hat\alpha) \}$, where $\widetilde{\mathcal{R}}(\hat\alpha)= C\bar{K}_2\sqrt{\frac{\log (M/\delta)}{N}}\max_{\alpha \in \mathcal{A}(\hat\alpha)}\|\alpha-\alpha_0\|_1+\frac{CK_{\alpha,2}}{\xi^2} \max_{\alpha \in \mathcal{A}(\hat\alpha)}\|\alpha - \alpha_0\|^2$. $\mathcal{E}_2:=\{\frac{\lambda}{c} \geq \left\|{\rm diag}^{-1}(w) \nabla_r Q^{FO}(\alpha_0,r_0)\right\|_\infty \}$, $\mathcal{E}_3:=\{\min_{i\in[N]} 1+W^{FO}(\alpha_0,Y_i)'r_0\geq\xi/2\}$. 
Assumption \ref{ass:smoothness-assumptions} implies that $\mathcal{E}_0$ holds with probability at least $1-\delta_\alpha$. $\mathcal{E}_1$ holds with probability $1-\epsilon-\delta_{\alpha}-\delta(1+\mathrm{o}(1))$ by Lemma \ref{lemma:MinMax:Claim1}. $\mathcal{E}_2$  holds with probability $1-\delta-\delta_{\alpha}$ by the conditions of Theorem~\ref{thm:main:FO} and Lemma \ref{lemma:Score:nocovariate}. Note that with probability $1-\epsilon-C\delta_{\alpha}$, 
$\min_{i\in[N]}1+W^{FO}(\alpha_0,Y_i)'r_0\geq\min_{i\in[N]}1+W(\hat{\alpha},Y_i)'r_0+(\alpha_0-\hat{\alpha})'\nabla_{\alpha}W(\hat{\alpha},Y_i)'r_0\geq_{(1)}\frac{3}{4}\xi-M^{1/2}K_{\alpha,1}\|\alpha_0-\hat{\alpha}\|\geq_{(2)}\frac{1}{2}\xi$, where inequality (1) follows from Holder's inequality and the fact that the proof of Proposition~\ref{prop:plug-in:no-covariate} shows that with probability $1-\epsilon-C'\delta_{\alpha}$, $1+W(\hat{\alpha},Y_i)'r_0\geq\frac{3}{4}\xi$, and inequality (2) follows from  Assumption~\ref{ass:smoothness-assumptions}. Let $\hat{r}=\hat{r}^{FO}$. It is then sufficient to prove that the result holds on $\mathcal{E}_0\cap\mathcal{E}_1\cap\mathcal{E}_2\cap\mathcal{E}_3$. Particularly on the intersection of all these events we have 
\begin{equation}\label{eq:main:steps}
\begin{array}{rl}
&\quad Q^{FO}(\alpha_0,\hat r) - \lambda \|\hat r\|_{w,1}\\
& \geq_{(1)} \displaystyle \min_{\alpha \in \mathcal{A}(\hat\alpha)} Q^{FO}(\alpha,\hat r) - \lambda \|\hat r\|_{w,1} \\
& \displaystyle  \geq_{(2)} \min_{\alpha \in \mathcal{A}(\hat\alpha)} Q^{FO}(\alpha,r_0) - \lambda \|r_0\|_{w,1}\\
&\displaystyle  = \min_{\alpha \in \mathcal{A}(\hat\alpha)} \{Q^{FO}(\alpha,r_0)-Q^{FO}(\alpha_0,r_0)\}\\
&\quad\displaystyle + Q^{FO}(\alpha_0,r_0)- \lambda \|r_0\|_{w,1} \\
&\displaystyle  \geq_{(3)} - \widetilde{\mathcal{R}}(\hat\alpha) + Q^{FO}(\alpha_0,r_0)- \lambda \|r_0\|_{w,1},
\end{array}
\end{equation}
where (1) holds by $\mathcal{E}_0$, (2) holds by the optimality of $\hat r$ and (3) holds by $\mathcal{E}_1$. Following similar proof steps as in Lemma~\ref{lemma:fast-rate:estimated-alpha}, $Q^{FO}(\alpha_0,\hat r)\leq Q^{FO}(\alpha_0, r_0) + \nabla_r Q^{FO}(\alpha_0, r_0)'(\hat r - r_0)  - \frac{c_0}{N} \sum_{i=1}^N\left\{\frac{W^{FO}(\alpha_0,Y_i)^\prime(\hat{r}-r_0)}{1+W^{FO}(\alpha_0,Y_i)^\prime r_0}\right\}^2$ for some universal constant $c_0>0$. Combining this and (\ref{eq:main:steps}), on event $\mathcal{E}_0\cap\mathcal{E}_1\cap\mathcal{E}_2\cap\mathcal{E}_3$, using (weighted) H\"older's inequality we have 
\begin{equation}\label{eq:thm1:main:main}\begin{array}{rl}
&\displaystyle\quad\frac{c_0}{N} \sum_{i=1}^N\left\{\frac{W^{FO}(\alpha_0,Y_i)^\prime(\hat{r}-r_0)}{1+W^{FO}(\alpha_0,Y_i)^\prime r_0}\right\}^2\\
&\displaystyle \leq   \lambda\|r_0\|_{w,1} -  \lambda\|\hat r\|_{w,1}   + \nabla_r Q^{FO}(\alpha_0,r_0)'(\hat r - r_0)\\
&\quad\displaystyle + \widetilde{\mathcal{R}}(\hat\alpha)\\ 
&\displaystyle \leq \lambda\|r_0\|_{w,1} -  \lambda\|\hat r\|_{w,1}+ \widetilde{\mathcal{R}}(\hat\alpha)\\
&\displaystyle\quad + \left\|{\rm diag}^{-1}(w) \nabla_r Q^{FO}(\alpha_0,r_0)\right\|_\infty\|\hat r - r_0\|_{w,1} \\
&\displaystyle \leq    \lambda\|r_0\|_{w,1} -  \lambda\|\hat r\|_{w,1} + \frac{\lambda}{c}\|\hat r - r_0\|_{w,1}   + \widetilde{\mathcal{R}}(\hat\alpha)\\
&\displaystyle \leq \left(1+\frac{1}{c}\right)\lambda \|\hat r_T-r_0\|_{w,1} -  \lambda \left(1-\frac{1}{c}\right)\|\hat r_{T^c}\|_{w,1}\\
&\displaystyle\quad + \widetilde{\mathcal{R}}(\hat\alpha).
\end{array}
\end{equation}
Let $\mathcal{W}:={\rm diag}(w)$. H\"older's inequality implies $v'\hat Uv \geq v'Uv - \|v\|_{w,1}^2\|\mathcal{W}^{-1}(U-\hat U)\mathcal{W}^{-1}\|_{\infty,\infty}$.
Define $\displaystyle U:=\frac{c_0}{N}\sum_{i=1}^N\frac{W(\alpha_0,Y_i)W(\alpha_0,Y_i)'}{(1+W(\alpha_0,Y_i)^\prime r_0)^2}$, $\displaystyle\hat U :=\frac{c_0}{N}\sum_{i=1}^N\frac{W^{FO}(\alpha_0,Y_i)W^{FO}(\alpha_0,Y_i)'}{(1+W^{FO}(\alpha_0,Y_i)^\prime r_0)^2}$, we have 
\begin{equation}\label{eq:DesignMatrixConnection}
\begin{array}{rl}
&\quad\displaystyle\frac{c_0}{N} \sum_{i=1}^N\left\{\frac{W^{FO}(\alpha_0,Y_i)^\prime(\hat{r}-r_0)}{1+W^{FO}(\alpha_0,Y_i)^\prime r_0}\right\}^2\\
&\quad-\displaystyle\frac{c_0}{N} \sum_{i=1}^N\left\{\frac{W(\alpha_0,Y_i)^\prime(\hat{r}-r_0)}{1+W(\alpha_0,Y_i)^\prime r_0}\right\}^2 \\
& \geq \displaystyle -\|\hat r - r_0\|_{w,1}^2 \|\mathcal{W}^{-1}(U-\hat U)\mathcal{W}^{-1}\|_{\infty,\infty} \\
\end{array}\end{equation}
Recall that $T$ is the support of $r_0$. For the rest of the proof, we show that the result holds under two exclusive cases: $\|\hat r_{T^c}\|_{w,1} \geq 2\frac{c+1}{c-1}\|\hat r_T - r_0\|_{w,1}$ and $\|\hat r_{T^c}\|_{w,1}<2\frac{c+1}{c-1}\|\hat r_T - r_0\|_{w,1}$. First assume that $\|\hat r_{T^c}\|_{w,1} \geq 2\frac{c+1}{c-1}\|\hat r_T - r_0\|_{w,1}$.
From (\ref{eq:thm1:main:main}) we have 
\begin{equation}\label{eq:nocovariate:UnRestrictedSet} 
\begin{array}{rl}
&\displaystyle\quad\frac{c_0}{N} \sum_{i=1}^N\left\{\frac{W^{FO}(\alpha_0,Y_i)^\prime(\hat{r}-r_0)}{1+W^{FO}(\alpha_0,Y_i)^\prime r_0}\right\}^2\\
&\displaystyle\leq \widetilde{\mathcal{R}}(\hat\alpha)-\frac{\lambda}{2}\left(1-\frac{1}{c}\right)\|\hat r_{T^c}\|_{w,1}.
\end{array}
\end{equation}
Thus from \eqref{eq:DesignMatrixConnection} and \eqref{eq:nocovariate:UnRestrictedSet} we have 
$$\begin{array}{rl}
&\displaystyle\quad\frac{c_0}{N} \sum_{i=1}^N\left\{\frac{W^{FO}(\alpha_0,Y_i)^\prime(\hat{r}-r_0)}{1+W^{FO}(\alpha_0,Y_i)^\prime r_0}\right\}^2  -   \frac{c_0}{N}\sum_{i=1}^N\left\{\frac{W(\alpha_0,Y_i)^\prime(\hat{r}-r_0)}{1+W(\alpha_0,Y_i)^\prime r_0}\right\}^2 \\
&\displaystyle \geq  -\|\hat r - r_0\|_{w,1}^2 \|\mathcal{W}^{-1}(U-\hat U)\mathcal{W}^{-1}\|_{\infty,\infty} \\
&\displaystyle=   -(\|\hat r_T - r_0\|_{w,1} + \| \hat r_{T^c}\|_{w,1})^2 \|\mathcal{W}^{-1}(U-\hat U)\mathcal{W}^{-1}\|_{\infty,\infty}\\
&\displaystyle \geq -\left(\frac{1}{2\frac{c+1}{c-1}} + 1\right)^2\| \hat r_{T^c}\|_{w,1}^2 \|\mathcal{W}^{-1}(U-\hat U)\mathcal{W}^{-1}\|_{\infty,\infty}\\
&\displaystyle\geq_{(a)}  - \frac{9}{4}\ \frac{4\widetilde{\mathcal{R}}^2(\hat\alpha)}{\mbox{$(1-\frac{1}{c})^2$}\lambda^2} \|\mathcal{W}^{-1}(U-\hat U)\mathcal{W}^{-1}\|_{\infty,\infty},
\end{array}$$
where (a) follows since \eqref{eq:nocovariate:UnRestrictedSet} implies that $\widetilde{\mathcal{R}}(\hat\alpha)\geq\frac{\lambda}{2}\left(1-\frac{1}{c}\right)\|\hat r_{T^c}\|_{w,1}$ when $\|\hat r_{T^c}\|_{w,1} \geq 2\frac{c+1}{c-1}\|\hat r_T - r_0\|_{w,1}$. Rearranging the terms above and using (\ref{eq:nocovariate:UnRestrictedSet}) again we have
$$\frac{c_0}{N}\sum_{i=1}^N\left\{\frac{W(\alpha_0,Y_i)^\prime(\hat{r}-r_0)}{1+W(\alpha_0,Y_i)^\prime r_0}\right\}^2 \leq  \frac{9\widetilde{\mathcal{R}}^2(\hat\alpha)}{\mbox{$(1-\frac{1}{c})^2$}\lambda^2} \|\mathcal{W}^{-1}(\hat U - U)\mathcal{W}^{-1} \|_{\infty,\infty} + \widetilde{\mathcal{R}}(\hat\alpha).$$ 
Next suppose that  $\|\hat r_{T^c}\|_{w,1} < 2\frac{c+1}{c-1}\|\hat r_T - r_0\|_{w,1}$. Then we have 
$$\begin{array}{rl}
&\quad\displaystyle\frac{c_0}{N} \sum_{i=1}^N\left\{\frac{W(\alpha_0,Y_i)^\prime(\hat{r}-r_0)}{1+W(\alpha_0,Y_i)^\prime r_0}\right\}^2\\
&\displaystyle \leq_{(1)}\lambda(1+\frac{1}{c})\|\hat r_{T} - r_0\|_{w,1}+\widetilde{\mathcal{R}}(\hat\alpha)+\|\hat r - r_0\|_{w,1}^2 \|\mathcal{W}^{-1}(\hat U - U)\mathcal{W}^{-1} \|_{\infty,\infty}\\
&\displaystyle \leq_{(2)}\frac{\lambda(1+c)\sqrt{s}}{\sqrt{c_0}c\kappa_{\bar c}}\!\!\sqrt{\frac{c_0}{N} \!\!\sum_{i=1}^N\left\{\frac{W(\alpha_0,Y_i)^\prime(\hat{r}-r_0)}{1+W(\alpha_0,Y_i)^\prime r_0}\right\}^2}  \\
&\displaystyle\quad+\frac{s\left(2\frac{c+1}{c-1}+1\right)^2\|\mathcal{W}^{-1}(\hat U - U)\mathcal{W}^{-1} \|_{\infty,\infty} }{c_0\kappa_{\bar c}^2}\frac{c_0}{N} \sum_{i=1}^N\left\{\frac{W(\alpha_0,Y_i)^\prime(\hat{r}-r_0)}{1+W(\alpha_0,Y_i)^\prime r_0}\right\}^2  + \widetilde{\mathcal{R}}(\hat\alpha),
\end{array}$$
where (1) follows by (\ref{eq:thm1:main:main}) and (\ref{eq:DesignMatrixConnection}),  and (2) follows from the definition of the restricted eigenvalue, $\|\hat{r}_{T}-r_0\|_{w,1}\leq\sqrt{s}\|\hat{r}_{T}-r_0\|_{w,2}$ and $\|\hat{r}-r_0\|_{w,1}^2=(\|\hat{r}_{T^c}\|_{w,1}+\|\hat{r}_T-r_0\|_{w,1})^2\leq\left(2\frac{c+1}{c-1}+1\right)^2\|\hat{r}_T-r_0\|_{w,1}^2$. Note that $\frac{s\left(2\frac{c+1}{c-1}+1\right)^2\|\mathcal{W}^{-1}(\hat U - U)\mathcal{W}^{-1} \|_{\infty,\infty} }{c_0\kappa_{\bar{c}}^2} \leq_{(i)} \frac{Cs\|\hat\alpha-\alpha_0\|^2 K_{\alpha,2}}{\xi\kappa_{\bar c}^2 \min_{j\in[p]} w_j^2}\leq_{(ii)}C\delta_N$, where $$C\geq\frac{\left(2\frac{c+1}{c-1}+1\right)^2}{c_0}\max_{k,l\in[p]}\frac{1}{N}\sum_{i=1}^N5\left(1+\frac{|W_k(\alpha_0,Y_i)|}{1+W(\alpha_0,Y_i)^\prime r_0}\right)\left(1+\frac{|W_l(\alpha_0,Y_i)|}{1+W(\alpha_0,Y_i)^\prime r_0}\right).$$
Here the inequality (i) holds by Lemma \ref{lemma:MatrixConnection}, and the inequality (ii) holds by the condition in Theorem \ref{thm:main:FO}. Since $\delta_N\rightarrow0$, so for $N$ sufficiently large, $\frac{s \left(2\frac{c+1}{c-1}+1\right)^2\|\mathcal{W}^{-1}(\hat U - U)\mathcal{W}^{-1} \|_{\infty,\infty} }{c_0\kappa_{\bar{c}}^2} \leq \frac{1}{2}$.
Furthermore, note that 
$\frac{x^2}{2} \leq \lambda(1+\frac{1}{c})\frac{\sqrt{s}}{\sqrt{c_0}\kappa_{\bar{c}}} x + \widetilde{\mathcal{R}}(\hat \alpha) \Rightarrow x \leq 2\lambda(1+\frac{1}{c})\frac{\sqrt{s}}{\sqrt{c_0}\kappa_{\bar{c}}} + \sqrt{2}\widetilde{\mathcal{R}}^{1/2}(\hat \alpha)$ for $x=\sqrt{\frac{c_0}{N}\sum_{i=1}^N\left\{\frac{W(\alpha_0,Y_i)'(\hat{r}-r_0)}{1+W(\alpha_0,Y_i)'r_0}\right\}^2}$.
So the result holds. 
\end{proof}
\medskip
\begin{proof}[Proof of Corollary~\ref{cor:FO}]
According to \eqref{def:A:no-covariate}, $\mathcal{A}(\hat\alpha) =\prod_{j=1}^M [\hat\alpha_j - cv_{\delta_\alpha} \hat s_j, \hat\alpha_j + cv_{\delta_\alpha} \hat s_j]$, where $\hat s_j= \sqrt{\hat\alpha_j(1-\hat\alpha_j)}$ and $cv_{\delta_\alpha} := (1-\delta_\alpha)\mbox{-quantile of} \max_{j\in[M]} \left| \frac{1}{N}\sum_{i=1}^N \xi_i \frac{(Y_{ij} - \hat \alpha_j)}{\sqrt{\hat\alpha_j(1-\hat\alpha_j)}}\right|$ conditioning on $\{Y_i\}_{i\in[N]}$. 
Denote $\hat{s}=(\hat{s}_j)_{j\in[M]}\in\mathbb{R}^M$. Then for any $t>0$, 
\begin{equation}\label{eq:key:cor:no-cov}
\begin{array}{rl}
&\quad\displaystyle\mathbb{P}\bigg(\displaystyle\max_{\alpha\in\mathcal{A}(\hat{\alpha})}\norm{\alpha -\alpha_0}_2\geq t\bigg)\\
&\leq\mathbb{P}\bigg(\norm{\hat{\alpha}-\alpha_0}_2+cv_{\delta_{\alpha}}\sqrt{\sum_{j\in[M]}\hat{s}_j^2}\geq t\bigg)\\
&\displaystyle\leq\mathbb{P}(\norm{\hat{\alpha}-\alpha_0}_2\geq t/2)+\mathbb{P}(cv_{\delta_{\alpha}}\sqrt{M/4}\geq t/2),
\end{array}
\end{equation}
where the second inequality above follows by triangular inequality and the fact that $\hat{s}_j^2\leq\frac{1}{4}$ for all $j\in[M]$. To bound $\mathbb{P}(\norm{\hat{\alpha}-\alpha_0}_2\geq t/2)$ of \eqref{eq:key:cor:no-cov}, note that by Proposition \ref{prop:concentration:no covariate}, with probability $\geq1-\delta$, $\norm{\hat{\alpha}-\alpha_0}_2\leq c_1\sqrt{\frac{M\log(1/\delta)}{N}}$, where $c_1$ is some absolute constant. To bound $\mathbb{P}(cv_{\delta_{\alpha}}\sqrt{M/4}\geq t/2)$ of \eqref{eq:key:cor:no-cov}, note that $\xi_i\overset{i.i.d.}{\sim}\mathcal{N}(0,1)$, therefore conditional on $\{Y_i\}_{i\in[N]}$, 
$\frac{1}{N}\sum_{i=1}^N\xi_i\frac{(Y_{ij}-\hat{\alpha}_j)}{\sqrt{\hat{\alpha}_j(1-\hat{\alpha}_j)}}\sim\mathcal{N}\left(0,\frac{1}{N^2}\sum_{i=1}^N\frac{(Y_{ij}-\hat{\alpha}_j)^2}{\hat{\alpha}_j(1-\hat{\alpha}_j)}\right)$ holds for all $j\in[M]$, so $\mathbb{E}\left[\max_{j\in[M]}\left|\frac{1}{N}\sum_{i=1}^N\xi_i\frac{(Y_{ij}-\hat{\alpha}_j)}{\sqrt{\hat{\alpha}_j(1-\hat{\alpha}_j)}}\right|\right]\leq\sqrt{\frac{2\sigma^2\log(2M)}{N}}$, where $\sigma^2=\max_{j\in[M]}\frac{1}{N}\sum_{i=1}^N\frac{(Y_{ij}-\hat{\alpha}_j)^2}{\hat{\alpha}_j(1-\hat{\alpha}_j)}$.
By Lemma~\ref{lemma:borell}, for some absolute constant $C_0$, with probability $1-\delta$ we have $\max_{j\in[M]}\left|\frac{1}{N}\sum_{i=1}^N\xi_i\frac{(Y_{ij}-\hat{\alpha}_j)}{\sqrt{\hat{\alpha}_j(1-\hat{\alpha}_j)}}\right|\leq\mathbb{E}\left[\max_{j\in[M]}\left|\frac{1}{N}\sum_{i=1}^N\xi_i\frac{(Y_{ij}-\hat{\alpha}_j)}{\sqrt{\hat{\alpha}_j(1-\hat{\alpha}_j)}}\right|\right]\!\!+\sigma\sqrt{\frac{2\log\left(2/\delta\right)}{N}}\leq C_0\sqrt{\frac{\log(M/\delta)}{N}}\sqrt{\max_{j\in[M]}\frac{1}{N}\sum_{i=1}^N\frac{(Y_{ij}-\hat{\alpha}_j)^2}{\hat{\alpha}_j(1-\hat{\alpha}_j)}}$. Let $\delta\leq\delta_{\alpha}$ and set $t=C\sqrt{\frac{M\log(M/\delta)}{N}}$ for some absolute constant $C$ in (\ref{eq:key:cor:no-cov}), then with probability $1-2\delta$, $\max_{\alpha\in\mathcal{A}(\hat{\alpha})}\norm{\alpha -\alpha_0}_2\leq C\sqrt{\frac{M\log(M/\delta)}{N}}$. Additionally, for any $t>0$, by triangular inequality, 
\begin{equation}\label{ineq:concentration:l1:no-cov}
\begin{array}{rl}
&\displaystyle\quad\mathbb{P}\big(\max_{\alpha\in\mathcal{A}(\hat{\alpha})}\|\alpha-\alpha_0\|_1\geq t\big)\\
&\displaystyle\leq\mathbb{P}(\|\hat{\alpha}-\alpha_0\|_1\geq t/2)+\mathbb{P}(cv_{\delta_{\alpha}}\|\hat{s}\|_1\geq t/2)\\
&\displaystyle\leq_{(a)}\mathbb{P}(\|\hat{\alpha}-\alpha_0\|_1\geq t/2)+\mathbb{P}(cv_{\delta_{\alpha}}(M/2)\geq t/2),
\end{array}
\end{equation}
where (a) follows because $\|\hat{s}\|_1\leq M/2$.
By Proposition \ref{prop:concentration:no covariate}, with probability $1-\delta$, we have $\norm{\hat\alpha-\alpha_0}_1\leq 
c_1\frac{M\sqrt{\log(1/\delta)}}{\sqrt{N}}$ for some constant $c_1>0$. Let $\delta\leq\delta_{\alpha}$ and set $t=CM\sqrt{\frac{\log (M/\delta)}{N}}$ for some absolute constant $C$, following (\ref{ineq:concentration:l1:no-cov}), with probability $1-2\delta$, we have $\max_{\alpha\in\mathcal{A}(\hat{\alpha})}\|\alpha-\alpha_0\|_1\leq CM\sqrt{\frac{\log(M/\delta)}{N}}$. Additionally, we choose $\lambda$ satisfying the condition of Theorem \ref{thm:main:FO}, so that for some absolute constant $\bar{C}$, $\lambda \leq \bar{C}N^{-1/2}\sqrt{\log(p/\delta)}$, then under the condition $\kappa_{\bar{c}}^{-1}+\bar K_2 + K_{\alpha,1} + K_{\alpha,2}/\xi^2 \leq C'$, we have $\widetilde{\mathcal{R}}(\hat\alpha) \leq \bar{C}'M\sqrt{\frac{\log(M/\delta)}{N}} \sqrt{\frac{\log(M/\delta)}{N}}+ \bar{C}'\frac{M\log(M/\delta)}{N}$ for some constant $\bar{C}'$. By Theorem \ref{thm:main:FO}, for an absolute constant $C$, $\sqrt{\frac{1}{N}\sum_{i=1}^N \frac{\{W(\alpha_0,Y_i)'(\hat r^{FO} - r_0 )\}^2}{(1+W(\alpha_0,Y_i)'r_0)^2}}\leq C\sqrt{\frac{s\log (p/\delta)}{N}}+ C\sqrt{\frac{M\log(M/\delta)}{N}}$ holds with probability $1-C\delta-C\delta_\alpha-C\epsilon$, hence the result follows since $p=2^M-M-1$. \end{proof}

\subsection{Additional Results and Technical Lemmas for Section \ref{Sec:Rates:no-covariates}}\label{appendix:no-covariate}
\begin{lemma}[Borell-TIS inequality \citep{borell1975brunn}]\label{lemma:borell} 
Let $X$ be a centered Gaussian random variable on $\mathbb{R}^n$ and $\displaystyle\sigma^2:=\max_{i\in[n]}\mathbb{E}[X_i^2]$. Then $\displaystyle\mathbb{P}\bigg(\bigg|\max_{i\in[n]}X_i-\mathbb{E}\big[\max_{i\in[n]}X_i\big]\bigg|>t\bigg)\leq 2e^{-t^2/2\sigma^2}$ holds for every $t>0$. \end{lemma}

\begin{lemma}\label{lemma:MinMax:Claim1}
Suppose Assumptions \ref{ass:iid-min-den}, \ref{ass:smoothness-assumptions} hold. Then $$\begin{array}{rl}
&\quad\displaystyle\min_{\alpha \in \mathcal{A}(\hat\alpha)} Q^{FO}(\alpha,r_0)-Q^{FO}(\alpha_0,r_0) \\
&\displaystyle\geq- C\max_{\alpha \in \mathcal{A}(\hat\alpha)}\|\alpha-\alpha_0\|_1 \left\| \frac{1}{N}\sum_{i=1}^N \frac{\nabla_\alpha W(\alpha_0,Y_i)'r_0}{1+W(\alpha_0,Y_i)'r_0} \right\|_\infty -C\frac{K_{\alpha,2}}{\xi^2}\max_{\alpha \in \mathcal{A}(\hat\alpha)} \|\alpha - \alpha_0\|^2
\end{array}$$ 
holds with probability $1-\epsilon-\delta_{\alpha}$, where $C$ is an absolute constant independent of $s,N,M$. Moreover, under Assumption \ref{ass:moment-assumptions}, $\left\| \frac{1}{N}\sum_{i=1}^N \frac{\nabla_\alpha W(\alpha_0,Y_i)'r_0}{1+W(\alpha_0,Y_i)'r_0} \right\|_\infty \leq \Phi^{-1}(1-\delta/2M) C'\bar K_2\sqrt{1/N}$ holds with probability $1-\delta((1+\mathrm{o}(1)))$ for some universal constant $C'$.
\end{lemma}
\begin{proof}[Proof of Lemma \ref{lemma:MinMax:Claim1}]
Recall $Q^{FO}(\alpha,r)= \frac{1}{N}\sum_{i=1}^N \log(1+ W^{FO}(\alpha,Y_i)'r)$, where $W^{FO}(\alpha,y)=W(\hat \alpha,y)+\nabla_\alpha W(\hat\alpha,y)(\alpha-\hat\alpha)$. Suppose all conditions of Assumption~\ref{ass:smoothness-assumptions} hold and condition (iii) of Assumption~\ref{ass:iid-min-den} holds. Then for any $\alpha\in\mathcal{A}(\hat{\alpha})$,
\begin{equation}\label{eq:Main:MinMax}
\begin{array}{rl}
&\quad Q^{FO}(\alpha,r_0)-Q^{FO}(\alpha_0,r_0) \\
&\displaystyle = \frac{1}{N}\sum_{i=1}^N \log\left( 1 + (\alpha - \alpha_0)'\frac{\nabla_\alpha W(\hat\alpha,Y_i)'r_0}{1+W^{FO}(\alpha_0,Y_i)'r_0} \right).
\end{array}
\end{equation}
By the mean value theorem, for some $\bar{\alpha}$, $\tilde \alpha$ in the line segment between $\hat\alpha$ and $\alpha_0$ we have $(\alpha-\alpha_0)'\nabla_\alpha W(\hat\alpha,Y_i)'r_0 = (\alpha-\alpha_0)'\nabla_\alpha W(\alpha_0,Y_i)'r_0 + (\alpha-\alpha_0)'\{\nabla_\alpha^2 W(\bar{\alpha},Y_i)'r_0\}(\hat\alpha-\alpha_0)$, and 
$$ \begin{array}{rl}
W^{FO}(\alpha_0,Y_i)'r_0 &= W(\hat \alpha,Y_i)'r_0 + (\alpha_0-\hat\alpha)'\nabla_\alpha W(\hat\alpha,Y_i)'r_0\\
&\displaystyle =_{(1)} W(\alpha_0,Y_i)'r_0 -\frac{1}{2}(\alpha_0-\hat\alpha)'\{\nabla_\alpha^2 W(\tilde\alpha,Y_i)'r_0\}(\alpha_0-\hat\alpha)\\
\\
\displaystyle \frac{1}{1+W^{FO}(\alpha_0,Y_i)'r_0} &\displaystyle= \frac{1}{1+W(\alpha_0,Y_i)'r_0}+\frac{W(\alpha_0,Y_i)'r_0 -W^{FO}(\alpha_0,Y_i)'r_0}{(1+W(\alpha_0,Y_i)'r_0)(1+W^{FO}(\alpha_0,Y_i)'r_0)}\\
&\displaystyle = \frac{1}{1+W(\alpha_0,Y_i)'r_0}\left(1 -\Delta R_0\right)
% R(\alpha_0,\hat\alpha,Y_i)
\end{array}$$
where (1) follows from the Taylor expansion of $W(\alpha_0,Y_i)'r_0$ around $\hat{\alpha}$, and $\Delta R_0:=-1/2[(\alpha_0-\hat\alpha)'\nabla_\alpha^2 W(\tilde\alpha,Y_i)'r_0(\alpha_0-\hat\alpha)]/\{1+W^{FO}(\alpha_0,Y_i)'r_0\}$. Thus we have
\begin{equation}\label{eq:Aux:MinMax}\begin{array}{rl}
&\quad\displaystyle\frac{(\alpha - \alpha_0)'\nabla_\alpha W(\hat\alpha,Y_i)'r_0}{1+W^{FO}(\alpha_0,Y_i)'r_0}\\
&\displaystyle = \frac{(\alpha - \alpha_0)'\nabla_\alpha W(\alpha_0,Y_i)'r_0}{1+W(\alpha_0,Y_i)'r_0}(1-\Delta R_0) \\
&\ \displaystyle +\frac{(\alpha-\alpha_0)'\{\nabla_\alpha^2 W(\bar{\alpha},Y_i)'r_0\}(\hat\alpha-\alpha_0)(1-\Delta R_0)}{1+W(\alpha_0,Y_i)'r_0}\\
&\displaystyle = \frac{(\alpha - \alpha_0)'\nabla_\alpha W(\alpha_0,Y_i)'r_0}{1+W(\alpha_0,Y_i)'r_0} +\widetilde R(\alpha,\hat\alpha,Y_i),\ \mbox{where}
\end{array}
\end{equation}
$$\begin{array}{rl}
|\widetilde R(\alpha,\hat\alpha,Y_i)| &\displaystyle\leq_{(i)}  \|\alpha-\alpha_0\| \frac{M^{1/2}K_{\alpha,1}}{\xi} |\Delta R_0| +\max_{\tilde\alpha \in \mathcal{A}(\hat\alpha)} 2\|\tilde\alpha-\alpha_0\|^2 \frac{K_{\alpha,2}}{\xi} (1+|\Delta R_0|)\\
&\displaystyle\leq_{(ii)} 3\max_{\tilde\alpha \in \mathcal{A}(\hat\alpha)} \|\tilde\alpha-\alpha_0\|^2 \frac{K_{\alpha,2}}{\xi^2}<_{(iii)}\frac{3}{16}.
\end{array}$$
Here inequality (i) above uses condition (iii) of Assumption~\ref{ass:iid-min-den} and (2), (4) of Assumption~\ref{ass:smoothness-assumptions}. Since $\hat{\alpha}\in\mathcal{A}(\hat{\alpha})$, then (2), (4), (5) of Assumption~\ref{ass:smoothness-assumptions} imply $|v'\nabla_\alpha^2 W(\tilde\alpha,Y_i)'r_0| \leq K_{\alpha,2}\|v\| $ and $ | W^{FO}(\alpha_0,Y_i)'r_0 - W(\alpha_0,Y_i)'r_0| \leq \frac{1}{16} \xi$, which then imply that $|\Delta R_0|\leq \frac{16}{15}\max_{\tilde \alpha \in \mathcal{A}(\hat \alpha)} \|\tilde\alpha-\alpha_0\|^2 \frac{K_{\alpha,2}}{\xi}$, and (5) of Assumption~\ref{ass:smoothness-assumptions} further implies inequality (ii). inequality (iii) follows from (5) of Assumption~\ref{ass:smoothness-assumptions}. Next, note that $\log(1+x)\geq \frac{x}{1+x} = x - \frac{x^2}{1+x}$ and conditions (2), (4), (5) of Assumption~\ref{ass:smoothness-assumptions} imply that $\left| \frac{(\alpha - \alpha_0)'\nabla_\alpha W(\alpha_0,Y_i)'r_0}{1+W(\alpha_0,Y_i)'r_0} +\widetilde R(\alpha,\hat\alpha,Y_i))\right| \leq 1/2$. Further by (\ref{eq:Main:MinMax}) and (\ref{eq:Aux:MinMax}), for a universal constant $C$,
$$\begin{array}{rl}
&\quad Q^{FO}(\alpha,r_0)-Q^{FO}(\alpha_0,r_0) \\
&\displaystyle \geq  \frac{1}{N}\sum_{i=1}^N \frac{(\alpha - \alpha_0)'\nabla_\alpha W(\alpha_0,Y_i)'r_0}{1+W(\alpha_0,Y_i)'r_0} +\widetilde R(\alpha,\hat\alpha,Y_i) - \!\frac{2}{N}\sum_{i=1}^N \left(\frac{(\alpha - \alpha_0)'\nabla_\alpha W(\alpha_0,Y_i)'r_0}{1+W(\alpha_0,Y_i)'r_0} +\widetilde R(\alpha,\hat\alpha,Y_i))\right)^2 \\
&\displaystyle \geq -C\max_{\alpha \in \mathcal{A}(\hat\alpha)} \|\alpha-\alpha_0\|_1 \left\|\frac{1}{N}\sum_{i=1}^N \frac{\nabla_\alpha W(\alpha_0,Y_i)'r_0}{1+W(\alpha_0,Y_i)'r_0}  \right\|_\infty - C\frac{K_{\alpha,2}}{\xi^2}\max_{\alpha \in \mathcal{A}(\hat\alpha)} \|\alpha-\alpha_0\|^2,
\end{array}$$
where the second inequality follows from H\"older's inequality, the bound derived above for $|\widetilde R(\alpha,\hat\alpha,Y_i)|$ and Assumption~\ref{ass:smoothness-assumptions}. The first statement then holds since (iii) of Assumption~\ref{ass:iid-min-den} holds with probability $1-\epsilon$ and the conditions of Assumption~\ref{ass:smoothness-assumptions} hold with probability $1-\delta_{\alpha}$. The second statement of the Lemma follows from an application of self-normalized moderate deviation theory.  Let $\tilde X_{ij} := \nabla_{\alpha_j} W(\alpha_0,Y_i)'r_0 / \{1+W(\alpha_0,Y_i)'r_0\}$. so under Assumptions \ref{ass:iid-min-den}, \ref{ass:moment-assumptions} we have $\min_{j\in[M]}\mathbb{E}[\tilde X_{ij}^2 ] \geq \bar\psi_2^2 $, $\max_{j\in[M]}\mathbb{E}[ |\tilde X_{ij}|^3 ]^{1/3} \leq \bar K_3$ and $\sqrt{\frac{\log^3p}{N}}\leq\sqrt{\frac{M^3}{N}}\leq\delta_N$, thus 
$$\begin{array}{rl}
&\displaystyle \mathbb{P} \left( \left| \frac{1}{N}\sum_{i=1}^N \tilde X_{ij} \right| > \frac{\Phi^{-1}(1-\delta/2M)}{N^{1/2}} \sqrt{\frac{1}{N}\sum_{i=1}^N \tilde X_{ij}^2}\mbox{ for some }j\in[M]\right)  \\
&\displaystyle \leq  \sum_{j\in [M]} \mathbb{P} \left( \left|\frac{1}{N}\sum_{i=1}^N\tilde X_{ij} \right| >\frac{\Phi^{-1}(1-\delta/2M)}{N^{1/2}}\sqrt{\frac{1}{N}\sum_{i=1}^N \tilde X_{ij}^2}\right)\leq \delta(1+\mathrm{o}(1)),
\end{array}$$
where the first inequality follows from the union bound, and the second inequality follows by applying Theorem 2.3 in \citet{jing2003self}. Hence the result follows.
\end{proof}

\begin{lemma}\label{lemma:Score:nocovariate}
Define $\displaystyle\beta_N := \max_{j\in[p]}\frac{1}{w_j}\frac{1}{N}\sum_{i=1}^N \frac{|W_j(\alpha_0,Y_i)|}{1+W(\alpha_0,Y_i)'r_0}$.
Suppose Assumption~\ref{ass:iid-min-den} and (i)--(iv):\\
(i) $\displaystyle\max_{\alpha \in\mathcal{A}(\hat{\alpha}), j\in[p]}|v'\nabla_\alpha W_j(\alpha,Y_i)| \leq M^{1/2}K_{\alpha,1}\|v\|$, (ii) $\max_{\tilde \alpha \in \mathcal{A}(\hat\alpha), j\in [p]} |v'\nabla_\alpha^2 W_j(\tilde\alpha,Y_i)u| \leq K_{\alpha,2}\|v\|\ \|u\| $, (iii) $\displaystyle|v'\nabla_\alpha W(\alpha_0,Y_i)'r_0| \leq M^{1/2}K_{\alpha,1}\|v\|$, $\max_{\tilde \alpha \in \mathcal{A}(\hat\alpha)} |v'\nabla_\alpha^2 W(\tilde\alpha,Y_i)'r_0v| \leq K_{\alpha,2}\|v\|^2$, \\
(iv) $\displaystyle\|\alpha-\alpha_0\|M^{1/2}K_{\alpha,1}/\xi < 1/4$,  $4\max_{\tilde\alpha \in \mathcal{A}(\hat\alpha)}\|\tilde \alpha-\alpha_0\|^2 K_{\alpha,2}/\xi \leq 1/4$, we have 
$$\begin{array}{rl}
&\displaystyle\left\|\frac{{\rm diag}^{-1}(w)}{N}\!\!\sum_{i=1}^N \frac{W(\hat\alpha,Y_i)}{1+W(\hat\alpha,Y_i)'r_0}\right\|_\infty\!\!\!\!\!\!\leq \left\|\frac{{\rm diag}^{-1}(w)}{N}\!\!\!\sum_{i=1}^N\frac{W(\alpha_0,Y_i)}{1+W(\alpha_0,Y_i)'r_0}\right\|_\infty\!\!\!\!\!\!+\!\frac{C(1+\beta_N)M^{1/2}K_{\alpha,1}\|\hat\alpha-\alpha_0\|}{\xi} \\
&\displaystyle\left\|\frac{{\rm diag}^{-1}(w)}{N}\!\!\sum_{i=1}^N\!\!\frac{W^{FO}(\alpha_0,Y_i)}{1+W^{FO}(\alpha_0,Y_i)'r_0}\right\|_\infty\!\!\!\!\!\!\leq \left\|\frac{{\rm diag}^{-1}(w)}{N}\!\!\sum_{i=1}^N \frac{W(\alpha_0,Y_i)}{1+W(\alpha_0,Y_i)'r_0}\right\|_\infty\!\!\!\!\!\!+\!\frac{C (1+\beta_N)K_{\alpha,2}\|\hat\alpha-\alpha_0\|^2}{\xi}
\end{array}$$
where $C$ is an absolute constant in the above, $w=(w_1,\ldots,w_p)$, and for each $j\in[p]$, $w_j$ is the weight for \eqref{eq:PMLE-no-covariates} and \eqref{eq:PMLE-no-covariates-bilinear}. Moreover, with probability $1-\delta(1+\mathrm{o}(1))$,  $\left\|\frac{{\rm diag}^{-1}(w)}{N}\!\!\sum_{i=1}^N \!\!\frac{W(\alpha_0,Y_i)}{1+W(\alpha_0,Y_i)'r_0}\right\|_\infty\!\!\!\!\!\!\leq\frac{\Phi^{-1}(1-\delta/2p)}{\sqrt{N}}\max_{j\in[p]}\frac{1}{w_j}\sqrt{\frac{1}{N}\sum_{i=1}^N \frac{W_j^2(\alpha_0,Y_i)}{(1+W(\alpha_0,Y_i)'r_0)^2}}$ holds under Assumptions \ref{ass:iid-min-den} and \ref{ass:moment-assumptions}.
\end{lemma}

\begin{proof}[Proof of Lemma \ref{lemma:Score:nocovariate}]
Note that by Taylor expansion, for each $j\in [p], i\in[N]$ there exist $\tilde \alpha^{ji}$, $\tilde{\alpha}$ in the line segment between $\hat \alpha$ and $\alpha_0$ such that
$W_j(\hat\alpha,Y_i) = W_j(\alpha_0,Y_i) + (\hat\alpha-\alpha_0)'\nabla_\alpha W_j(\tilde \alpha^{ji},Y_i)$. Thus 
$$\begin{array}{rl}
\displaystyle\frac{W_j(\hat\alpha,Y_i)}{1+W(\hat\alpha,Y_i)'r_0}&\displaystyle= \frac{W_j(\alpha_0,Y_i) }{1+W(\alpha_0,Y_i)'r_0} \left( 1 -  \frac{(\hat\alpha-\alpha_0)'\nabla_\alpha W(\tilde \alpha,Y_i)'r_0}{1+W(\alpha_0,Y_i)'r_0 + (\hat\alpha-\alpha_0)'\nabla_\alpha W(\tilde \alpha,Y_i)'r_0}\right)\\
&\displaystyle\quad + \frac{(\hat\alpha-\alpha_0)'\nabla_\alpha W_j(\tilde \alpha^{ji},Y_i)}{1+W(\alpha_0,Y_i)'r_0 + (\hat\alpha-\alpha_0)'\nabla_\alpha W(\tilde \alpha,Y_i)'r_0}\\
&\displaystyle = \frac{W_j(\alpha_0,Y_i) }{1+W(\alpha_0,Y_i)'r_0} + \widetilde R(\hat\alpha,Y_i),
\end{array}$$ 
where $\displaystyle|\widetilde R(\hat\alpha,Y_i)| \leq 2\left(1+\frac{|W_j(\alpha_0,Y_i)|}{1+W(\alpha_0,Y_i)'r_0}\right) \frac{\|\hat\alpha-\alpha_0\| M^{1/2}K_{\alpha,1}}{\xi}$ following from (iii) and (iv). So the first inequality holds.

Next we prove the second inequality. By Taylor expansion, for each $j\in [p], i\in[n]$ there is $\bar\alpha_{ji}$ in the line segment between $\hat \alpha$ and $\alpha_0$ such that
$W_j(\alpha_0,Y_i) = W_j(\hat{\alpha},Y_i) + (\alpha_0-\hat{\alpha})'\nabla_\alpha W_j(\hat{\alpha},Y_i) + \frac{1}{2}(\alpha_0-\hat\alpha)'\nabla_\alpha^2 W_j(\bar{\alpha}_{ji},Y_i)(\alpha_0-\hat\alpha)$. Recall that $W^{FO}(\alpha,y)=W(\hat \alpha,y)+\nabla_\alpha W(\hat\alpha,y)(\alpha-\hat\alpha)$. Thus 
$$ \begin{array}{rl}
W^{FO}_j(\alpha_0,Y_i) &= W_j(\hat \alpha,Y_i) + (\alpha_0-\hat\alpha)'\nabla_\alpha W_j(\hat\alpha,Y_i)\\
& = W_j(\alpha_0,Y_i)-\frac{1}{2}(\alpha_0-\hat\alpha)'\nabla_\alpha^2 W_j(\bar{\alpha}_{ji},Y_i)(\alpha_0-\hat\alpha)\\
\displaystyle\frac{1}{1+W^{FO}(\alpha_0,Y_i)'r_0} &\displaystyle= \frac{1}{1+W(\alpha_0,Y_i)'r_0}+\frac{W(\alpha_0,Y_i)'r_0 -W^{FO}(\alpha_0,Y_i)'r_0}{(1+W(\alpha_0,Y_i)'r_0)(1+W^{FO}(\alpha_0,Y_i)'r_0)}.\ \mbox{So}
\end{array}$$ 
\begin{equation}\label{eq:Wj:WjFO:nocovariate}  \frac{W^{FO}_j(\alpha_0,Y_i) }{1+W^{FO}(\alpha_0,Y_i)'r_0} =  \frac{W_j(\alpha_0,Y_i)}{1+W(\alpha_0,Y_i)'r_0} + \widetilde R_0.
\end{equation}
Using Assumption~\ref{ass:iid-min-den} and conditions (i) -- (iv), 
$$\begin{array}{rl}
\displaystyle|\widetilde R_0| 
& \displaystyle\leq C\left(1+\frac{|W_j(\alpha_0,Y_i)| }{1+W(\alpha_0,Y_i)'r_0}\right) \|\hat\alpha - \alpha_0\|^2 K_{\alpha,2}/\xi\ \ \mbox{for some absolute constant }C.
\end{array}$$
The third inequality follows by applying the self-normalized moderate deviation theory. Let $\tilde X_{ij} := W_j(\alpha_0,Y_i) / \{1+W(\alpha_0,Y_i)'r_0\}$ so under Assumptions \ref{ass:iid-min-den} and \ref{ass:moment-assumptions}, ${\displaystyle\min_{j\in[p]}}\mathbb{E}[\tilde X_{ij}^2 ]^{1/2}\geq \bar\psi_2$, ${\displaystyle\max_{j\in[p]}}\mathbb{E}[ |\tilde X_{ij}|^3 ]^{1/3} \leq \bar K_3$, and  $\sqrt{\frac{\log^3p}{N}}\leq\sqrt{\frac{M^3}{N}}\leq\delta_N$, so we have 
$$\begin{array}{rl}
&\displaystyle \mathbb{P} \left( \left| \frac{1}{N}\sum_{i=1}^N \frac{\tilde X_{ij}}{w_j} \right| > \frac{\Phi^{-1}(1-\delta/(2p))}{w_j N^{1/2}} \sqrt{\frac{1}{N}\sum_{i=1}^N \tilde X_{ij}^2} \mbox{ for some }j\in[p] \right)\\
&\displaystyle \leq  \sum_{j\in [p]} \mathbb{P} \left( \left|\frac{1}{N}\sum_{i=1}^N\frac{\tilde X_{ij}}{w_j} \right| > \frac{\Phi^{-1}(1-\delta/(2p))}{w_jN^{1/2}} \sqrt{\frac{1}{N}\sum_{i=1}^N \tilde X_{ij}^2 } \right)\leq \delta(1+\mathrm{o}(1)).
\end{array}$$
where the first inequality follows from the union bound, and the second inequality follows by applying Theorem 2.3 of \citet{jing2003self}.
\end{proof}

\begin{lemma}\label{lemma:MatrixConnection}
Suppose conditions (i) -- (iv) of Lemma~\ref{lemma:Score:nocovariate} hold. Let $U:= \frac{1}{N}\sum_{i=1}^N \frac{W(\alpha_0,Y_i)W(\alpha_0,Y_i)'}{(1+W(\alpha_0,Y_i)'r_0)^2}$ and 
$\hat U:= \frac{1}{N}\sum_{i=1}^N \frac{W^{FO}(\alpha_0,Y_i)  W^{FO}(\alpha_0,Y_i)'}{(1+W^{FO}(\alpha_0,Y_i)'r_0)^2}$. 
Then $\left| U_{lk} - \hat U_{lk} \right| \leq 5\|\hat\alpha - \alpha_0\|^2 \frac{K_{\alpha,2}}{\xi} \max_{l,k\in[p]}\frac{1}{N}\sum_{i=1}^N(1+|W_{il}|)(1+|W_{ik}|)$, where $W_i:=W(\alpha_0,Y_i)/(1+W(\alpha_0,Y_i)'r_0)$.
\end{lemma}
\begin{proof}[Proof of Lemma \ref{lemma:MatrixConnection}]
Note that by (\ref{eq:Wj:WjFO:nocovariate}) we have $\frac{W^{FO}(\alpha_0,Y_i) }{1+W^{FO}(\alpha_0,Y_i)'r_0} =  \frac{W(\alpha_0,Y_i)}{1+W(\alpha_0,Y_i)'r_0} + \widetilde R(\hat\alpha,\alpha_0,Y_i)$,
where $|\widetilde R_j(\hat{\alpha},\alpha_0,Y_i)| \leq  2\left(1+\frac{|W_j(\alpha_0,Y_i)| }{1+W(\alpha_0,Y_i)'r_0}\right) \|\hat\alpha - \alpha_0\|^2 K_{\alpha,2}/\xi$, for any $j\in[p]$. Define 
$\hat W_i := \frac{W^{FO}(\alpha_0,Y_i) }{1+W^{FO}(\alpha_0,Y_i)'r_0}$, $\tilde R_{il}:=\tilde R_l(\hat{\alpha},\alpha_0,Y_i)$. Thus 
$\left|\hat W_{il}\hat W_{ik} -W_{il}W_{ik}\right| \leq \left| W_{il}\tilde R_{ik} + \tilde R_{il} W_{ik} + \tilde R_{il}\tilde R_{ik}\right|$. By (iv) of Lemma~\ref{lemma:Score:nocovariate}, $\frac{1}{N}\sum_{i=1}^N|\tilde{R}_{ik}|\left(1+\frac{|W_l(\alpha_0,Y_i)|}{1+W(\alpha_0,Y_i)'r_0}\right)2\|\hat\alpha - \alpha_0\|^2 \frac{K_{\alpha,2}}{\xi} \leq \frac{1}{2N}\sum_{i=1}^N|\tilde{R}_{ik}|\big(1+\frac{|W_l(\alpha_0,Y_i)|}{1+W(\alpha_0,Y_i)'r_0}\big)$, so we have 
$$\frac{1}{N}\sum_{i=1}^N |\tilde R_{il}\tilde R_{ik}|\leq\frac{1}{2N}\sum_{i=1}^N (1+|W_{il}|)|\tilde R_{ik}|.$$ 
So the upper bound for $\left| U_{lk} - \hat U_{lk} \right|$ holds by the upper bound of $\tilde{R}_{il}$ for any $l\in[p]$.
\end{proof}

\begin{lemma}\label{lem:diff-restricted}
If {$\displaystyle\norm{{\rm diag}^{-1}(w)\frac{1}{N}\sum_{i=1}^N\frac{W(\hat\alpha,Y_i)}{1+W(\hat\alpha,Y_i)^\prime r_0}}_\infty\leq\frac{\lambda}{c}$} for some constant $c>1$, then for $\nu=\hat r^{PI}-r_0$, $T$ as the support of $r_0$, we have $\displaystyle\norm{\nu_{T^c}}_{w,1}\leq\frac{c+1}{c-1}\norm{\nu_T}_{w,1}$. 
\end{lemma}
\begin{proof}[Proof of Lemma \ref{lem:diff-restricted}]
For a given $\lambda$, $\frac{1}{N}\sum_{i=1}^N\left[\log\left(1+W(\hat\alpha,Y_i)^\prime\hat r^{PI}\right)-\log\left(1+W\left(\hat\alpha,Y_i\right)^\prime r_0\right)\right]\geq\lambda\left(\norm{\hat r^{PI}}_{w,1}-\norm{r_0}_{w,1}\right)$.
Let $\hat Q(\alpha,r)=\frac{1}{N}\sum_{i=1}^N\log\left(1+W\left(\alpha,Y_i\right)^\prime r\right)$, we can write the above inequality as $\hat Q(\hat\alpha,\hat r^{PI})-\hat Q(\hat\alpha,r_0)\geq\lambda\left(\norm{\hat r^{PI}}_{w,1}-\norm{r_0}_{w,1}\right)$,
so using concavity of $\hat Q$ in parameter $r$, we have $\lambda\left(\norm{r_0}_{w,1}-\norm{\hat r^{PI}}_{w,1}\right)\geq\hat Q\left(\hat\alpha,r_0\right)-\hat Q\left(\hat\alpha,\hat r^{PI}\right)\geq\nabla_r\hat Q(\hat\alpha,r_0)^T(r_0-\hat r^{PI})\geq_{(a)}-\frac{\lambda}{c}\norm{\nu}_{w,1}$,
where (a) follows from H\"older's inequality. Note that $\norm{\hat r^{PI}}_{w,1}=\norm{r_0+\nu_T}_{w,1}+\norm{\nu_{T^c}}_{w,1}\geq\norm{r_0}_{w,1}-\norm{\nu_T}_{w,1}+\norm{\nu_{T^c}}_{w,1}$, so $\lambda(\norm{r_0}_{w,1}-\norm{\hat r^{PI}}_{w,1})\leq\lambda\left(\norm{\nu_T}_{w,1}-\norm{\nu_{T^c}}_{w,1}\right)$. Thus $\lambda\left(\norm{\nu_T}_{w,1}-\norm{\nu_{T^c}}_{w,1}\right)\geq-\frac{\lambda}{c}\norm{\nu}_{w,1}=-\frac{\lambda}{c}\left(\norm{\nu_T}_{w,1}+\norm{\nu_{T^c}}_{w,1}\right)$,
implying $\frac{c+1}{c-1}\norm{\nu_T}_{w,1}\geq\norm{\nu_{T^c}}_{w,1}$.
\end{proof}

\begin{lemma}\label{lemma:fast-rate:estimated-alpha}
    Suppose Assumptions \ref{ass:iid-min-den}  and \ref{ass:moment-assumptions} hold. Suppose $\norm{{\rm diag}^{-1}(w)\frac{1}{N}\sum_{i=1}^N\frac{W(\hat\alpha,Y_i)}{1+W(\hat\alpha,Y_i)^\prime r_0}}_\infty\leq\frac{\lambda}{c}$ for some $c>1$ and $\min_{i\in[N]}1+W(\hat\alpha,Y_i)^\prime r_0\geq \underline{c}\xi$ for an absolute constant $\underline{c}$. Then we have $$\frac{c_0}{N}\sum_{i=1}^N\left\{\frac{W(\hat\alpha,Y_i)^\prime(\hat{r}^{PI}-r_0)}{(1+W(\hat\alpha,Y_i)^\prime r_0)}\right\}^2\leq2\lambda\frac{c+1}{c-1}\norm{\hat{r}_T^{PI}-r_0}_{w,1},$$ 
    where $c_0>0$ is a universal constant.
\end{lemma}
\begin{proof}[Proof of Lemma \ref{lemma:fast-rate:estimated-alpha}]
Let $T=\mathrm{supp}(r_0)$ and $f(Y_i,\alpha,r)=1+W(\alpha,Y_i)'r$. 
By optimality of $\hat{r}^{PI}$, we have $\frac{1}{N}\sum_{i=1}^N\log f(Y_i,\hat\alpha,
\hat{r}^{PI})-\lambda\norm{\hat{r}^{PI}}_{w,1}\geq\frac{1}{N}\sum_{i=1}^N\log f(Y_i,\hat\alpha,r_0)-\lambda\norm{r_0}_{w,1}$. So $$\lambda\left(\norm{\hat r^{PI}}_{w,1}-\norm{r_0}_{w,1}\right)\leq\frac{1}{N}\sum_{i=1}^N\log f(Y_i,\hat\alpha,\hat r^{PI})-\log f(Y_i,\hat\alpha,r_0).$$ 
Moreover, $\displaystyle\log f(Y_i,\hat\alpha,\hat{r}^{PI})-\log f(Y_i,\hat\alpha, r_0)=\log\left(1+\frac{W(\hat\alpha,Y_i)^\prime(\hat{r}^{PI}-r_0)}{1+W(\hat\alpha,Y_i)^\prime r_0}\right)$.
Additionally, note that given any $\bar{y}>-1$, let $c'=\frac{1}{2(1+\max\{\bar{y},0\})}$, then $\log(1+y)\leq y-c'y^2$ holds for $y\in(-1,\bar{y}]$. Hence we have $\log f(Y_i,\hat\alpha,\hat{r}^{PI})-\log f(Y_i,\hat\alpha, r_0)\leq\frac{W(\hat\alpha,Y_i)^\prime(\hat{r}^{PI}-r_0)}{1+W(\hat\alpha,Y_i)^\prime r_0}-c_0\left\{\frac{W(\hat\alpha,Y_i)^\prime(\hat{r}^{PI}-r_0)}{1+W(\hat\alpha,Y_i)^\prime r_0}\right\}^2$, where $c_0>0$ is a universal constant.
Thus $\displaystyle\frac{c_0}{N}\sum_{i=1}^N\left\{\frac{W(\hat\alpha,Y_i)^\prime(\hat{r}^{PI}-r_0)}{1+W(\hat\alpha,Y_i)^\prime r_0}\right\}^2\leq\frac{1}{N}\sum_{i=1}^N\frac{W(\hat\alpha,Y_i)^\prime(\hat{r}^{PI}-r_0)}{1+W(\hat\alpha,Y_i)^\prime r_0}+\lambda\left(\norm{r_0}_{w,1}-\norm{\hat r^{PI}}_{w,1}\right)$. Since $\norm{{\rm diag}^{-1}(w) \frac{1}{N}\sum_{i=1}^N\frac{W(\hat\alpha,Y_i)}{1+W(\hat\alpha,Y_i)^\prime r_0}}_\infty\leq\frac{\lambda}{c}$, by H\"older's inequality, we have 
$$\begin{array}{rl}
&\quad\displaystyle\frac{c_0}{N}\sum_{i=1}^N\frac{\left[W(\hat\alpha,Y_i)^\prime(\hat{r}^{PI}-r_0)\right]^2}{(1+W(\hat\alpha,Y_i)^\prime r_0)^2}\\
&\displaystyle\leq\norm{\frac{{\rm diag}^{-1}(w)}{N}\sum_{i=1}^N\frac{W(\hat\alpha,Y_i)}{1+W(\hat\alpha,Y_i)^\prime r_0}}_\infty\norm{\hat{r}^{PI}-r_0}_{w,1}+\lambda\norm{r_0}_{w,1}-\lambda\norm{\hat{r}^{PI}}_{w,1}\\
&\displaystyle\leq\lambda(1+1/c)\norm{\hat{r}^{PI}-r_0}_{w,1}\leq\lambda(1+1/c)(\norm{\hat{r}_T^{PI}-r_0}_{w,1}+\norm{\hat{r}_{T^c}^{PI}}_{w,1})\\
&\displaystyle\leq_{(a)}\lambda(1+1/c)\norm{\hat{r}_T^{PI}-r_0}_{w,1}\left(1+\frac{c+1}{c-1}\right)=2\lambda\frac{c+1}{c-1}\norm{\hat{r}_T^{PI}-r_0}_{w,1},
\end{array}$$ 
where inequality (a) follows from Lemma~\ref{lem:diff-restricted}. Hence the result follows.
\end{proof}
\medskip
\begin{example}\label{ex:plug-in}
Let $M^6=\mathrm{o}(N)$, $\alpha_{0j}=\frac{1}{2}, j\in[M]$. Here we use $r_{0\ell}$ to denote the coordinate of $r_0$ corresponding to assortment $\ell$. For a fixed constant $a>0$, let $r_{0\ell}= a\,M/\sqrt N$ when $|\ell|=2$, and $r_{0\ell}= 0$ when $|\ell|\ge 3$. So $s=\binom{M}{2}$. So $\mathbb P(Y=y)
=2^{-M}[1+\sum_{1\le j<k\le M}r_{0,jk}(2y_j-1)(2y_k-1)],\forall y\in\{0,1\}^M$. Particularly, the probability is well defined since $1+W(\alpha_0,Y_i)'r_0=(1+a\frac{M}{\sqrt{N}}\sum_{j<k}(2Y_{ij}-1)(2Y_{ik}-1))=1+\mathrm{o}(1)$. Also $\sup_y|W(\alpha_0,y)'r_0|\le\frac{aM}{\sqrt{N}}\frac{M^2}{2}=\frac{aM^3}{2\sqrt N}=\mathrm{o}(1)$. Proposition~\ref{prop:concentration:no covariate} implies that $\|\hat{\alpha}-\alpha_0\|\asymp\sqrt{M/N}$. Proposition~\ref{prop:plug-in:no-covariate} requires $\lambda=\bar{C}M/\sqrt{N}$ for some absolute constant $\bar{C}$. By smoothness conditions, $\bigl\|\mathbb E[W(\hat\alpha,Y)]-\mathbb E[W(\alpha_0,Y)]\bigr\|_\infty\lesssim \sqrt{\frac{M}{N}}$. Additionally, $\|\mathbb E[W(\alpha_0,Y)]\|_\infty=\max_{\ell}|r_{0\ell}|=
a\,\frac{M}{\sqrt N}$. Hence $\|\mathbb E[W(\hat\alpha,Y)]\|_\infty
\le
a\,\frac{M}{\sqrt N}+C\sqrt{\frac{M}{N}}
\le \frac{\lambda}{4}$ for $\bar{C}$ sufficiently large. Let $\hat Q(r)=\frac1N\sum_{i=1}^N \log\bigl(1+W(\hat\alpha,Y_i)'r\bigr)-\lambda\|r\|_1$. Since the coordinates of $W(\alpha_0,Y_i)$ are uniformly bounded,
Hoeffding's inequality and a union bound over $p\le 2^M$ coordinates yield $\left\|
\frac{1}{N}\sum_{i=1}^N W(\alpha_0,Y_i)-\mathbb{E}[W(\alpha_0,Y_i)]\right\|_\infty=\mathrm{O}_p\!(\sqrt{M/N})=\mathrm{o}_p\!(M/\sqrt{N})$. Then the smoothness conditions imposed imply $\left\|
\frac{1}{N}\sum_{i=1}^N W(\hat\alpha,Y_i)-\mathbb{E}[W(\hat\alpha,Y)]\right\|_\infty=\mathrm{O}_p\!(\sqrt{M/N})=\mathrm{o}_p\!(M/\sqrt{N})$. Therefore, with probability tending to one, $\left\|\frac{1}{N}\sum_{i=1}^N W(\hat\alpha,Y_i)\right\|_\infty\le\frac{\lambda}{2}$. On this event, for any feasible $r\neq 0$, $\hat Q(r)-\hat Q(0)\le\left(\frac1N\sum_{i=1}^N W(\hat\alpha,Y_i)\right)'r-\lambda\|r\|_1\le
\left\|\frac{1}{N}\sum_{i=1}^N W(\hat\alpha,Y_i)
\right\|_\infty \|r\|_1-\lambda\|r\|_1\le-\frac{\lambda}{2}\|r\|_1<0$. Hence $\hat r^{PI}=0$ with probability tending to one. Further, on this event, $\sqrt{\frac{1}{N}\sum_{i=1}^N\frac{\{W(\alpha_0,Y_i)'(\hat{r}^{PI}-r_0)\}^2}{(1+W(\alpha_0,Y_i)'r_0)^2}}=\sqrt{\frac{1}{N}\sum_{i=1}^N\frac{\{W(\alpha_0,Y_i)'r_0\}^2}{(1+W(\alpha_0,Y_i)'r_0)^2}}$. Note that $r_0$'s support corresponds to all the assortments $|\ell|=2$. Let $\Sigma=\mathbb{E}\left[\widetilde{W}_i\widetilde{W}_i'\right]\in\mathbb{R}^{s\times s}$, where $\widetilde{W}_i$ denotes the subvector of $W(\alpha_0,Y_i)$ restricted to coordinates corresponding to assortments with with $|\ell|=2$ only, and $s={M\choose 2}$. Index the $s$ assortments with $|\ell|=2$ by edges $e=(j,k)$ of the complete graph $K_M$, where $j,k\in[M]$ and $j\neq k$. Then $\Sigma=\mathbf{I}_s+\frac{aM}{\sqrt{N}}A$. Here $\mathbf{I}_s$ is the $s\times s$ identity matrix, and $A$ is the adjacency matrix for the line graph of $K_M$, where $A_{e,e'}=1$ iff the two edges $e,e'$ share a vertex. For $M$ large, $\lambda_{\min}(\Sigma)$(i.e. the minimum eigenvalue of $\Sigma$) satisfies $\lambda_{\min}(\Sigma)=1-\frac{2aM}{\sqrt{N}}\geq1/2$ since $M^6=\mathrm{o}(N)$. By applying a standard matrix Chernoff lower-tail inequality, with high probability we get $\lambda_{\min}\left(\frac{1}{N}\sum_{i=1}^N\tilde{W}_i\tilde{W}_i'\right)\geq c_0$ for some absolute constant $c_0>0$. Note $1+W(\alpha_0,Y_i)'r_0=1+\frac{aM}{\sqrt{N}}\sum_{j<k}(2Y_j-1)(2Y_k-1)\in[\frac{1}{2},\frac{3}{2}]$ for $M$ sufficiently large. Thus $\sqrt{\frac{1}{N}\sum_{i=1}^N\frac{\{W(\alpha_0,Y_i)'r_0\}^2}{(1+W(\alpha_0,Y_i)'r_0)^2}}\geq\frac{2}{3}\sqrt{r_{0T}'\left(\frac{1}{N}\sum_{i=1}^N\widetilde{W}_i\widetilde{W}_i'\right)r_{0T}}\geq\Omega_p\left(M\sqrt{s/N}\right)$. 
\end{example}

\section{Proofs of Results for Section \ref{Sec:Rates:covariates}}\label{sec:proof:cov}
\subsection{Proof of Theorem \ref{thm:main:covariate}}
It is sufficient to verify each condition of Lemma \ref{lemma:main:FO-covariate} to prove Theorem~\ref{thm:main:covariate}. 
Firstly, the conditions (i)-(iv) hold as follows:
\begin{enumerate}
    \item Condition (i): By Lemma \ref{lemma:MinMax:Claim1:covariate}, with probability $1-\delta$ where $\delta\geq\max\{\exp(-\bar{C}_0\{\sqrt{\log M}\sqrt{Nh^d}+Nh^d\}),\log(M)/(\delta_NNh^d),1/2^M\}$ for some universal constant $C$,
$$\begin{array}{rl}
&\quad\displaystyle \min_{\alpha \in \mathcal{A}_x(\hat\alpha)} Q^{FO}(\alpha,r_*)-Q^{FO}(\alpha_{*x},r_*)\\
& \displaystyle 
 \geq-\max_{\alpha \in \mathcal{A}_x(\hat\alpha), \|X_i-x\|\leq h}\|\alpha(X_i)-\alpha_0(X_i)\|_1 C\sqrt{\frac{\log (M/\delta)}{Nh^d}}\\
&\quad \displaystyle {-} C\max_{\|X_i-x\|\leq h,\alpha \in \mathcal{A}_x(\hat\alpha)}\!\!\|\alpha(X_i)\! - \!\alpha_0(X_i)\|^2 \!\!-Csdh^2\left(\frac{Md}{\sqrt{Nh^{d}}}+Mdh^2\right).
\end{array}$$
Thus with 
 $\displaystyle\mathcal{R}_x(\hat{\alpha})=\max_{\|X_i-x\|\leq h,\alpha \in \mathcal{A}_x(\hat\alpha)}\|\alpha(X_i)-\alpha_0(X_i)\|_1 C\sqrt{\frac{\log(M/\delta)}{Nh^d}}+C\max_{\|X_i-x\|\leq h,\alpha \in \mathcal{A}_x(\hat\alpha)}\|\alpha(X_i) - \alpha_0(X_i)\|^2+Csdh^2\left(\frac{Md}{\sqrt{Nh^{d}}}+Mdh^2\right)$, the first term in (i) of Lemma \ref{lemma:main:FO-covariate} holds. Denote $\widetilde{W}_i:=\frac{W(\alpha_0(X_i),Y_i)}{1+W(\alpha_0(X_i),Y_i)'r_0(X_i)}$ and $T_i:=\mathrm{supp}(r_0(X_i)-r_{*x}(X_i))$, where $r_{*x}(X_i):=r_0(x)+\nabla r_0(x)(X_i-x)$. By Assumption \ref{ass:iid-min-den} $|T_{i}|\leq 2s$. Let $r_{0j}(\cdot)$ denote the $j$-th element of $r_0(\cdot)$ and $r_{*xj}(\cdot)$ denote the $j$-th element of $r_{*x}(\cdot)$. Since $r_{*xj}(X_i)-r_{0j}(X_i)=r_{0j}(x)+\nabla_x r_{0j}(x)(X_i-x) - r_{0j}(X_i)$, we have 
\begin{equation}\label{eq:r_0 estimate:ell_1 error}
\left|r_{*xj}(X_i)-r_{0j}(X_i)\right|\leq\bar{K}_{x,2}dh^2\ \mbox{for any $j\in T_i$},
\end{equation}
where we use the definition of $\bar{K}_{x,2}$ in the (component wise) Taylor expansion, and $\bar{K}_{x,2}$ is independent of $s,h,N,M,d$. Then with probability $1-\delta_{\alpha}$, $$\begin{array}{rl}
&\displaystyle\quad\frac{1}{N}\sum_{i=1}^N\!K_h(X_i-x)\left\{\frac{W(\alpha_0(X_i),Y_i)'(r_0(X_i)-r_{*x}(X_i))}{1+ W(\alpha_0(X_i),Y_i)'r_0(X_i)}\right\}^2\\
&\displaystyle\leq_{(1)}C's^2\{\max_{i\in[N], j\in T_i}(r_{0j}(X_i)-r_{*xj}(X_i))^2\}\leq_{(2)}Cs^2d^2h^4,
\end{array}$$ 
where (1) holds with probability $1-\delta_{\alpha}$ by Assumption \ref{ass:smoothness-assumptions} and sparsity condition in Assumption \ref{ass:iid-min-den}, (2) holds by (\ref{eq:r_0 estimate:ell_1 error}), and $C,C'$ are universal constants. Thus the second term in (i) of Lemma \ref{lemma:main:FO-covariate} holds with $\mathcal{G}(x)=Cs^2d^2h^4$. 
\item Condition (ii): By Lemma \ref{lemma:choice of lambda:covariate}, the choice of $\lambda_x$ in Theorem \ref{thm:main:covariate} makes (ii) of Lemma \ref{lemma:main:FO-covariate} holds. 
\item Condition (iii) of Lemma \ref{lemma:main:FO-covariate} follows immediately from Lemma \ref{lemma:covariate:score*approx}, which holds with probability $1-\delta-\frac{C(\bar{K}_{\max}/\bar{K}_2)^2}{Nh^d}$ for some universal constant $C$. 
\item Condition (iv): note that by Taylor expansion, for each $j\in [p], i\in[n]$ there is $\tilde \alpha^{ji}$ in the line segment between $\hat \alpha(X_i)$ and $\alpha_0(X_i)$ such that
$$ \begin{array}{rl}
W^{FO}_j(\alpha_{*x}(X_i),Y_i) &= W_j(\hat \alpha(X_i),Y_i) + (\alpha_{*x}(X_i)-\hat\alpha(X_i))'\nabla_\alpha W_j(\hat\alpha(X_i),Y_i)\\ 
& = W_j(\alpha_{0}(X_i),Y_i)+(\hat\alpha(X_i)-\alpha_{0}(X_i))'\nabla_\alpha W_j(\alpha_{0}(X_i),Y_i)\\
&\quad\displaystyle+\frac{1}{2}(\hat\alpha(X_i)-\alpha_{0}(X_i))'\nabla_\alpha^2W_j(\tilde \alpha^{ji},Y_i)(\hat\alpha(X_i)-\alpha_{0}(X_i))\\
&\quad + (\alpha_{*x}(X_i)-\hat\alpha(X_i))'\nabla_\alpha W_j(\hat\alpha(X_i),Y_i)\\ 
& = W_j(\alpha_{0}(X_i),Y_i)+(\alpha_{*x}(X_i)-\alpha_{0}(X_i))'\nabla_\alpha W_j(\alpha_{0}(X_i),Y_i) + R_{ij}^{num}(x,\hat\alpha),
\end{array}$$
where $R_{ij}^{num}(x,\hat\alpha)$ above contains second order terms. Following similar proof steps as above, for $R_i^{dem}(x,\hat \alpha)$ containing second order terms, we have  $$\frac{1+W(\alpha_0(X_i),Y_i)'r_0(X_i)}{1+W^{FO}(\alpha_{*x}(X_i),Y_i)'r_0(X_i)}=1 - \frac{(\alpha_{*x}(X_i)-\alpha_0(X_i))'\nabla_\alpha W(\alpha_0(X_i),Y_i)'r_0(X_i)+ R_i^{dem}(x,\hat \alpha)}{1+W^{FO}(\alpha_{*x}(X_i),Y_i)'r_0(X_i)}.$$ 
Particularly, Taylor expansion implies that for some $\tilde \alpha(X_i)$, $\bar \alpha(X_i)$, 
{\small$$ \begin{array}{rl}
W(\hat \alpha(X_i),Y_i)'r_0(X_i)&= W(\alpha_0(X_i),Y_i)'r_0(X_i) + (\hat \alpha(X_i)-\alpha_0(X_i))'\nabla_\alpha W(\alpha_0(X_i),Y_i)'r_0(X_i)\\
&\quad\displaystyle + \frac{1}{2}(\hat \alpha(X_i)-\alpha_0(X_i))'\nabla_\alpha^2 \{W(\tilde \alpha(X_i),Y_i)'r_0(X_i)\}(\hat \alpha(X_i)-\alpha_0(X_i)),
\end{array}$$}
{\small$$\begin{array}{rl}
&\quad(\alpha_{*x}(X_i)-\hat\alpha(X_i))'\nabla_\alpha \{W(\hat\alpha(X_i),Y_i)'r_0(X_i)\}\\  
&= (\alpha_{*x}(X_i)-\hat\alpha(X_i))'\left[\nabla_\alpha \{W(\alpha_0(X_i),Y_i)'r_0(X_i)\}+ \nabla_\alpha^2 \{W(\bar{\alpha}(X_i),Y_i)'r_0(X_i)\}(\hat\alpha(X_i)-\alpha_0(X_i))\right].
\end{array}$$}
Combining all these arguments above, we have $\frac{W^{FO}_j(\alpha_{*x}(X_i),Y_i) }{1+W^{FO}(\alpha_{*x}(X_i),Y_i)'r_0(X_i)} =  \frac{W_j(\alpha_0(X_i),Y_i)}{1+W(\alpha_0(X_i),Y_i)'r_0(X_i)} + \widetilde R_{ij}(x,\hat\alpha)$, $|\widetilde R_{ij}(x,\hat\alpha)|\leq C'(\|\alpha_{*x}(X_i)-\alpha_0(X_i)\|+\|\hat{\alpha}(X_i)-\alpha_0(X_i)\|^2)$ for some universal constant $C'$. Thus (iv) of Lemma \ref{lemma:main:FO-covariate} holds with $\widetilde{R}_i=C'(\|\alpha_{*x}(X_i)-\alpha_0(X_i)\|+\|\hat{\alpha}(X_i)-\alpha_0(X_i)\|^2)$. 
\end{enumerate}
Hence from the argument above and statement of Theorem \ref{thm:main:covariate}, conditions (i)-(iv) of Lemma \ref{lemma:main:FO-covariate} hold with probability $1-\delta$. Next, the remaining conditions of Lemma \ref{lemma:main:FO-covariate} hold with probability $1-\delta_{\alpha}$ by Assumption \ref{ass:smoothness-assumptions}. Thus Theorem \ref{thm:main:covariate} holds. \proofend
 
\subsection{Proof of Corollary \ref{cor:FO:covariate}}
Recall  $\displaystyle\mathcal{A}_x(\hat\alpha) =\prod_{j=1}^M \left\{ \alpha_j+\beta_j(X-x) : \begin{array}{l}\alpha_j \in [\hat \alpha_j(x) - cv_x \hat s_{j,0}(x), \hat\alpha_j(x) + cv_x \hat s_{j,0}(x)], \\ 
\beta_{jk} \in [\tilde \beta_{jk}(x) - cv_x \hat s_{j,k}(x), \tilde \beta_{jk}(x) + cv_x \hat s_{j,k}(x)], k \in [d] 
\end{array}\right\}$,
where $cv_x := (1-\delta_\alpha)\mbox{-quantile of} \max_{ j \in  [M], \ell \in\{0\}\cup[d]}\left|\sum_{i=1}^N \xi_i \frac{\hat \epsilon_{ij} \hat{\zeta}_{i\ell}}{\hat s_{j,\ell}(x)}\right|$
$\mbox{conditional on }\{Y_i,X_i\}_{i\in[N]}$, $\{\xi_i\}_{i\in[N]}$ are i.i.d. $\mathcal{N}(0,1)$ random variables. Here $\hat\epsilon_{ij}:=Y_{ij} - \hat{\alpha}_j(x)-\tilde\beta_j(x)(X_i-x)$, and $\hat{\zeta}_{i} := (Z'W_xZ)^{-1}Z_i K_h(X_i-x)$  with $Z_i=(1, (X_i-x)')'$ and $W_{x}$ is a diagonal $N\times N$ matrix with $W_{x,ii}:=K_h(X_i-x)$, $\hat s_{j,\ell}(x) := \sqrt{\sum_{i=1}^N \hat \epsilon_{ij}^2 \hat{\zeta}_{i\ell}^2}$, where $\hat{\zeta}_{i\ell}$ is the $\ell$-th entry of $\hat{\zeta}_i$, $\forall \ell\in\{0\}\cup[d]$.  Note that by triangular inequality we have 
$$\begin{array}{rl}
&\quad\displaystyle\mathbb{P}\left(\max_{\alpha \in \mathcal{A}_x(\hat\alpha), \|X_i-x\|\leq h}\norm{\alpha(X_i)-\alpha_0(X_i)}\geq t\right)\\
&\displaystyle\leq\mathbb{P}\left(\max_{\alpha \in \mathcal{A}_x(\hat\alpha), \|X_i-x\|\leq h}\norm{\hat{\alpha}(x)+\tilde{\beta}(x)(X_i-x)-\alpha_0(X_i)}\geq\frac{t}{2}\right)\\
&\displaystyle\ +\mathbb{P}\left(cv_{x}\sqrt{\sum_{j\in[M]}\left(\sum_{\ell\in[d]}\hat{s}_{j,\ell}(x)|X_{il}-x_{l}|+\hat{s}_{j,0}(x)\right)^2}\geq\frac{t}{2}\right)\\
&\displaystyle\leq_{(a)}\mathbb{P}\left(\max_{\alpha \in \mathcal{A}_x(\hat\alpha), \|X_i-x\|\leq h}\norm{\hat{\alpha}(x)+\tilde{\beta}(x)(X_i-x)-\alpha_0(X_i)}\geq\frac{t}{2}\right)\\
&\displaystyle\ +\mathbb{P}\left(cv_{x}\sqrt{2dh^2\sum_{j=1}^M\max_{\ell\in[d]}\hat{s}_{j,\ell}(x)^2+2\sum_{j=1}^M\hat{s}_{j,0}(x)^2}\geq\frac{t}{2}\right),
\end{array}$$
where the first inequality follows from the triangular inequality, (a) above follows because $\sum_{\ell\in[d]}(X_{i\ell}-x_{\ell})^2\leq h^2$, $\sum_{j\in[M]}\bigg(\sum_{\ell\in[d]}\hat{s}_{j,\ell}(x)(X_{il}-x_{l})+\hat{s}_{j,0}(x)\bigg)^2\leq 2dh^2\sum_{j=1}^M\max_{\ell\in[d]}\hat{s}_{j,\ell}(x)^2+2\sum_{j=1}^M\hat{s}_{j,0}(x)^2$. Note that 
$$\begin{array}{rl}
&\displaystyle\quad\mathbb{P}\left(\max_{\alpha \in \mathcal{A}_x(\hat\alpha), \|X_i-x\|\leq h}\norm{\hat{\alpha}(x)+\tilde{\beta}(x)(X_i-x)-\alpha_0(X_i)}\geq\frac{t}{2}\right)\\
&\displaystyle\leq_{(1)}\mathbb{P}\left(\max_{\alpha \in \mathcal{A}_x(\hat\alpha), \|X_i-x\|\leq h}\norm{(\hat{\alpha}(x)-\alpha_0(x))+(\tilde{\beta}(x)-\nabla\alpha_0(x))(X_i-x)}+\sqrt{M}\bar{K}_{x,2}\|X_i-x\|^2\geq\frac{t}{2}\right),
\end{array}$$
where (1) holds by Taylor expansion and Assumption~\ref{ass:smoothness-assumptions}. Furthermore, Proposition \ref{prop:concentration:covariate} implies that both $\norm{\hat{\alpha}(x)-\alpha_0(x)}_2\leq\tilde{C}_0(x)\left(\sqrt{\frac{M}{Nh^d}}+\sqrt{M}dh^2+\sqrt{\left(\frac{1}{Nh^d}+h^4d^2\right)\log(1/\delta)}\right)$ and $\norm{\tilde{\beta}(x)-\nabla\alpha_0(x)}_{2,2}\leq\displaystyle\tilde{C}_1(x)\sqrt{d}\bigg(\sqrt{\frac{M}{Nh^{d+2}}}+\sqrt{M}dh^2+\sqrt{\left(\frac{1}{Nh^{d+2}}+h^4d^4\right)\log(1/\delta)}\bigg)$ hold with probability $1-\delta$ for some absolute constants $\tilde{C}_0(x)$ and $\tilde{C}_1(x)$. Thus for $X_i$ such that $\|X_i-x\|\leq h$, with probability $1-\delta$, for some absolute constant $\tilde{C}(x)$, 
$$\begin{array}{rl}
&\displaystyle\quad\norm{(\hat{\alpha}(x)-\alpha_0(x))+(\tilde{\beta}(x)-\nabla\alpha_0(x))(X_i-x)}\\
&\displaystyle\leq\tilde{C}(x)\left[\sqrt{\frac{M}{Nh^d}}+\sqrt{M}dh^2+\sqrt{\left(\frac{1}{Nh^d}+h^4d^2\right)\log(1/\delta)}\right]\\
&\displaystyle\quad+\tilde{C}(x)\left[h\sqrt{d}\left(\sqrt{\frac{M}{Nh^{d+2}}}+\sqrt{M}dh^2+\sqrt{\left(\frac{1}{Nh^{d+2}}+h^4d^4\right)\log(1/\delta)}\right)\right].
\end{array}$$
Additionally, $\sum_{i=1}^N\frac{\hat \epsilon_{ij}^2\hat{\zeta}_{i\ell}^2}{\hat s_{j,\ell}(x)^2}=1$ for any $j\in[M], \ell\in[d]\cup\{0\}$, so 
conditional on $\{Y_i,X_i\}$, $\sum_{i=1}^N \xi_i \frac{\hat \epsilon_{ij} \hat{\zeta}_{i\ell}}{\hat s_{j,\ell}(x)}\sim \mathcal{N}\left(0,1\right)$, thus 
$\mathbb{E}\left[\max_{ j \in  [M], \ell \in\{0\}\cup[d]}\sum_{i=1}^N \xi_i \frac{\hat \epsilon_{ij} \hat{\zeta}_{i\ell}}{\hat s_{j,\ell}(x)}\right]\leq\sqrt{2\log\{M(d+1)\}}$.
By Lemma~\ref{lemma:borell} (Borell-TIS inequality \citep{borell1975brunn}), with probability $1-\delta$, we have 
\begin{equation}\label{eq:cov:concentration:normal}
    \max_{ j \in  [M], \ell \in [d]\cup\{0\}} \left|\sum_{i=1}^N\frac{ \xi_i \hat \epsilon_{ij} \hat{\zeta}_{i\ell}}{\hat s_{j,\ell}(x)}\right|\leq2\sqrt{\log\left(\frac{M(d+1)}{\delta}\right)}.
\end{equation}
Let $\delta\leq\delta_{\alpha}$ and recall that $\delta\geq\exp(-C\{\sqrt{\log M}\sqrt{Nh^d}+Nh^d\})$, and set $t=C\sqrt{\frac{Md\log\{M(d+1)/\delta\}}{Nh^d}}+C\sqrt{M}dh^2+C\sqrt{\frac{\log(1/\delta)}{Nh^d}}$ for some absolute constant $C$ in the above, then the above results imply that $\max_{\alpha \in \mathcal{A}_x(\hat\alpha), \|X_i-x\|\leq h}\norm{\alpha(X_i)-\alpha_0(X_i)}\leq C\sqrt{Md}\left[\sqrt{\frac{\log\{M(d+1)/\delta\}}{Nh^d}}+\sqrt{d}h^2\right]+C\sqrt{\frac{\log(1/\delta)}{Nh^d}}$ holds with probability $1-2\delta$. Additionally, by triangular inequality, for any $t'>0$ we have 
$$\begin{array}{rl}
&\displaystyle\quad\mathbb{P}\left(\max_{\alpha \in \mathcal{A}_x(\hat\alpha), \|X_i-x\|\leq h}\norm{\alpha(X_i)-\alpha_0(X_i)}_1\geq t'\right)\\
&\displaystyle\leq\mathbb{P}\left(\max_{\alpha \in \mathcal{A}_x(\hat\alpha), \|X_i-x\|\leq h}\norm{\hat{\alpha}(x)+\tilde{\beta}(x)(X_i-x)-\alpha_0(X_i)}_1\!\!\geq\!\frac{t'}{2}\right)\\
&\displaystyle\quad+\mathbb{P}\left(\!cv_{x}\!\!\left[\sum_{j\in[M],\ell\in[d]}\!\!\!\!\!\!\!\!\!\!\!|\hat{s}_{j,\ell}(x)(X_{i\ell}-x_{\ell})|+\!\!\!\!\!\sum_{j\in[M]}\!\!\!\!|\hat{s}_{j,0}(x)|\right]\geq\frac{t'}{2}\right)\\
&\displaystyle\leq_{(b)}\mathbb{P}\left(\max_{\alpha \in \mathcal{A}_x(\hat\alpha), \|X_i-x\|\leq h}\norm{\hat{\alpha}(x)+\tilde{\beta}(x)(X_i-x)-\alpha_0(X_i)}_1\geq\frac{t'}{2}\right)\\
&\displaystyle\quad\quad+\mathbb{P}\left(cv_{x}\left[\sqrt{Md}h\sqrt{\sum_{j=1}^M\max_{\ell\in[d]}\hat{s}_{j,\ell}^2(x)}+\sqrt{M}\sqrt{\sum_{j\in[M]}\hat{s}_{j,0}^2(x)}\right]\geq\frac{t'}{2}\right),
\end{array}$$
where (b) follows since
$$\sum_{j\in[M],\ell\in[d]}|\hat{s}_{j,\ell}(x)(X_{il}-x_{l})|+\sum_{j\in[M]}|\hat{s}_{j,0}(x)|\leq\sqrt{Md}h\sqrt{\sum_{j=1}^M\max_{\ell\in[d]}\hat{s}_{j,\ell}^2(x)}+\sqrt{M}\sqrt{\sum_{j\in[M]}\hat{s}_{j,0}^2(x)}.$$ 
Note 
$$\begin{array}{rl}
&\displaystyle\quad\mathbb{P}\left(\max_{\alpha \in \mathcal{A}_x(\hat\alpha), \|X_i-x\|\leq h}\norm{\hat{\alpha}(x)+\tilde{\beta}(x)(X_i-x)-\alpha_0(X_i)}_1\geq\frac{t'}{2}\right)\\
&\displaystyle\leq_{(1)}\mathbb{P}\left(\max_{\alpha \in \mathcal{A}_x(\hat\alpha), \|X_i-x\|\leq h}\norm{(\hat{\alpha}(x)-\alpha_0(x))+(\tilde{\beta}(x)-\nabla\alpha_0(x))(X_i-x)}_1+M\bar{K}_{x,2}h^2\geq\frac{t'}{2}\right),
\end{array}$$
where (1) holds by Taylor expansion. By Proposition \ref{prop:concentration:covariate}, with probability $1-\delta$, 
$$\norm{\hat{\alpha}(x)-\alpha_0(x)}_1\leq\tilde{C}_0(x)\left(\frac{M}{\sqrt{Nh^d}}+Mdh^2+\sqrt{M\left(\frac{1}{Nh^d}+h^4d^2\right)\log(1/\delta)}\right),$$ $$\norm{\tilde{\beta}(x)-\nabla\alpha_0(x)}_{1,1}\leq\tilde{C}_1(x)d\left(\frac{M}{\sqrt{Nh^{d+2}}}+Md^2h^2+\sqrt{M\left(\frac{1}{Nh^{d+2}}+h^4d^4\right)\log(1/\delta)}\right).$$
Hence for some absolute constant $\tilde{C}'(x)$, for $X_i$ such that $\|X_i-x\|\leq h$, with probability $1-\delta$ we have 
$$\begin{array}{rl}
&\displaystyle\quad\norm{(\hat{\alpha}(x)-\alpha_0(x))+(\tilde{\beta}(x)-\nabla\alpha_0(x))(X_i-x)}_1\\
&\displaystyle\leq\tilde{C}'(x)\left[\frac{M}{\sqrt{Nh^d}}+Mdh^2+\sqrt{M\left(\frac{1}{Nh^d}+h^4d^2\right)\log(1/\delta)}\right]\\
&\displaystyle\quad+\tilde{C}'(x)\left[hd\left(\frac{M}{\sqrt{Nh^{d+2}}}+Md^2h^2+\sqrt{M\left(\frac{1}{Nh^{d+2}}+h^4d^4\right)\log(1/\delta)}\right)\right].
\end{array}$$
Let $\delta\leq\delta_{\alpha}$, set $t'=C'\frac{Md}{\sqrt{Nh^d}}+C'Mdh^2+C'd\sqrt{\frac{M\log(1/\delta)}{Nh^d}}$ in the above, where $C'$ is an absolute constant. Note that Assumption~\ref{ass:moment-assumptions} implies $dh^2<1$ and note that $\delta\geq2^{-M}$, then the previous results imply that with probability $1-2\delta$, $\max_{\alpha \in \mathcal{A}_x(\hat\alpha), \|X_i-x\|\leq h}\norm{\alpha(X_i)-\alpha_0(X_i)}_1\leq C'M\left(\frac{d}{\sqrt{Nh^d}}+dh^2\right)+C'd\sqrt{\frac{M\log(1/\delta)}{Nh^d}}$. Note that $\lambda_x=C\sqrt{\log(p/\delta)/Nh^d}+Ch^2$ under the given condition, Theorem~\ref{thm:main:covariate} implies that with probability $1-C\delta-C\delta_\alpha-\frac{C(\bar{K}_{\max}/\bar{K}_2)^2}{Nh^d}-C\epsilon$,
$$\begin{array}{rl}
\widetilde{\mathcal{R}}(\hat\alpha) &\displaystyle\leq C_1\max_{\alpha \in \mathcal{A}_x(\hat\alpha), \|X_i-x\|\leq h}\|\alpha(X_i)-\alpha_0(X_i)\|_1 \sqrt{\frac{\log(M/\delta)}{Nh^d}}+C_1sdh^2\sqrt{\frac{\log(1/\delta)}{Nh^d}}+C_1s^2d^2h^4\\
&\displaystyle\quad + C_1\max_{\alpha \in \mathcal{A}_x(\hat\alpha), \|X_i-x\|\leq h} \|\alpha(X_i)-\alpha_0(X_i)\|^2+C_1sdh^2\left(\frac{Md}{\sqrt{Nh^{d}}}+Mdh^2\right)\\
&\displaystyle\leq_{(1)}C_2(M+s^2)d^2h^4,
\end{array}$$
where $C,C_1,C_2$ are constants independent of $s,N,M,h,d$ above. Here (1) above uses Assumption~\ref{ass:moment-assumptions}, and the high-probability bounds derived for ${\displaystyle\max_{\alpha \in \mathcal{A}_x(\hat\alpha), \|X_i-x\|\leq h}}\|\alpha(X_i)-\alpha_0(X_i)\|_1$ and ${\displaystyle\max_{\alpha \in \mathcal{A}_x(\hat\alpha), \|X_i-x\|\leq h}} \|\alpha(X_i)-\alpha_0(X_i)\|^2$. Thus, Theorem~\ref{thm:main:covariate} implies $\sqrt{\frac{1}{N}\sum_{i=1}^N K_h(X_i-x)\Delta_{x,i}^{FO}}\leq C\sqrt{\frac{sM+s\log(1/\delta)}{Nh^d}}+C(\sqrt{M}+s)dh^2$. 

Furthermore, Assumption \ref{ass:moment-assumptions} requires $h$ to satisfy 
$\left(\frac{M^3}{N\delta_N^2}\right)^{1/d}\leq h\leq\sqrt{\frac{\delta_N\xi}{sd}}$. By first order condition, we have $h^*=\min\left\{\sqrt{\frac{\delta_N\xi}{sd}},\max\left\{\left(\frac{M^3}{N\delta_N^2}\right)^{1/d},4^{-\frac{2}{d+4}}\left(\frac{sM+s\log(1/\delta)}{N(s+\sqrt{M})^2}\right)^{\frac{1}{d+4}}\right\}\right\}$. Particularly, when $h=h^*$, by direct calculation, we have $\sqrt{\frac{sM+s\log(1/\delta)}{Nh^d}}+(\sqrt{M}+s)dh^2\leq\bar{C}\min\left\{B_{1},B_2\right\}$ for some universal constant $\bar{C}$, where $B_1:=\delta_N\sqrt{\frac{s(M+\log(1/\delta))}{M^3}}+(\sqrt{M}+s)d\left(\frac{M^3}{N\delta_N^2}\right)^{\frac{2}{d}}$,
$B_2:=\sqrt{\frac{s(M+\log(1/\delta))}{N}}\left(\frac{sd}{\delta_N\xi}\right)^{d/4}\!\!+\frac{(\sqrt{M}+s)\delta_N\xi}{s}$.
\proofend

\subsection{Technical Lemmas for Section \ref{Sec:Rates:covariates}}
\begin{lemma}\label{lemma:main:FO-covariate}
For any $x\in \mathcal{X}$, define $\alpha_{*x}(X):= \alpha_0(x)+\nabla_x\alpha_0(x)(X-x)$ and define $r_{*x}(X):= r_0(x)+\nabla_xr_0(x)(X-x)$. Moreover for any $\bar \alpha \in [0,1]^M, \bar y \in \{0,1\}^M$, let 
$$W^{FO}_i(\bar \alpha,\bar{y})=W(\hat \alpha(X_i),\bar y)+\nabla_\alpha W(\hat\alpha(X_i),\bar y)(\bar \alpha-\hat\alpha(X_i)).$$ 
Suppose\\
(i) $\min_{\alpha\in\mathcal{A}_x(\hat \alpha)}Q^{FO}(\alpha,r_{*x})-Q^{FO}(\alpha_{*x},r_{*x}) \geq -\mathcal{R}_x(\hat{\alpha})$, \\
and $\displaystyle\frac{1}{N}\sum_{i=1}^NK_h(X_i-x) \left\{\frac{W(\alpha_0(X_i),Y_i)'(r_0(X_i)-r_{*x}(X_i))}{1+ W(\alpha_0(X_i),Y_i)'r_0(X_i)}\right\}^2 \leq \mathcal{G}(x)$,\\
(ii)$\displaystyle\frac{\lambda_x}{c} \geq \left\|  \mathrm{diag}^{-1}(w_x)\frac{1}{N}\sum_{i=1}^N\frac{  K_h(X_i-x) }{1+W^{FO}(\alpha_{*x}(X_i),Y_i)^\prime r_0(X_i)} \begin{pmatrix} W^{FO}(\alpha_{*x}(X_i),Y_i)\\{\rm vec}( W^{FO}(\alpha_{*x}(X_i),Y_i)\frac{(X_i-x)'}{h} )
\end{pmatrix}\right\|_\infty $\\
(iii) $\displaystyle\left|\frac{1}{N}\sum_{i=1}^N K_h(X_i-x)\frac{ W(\alpha_0(X_i),Y_i)^\prime(r_{*x}(X_i)-r_0(X_i))}{1+W(\alpha_0(X_i),Y_i)^\prime r_0(X_i)}\right|\leq Csdh^2\sqrt{\frac{\log(1/\delta)}{Nh^d}}$\\ 
(iv) $\displaystyle\left\|\frac{ W^{FO}(\alpha_{*x}(X_i),Y_i)}{1+W^{FO}(\alpha_{*x}(X_i),Y_i)^\prime r_0(X_i)}-\frac{ W(\alpha_0(X_i),Y_i)}{1+W(\alpha_0(X_i),Y_i)^\prime r_0(X_i)}\right\|_\infty \leq \widetilde{R}_{i}$. (v) $\alpha_{*x} \in \mathcal{A}_x(\hat\alpha)$. \\
(vi) $\displaystyle \max_{\|X_i-x\|\leq h, j\in [p], \tilde \alpha \in \mathcal{A}_x(\hat\alpha)}|v'\nabla_\alpha W_j(\tilde \alpha(X_i),Y_i)|  +|v'\nabla_\alpha W(\tilde \alpha(X_i),Y_i)'r_0(X_i)| \leq M^{1/2}K_{\alpha,1}\|v\|$, \\
(vii) $\displaystyle\max_{\|X_i-x\|\leq h, j \in [p], \tilde \alpha \in \mathcal{A}_x(\hat\alpha)}|u'\nabla_\alpha^2 W(\tilde\alpha(X_i),Y_i)'r_0(X_i)v|+|u'\nabla_\alpha^2 W_j(\tilde\alpha(X_i),Y_i)v|\leq K_{\alpha,2}\|u\|\|v\|$, and\\
${\displaystyle\max_{\tilde{x}\in\mathcal{X},i\in[N]}}|u'\nabla_{\alpha}W(\alpha_0(\tilde{x}),Y_i)'r_0(X_i)|\leq\bar{K}_w\|u\|_1$, ${\displaystyle\max_{\tilde{x}\in\mathcal{X},i\in[N]}}|u'\nabla_{\alpha}W(\alpha_0(\tilde{x}),Y_i)'v|\leq \bar{K}_w\norm{u}_{\infty}\norm{v}_{1}$. (viii) ${\displaystyle\max_{\|\tilde{x}-x\|\leq h, \alpha \in \mathcal{A}_x(\hat\alpha)}}\|\alpha(x)-\alpha_0(x)\|M^{1/2}K_{\alpha,1}/\xi < 1/4$, ${\displaystyle\max_{\|\tilde{x}-x\|\leq h, \tilde\alpha \in \mathcal{A}_x(\hat\alpha)}}\|\tilde \alpha(x)-\alpha_0(x)\|^2 K_{\alpha,2}/\xi \leq 1/16$. 
(ix) ${\displaystyle\max_{j\in [p], k\in[M], x\in \mathcal{X}}}|v'\nabla_{x}^2 r_{0j}(x)v|+ |v'\nabla^2\alpha_{0k}(x)v|\leq \bar{K}_{x,2}\|v\|^2$. \\
(x) $\delta_N\geq s{\displaystyle\max_{\alpha \in \mathcal{A}_x(\hat\alpha), \|X_i-x\|\leq h}}\|\tilde\alpha(X_i)-\alpha_0(X_i)\|^2/ \{\xi\kappa_{x,\bar c}^2{\displaystyle\min_{j\in[p], k\in[d]}}\{w_{xjk}^2,w_{xj}^2\}\}$. \\
$\delta_N\geq{\displaystyle\max_{\alpha \in \mathcal{A}_x(\hat\alpha), \|X_i-x\|\leq h}}\|\tilde\alpha(X_i)-\alpha_0(X_i)\|/ \{\xi{\displaystyle\min_{j\in[p], k\in[d]}}\{w_{xjk},w_{xj}\}\}$. 

Then for all $x \in \mathcal{X}$, for some universal constant $C$, with probability $1-C\epsilon$, 
$$\sqrt{\frac{1}{N} \sum_{i=1}^NK_h(X_i-x)\left\{\frac{W(\alpha_0(X_i),Y_i)^\prime(\hat{r}^{FO}(X_i,x)-r_{*x}(X_i))}{1+W(\alpha_0(X_i),Y_i)^\prime r_0(X_i)}\right\}^2}\leq C\left(\lambda_x \left( 1 + \frac{1}{c}\right) \frac{\sqrt{s}}{\kappa_{x,\bar{c}}} + \widetilde{\mathcal{R}}(x)^{1/2}\right),$$
where $\displaystyle\widetilde{\mathcal{R}}(x)\leq \mathcal{R}_x(\hat\alpha) + \mathcal{G}(x) + Csdh^2\left(\!\!\sqrt{\frac{\log(1/\delta)}{Nh^d}}\! + \! \frac{1}{N}\sum_{i=1}^NK_h(X_i-x)|\widetilde{R}_i|\!+\!\frac{1}{N}\sum_{i=1}^NK_h(X_i-x)|\widetilde{R}_i|^2\right)$.
\end{lemma}
\begin{proof}[Proof of Lemma \ref{lemma:main:FO-covariate}]
Let $Q^{FO}(\alpha,r_{a,b}):= \frac{1}{N}\sum_{i=1}^N K_h(X_i-x)\log(1+ W^{FO}(\alpha(X_i),Y_i)'r_{a,b}(X_i,x))$. 
Following similar proof steps as in Proposition~\ref{prop:plug-in:no-covariate} and Theorem~\ref{thm:main:FO}, using Assumption~\ref{ass:iid-min-den} and the given conditions, with probability $1-C\epsilon$, we have 
$$1 +  W^{FO}(\alpha_{*x}(X_i),Y_i)'\hat r_{\hat a,\hat b}(X_i))>c_1\xi, 1 + W^{FO}(\alpha_{*x}(X_i),Y_i)' r_0(X_i)>c_1\xi$$ 
for some $c_1\in(0,1)$. Additionally, note that given any $\bar{y}>-1$, there exists a constant $c'=\frac{1}{2(1+\max\{\bar{y},0\})}$ such that $\log(1+y)\leq y-c'y^2$ holds for $y\in(-1,\bar{y}]$. Hence for some universal constant $c_0>0$ we have 
$$\begin{array}{rl}
       &\displaystyle\quad \log ( 1 +  W^{FO}(\alpha_{*x}(X_i),Y_i)'\hat r_{\hat a,\hat b}(X_i))   -\log( 1 + W^{FO}(\alpha_{*x}(X_i),Y_i)' r_0(X_i)) \\
       &\displaystyle=\log\left(1+\frac{ W^{FO}(\alpha_{*x}(X_i),Y_i)^\prime(\hat{r}_{\hat a,\hat b}(X_i,x)-r_0(X_i))}{1+ W^{FO}(\alpha_{*x}(X_i),Y_i)^\prime r_0(X_i)}\right)\\
       &\displaystyle\leq\frac{ W^{FO}(\alpha_{*x}(X_i),Y_i)^\prime(\hat{r}_{\hat a,\hat b}(X_i,x)-r_0(X_i))}{1+W^{FO}(\alpha_{*x}(X_i),Y_i)^\prime r_0(X_i)}-c_0\left\{\frac{W^{FO}(\alpha_{*x}(X_i),Y_i)^\prime(\hat{r}_{\hat a,\hat b}(X_i,x)-r_0(X_i))}{1+W^{FO}(\alpha_{*x}(X_i),Y_i)^\prime r_0(X_i)}\right\}^2,
\end{array}$$
Define $\Delta_{0,ab}^{FO}:=\frac{W^{FO}(\alpha_{*x}(X_i),Y_i)^\prime(\hat{r}_{\hat a, \hat b}(X_i,x)-r_0(X_i))}{1+W^{FO}(\alpha_{*x}(X_i),Y_i)^\prime r_0(X_i)}$, then 
\begin{equation}\label{eq:covariate:quadratic:Q}
\begin{array}{rl} 
&\quad \displaystyle Q^{FO}(\alpha_{*x},\hat r_{\hat a,\hat b})-Q^{FO}(\alpha_{*x}, r_0)  \\
&\displaystyle \leq \frac{1}{N}\sum_{i=1}^N K_h(X_i-x)\left\{\Delta_{0,ab}^{FO}-c_0(\Delta_{0,ab}^{FO})^2\right\}.
\end{array}
\end{equation}
$$\begin{array}{rl}
\mbox{Next note that}&\displaystyle\quad \frac{1}{N}\sum_{i=1}^N \frac{K_h(X_i-x)W^{FO}(\alpha_{*x}(X_i),Y_i)^\prime(\hat{r}_{\hat a,\hat b}(X_i,x)-r_0(X_i))}{1+W^{FO}(\alpha_{*x}(X_i),Y_i)^\prime r_0(X_i)} \\
&\displaystyle = \underbrace{\frac{1}{N}\sum_{i=1}^N\frac{K_h(X_i-x)W^{FO}(\alpha_{*x}(X_i),Y_i)^\prime( \hat a - r_0(x) + \hat b(X_i-x) - \nabla_x r_0(x)(X_i-x))}{1+W^{FO}(\alpha_{*x}(X_i),Y_i)^\prime r_0(X_i)}}_{\textrm{(I)}} \\
&\displaystyle\quad + \underbrace{\frac{1}{N}\sum_{i=1}^N\frac{K_h(X_i-x)W^{FO}(\alpha_{*x}(X_i),Y_i)^\prime(r_{*x}(X_i)-r_0(X_i))}{1+W^{FO}(\alpha_{*x}(X_i),Y_i)^\prime r_0(X_i)}}_{\textrm{(II)}} = \textrm{(I)} + \textrm{(II)}.\ \ \mbox{Note that}
\end{array}$$
$$\begin{array}{rl}
|\textrm{(I)}|
&\displaystyle\leq\left\|\mathrm{diag}^{-1}(w_x)\frac{1}{N}\sum_{i=1}^N\frac{  K_h(X_i-x) }{1+W^{FO}(\alpha_{*x}(X_i),Y_i)^\prime r_0(X_i)} \left( { W^{FO}(\alpha_{*x}(X_i),Y_i)}\atop{ {\rm vec}(W^{FO}(\alpha_{*x}(X_i),Y_i)\frac{(X_i-x)'}{h} )}\right)\right\|_\infty\\
&\displaystyle\quad\quad\times\left\|\left( {\hat a - r_0(x)}\atop{ {\rm vec}\big(h(\hat b - \nabla_xr_0(x)})\big)\right) \right\|_{w_x,1}\\
&\displaystyle\leq\frac{\lambda_x}{c}\left\|\left( {\hat a - r_0(x)}\atop{ {\rm vec}\big(h(\hat b - \nabla_xr_0(x)})\big)\right) \right\|_{w_x,1},
\end{array}$$
where in the above, the first inequality holds by H\"older's inequality and the second inequality holds by condition (ii). Since $r_{*x}(X_i)-r_0(X_i)=r_0(x)+\nabla_x r_0(x)(X_i-x) - r_0(X_i)$, by Taylor expansion, $\left\|  r_{*x}(X_i)-r_0(X_i) \right\|_1=\| r_0(x)+\nabla_x r_0(x)(X_i-x) - r_0(X_i) \|_1\leq\bar{K}sdh^2$ for some universal constant $\bar{K}$ by condition (ix). Define $\Delta^{FO}(X_i,Y_i):=\left\{\frac{W^{FO}(\alpha_{*x}(X_i),Y_i)^\prime}{1+W^{FO}(\alpha_{*x}(X_i),Y_i)^\prime r_0(X_i)}-\frac{W(\alpha_0(X_i),Y_i)^\prime}{1+W(\alpha_0(X_i),Y_i)^\prime r_0(X_i)}\right\}$.
So for some universal constant $C$,
$$\begin{array}{rl}
|\textrm{(II)}|&\displaystyle \leq_{(1)} \left|\frac{1}{N}\sum_{i=1}^N K_h(X_i-x)\frac{ W(\alpha_0(X_i),Y_i)^\prime(r_{*x}(X_i)-r_0(X_i))}{1+W(\alpha_0(X_i),Y_i)^\prime r_0(X_i)}\right|\\
&\displaystyle\quad+\bigg|\frac{1}{N}\sum_{i=1}^N K_h(X_i-x)\Delta^{FO}(X_i,Y_i)(r_{*x}(X_i)-r_0(X_i)) \bigg|\\
&\displaystyle\leq_{(2)} Csdh^2\sqrt{\frac{\log(1/\delta)}{Nh^d}}+ \frac{1}{N}{\displaystyle\sum_{i=1}^N} K_h(X_i-x)\left\|\Delta^{FO}(X_i,Y_i)\right\|_\infty\|r_{*x}(X_i)-r_0(X_i)\|_1 \\
&\displaystyle \leq_{(3)} Csdh^2\sqrt{\frac{\log(1/\delta)}{Nh^d}} +  \frac{Csdh^2}{N}\sum_{i=1}^N K_h(X_i-x)\widetilde{R}_i,
\end{array}$$
where (1) holds by the triangle inequality, the first term of (2) holds by condition (iii), and the second term of (2) holds by H\"older's inequality, (3) holds by condition (iv) and the bound for $\left\|  r_{*x}(X_i)-r_0(X_i) \right\|_1$ above. So the arguments above imply
\begin{equation}\label{eq:bound:I:II}
\begin{array}{rl}
    &\quad\displaystyle\frac{1}{N}\sum_{i=1}^N K_h(X_i-x)\Delta_{0,ab}^{FO}\\
    &\displaystyle\leq Csdh^2\left[\sqrt{\frac{\log(1/\delta)}{Nh^d}} +  \frac{1}{N}\sum_{i=1}^N K_h(X_i-x)\widetilde{R}_i\right]\\
    &\displaystyle\quad+\frac{\lambda_x}{c}\left\|\left( {\hat a - r_0(x)}\atop{ {\rm vec}\big(h(\hat b - \nabla_xr_0(x)})\big)\right) \right\|_{w_x,1}.
\end{array}
\end{equation}
Let $a^*:=r_0(x), b^*:=\nabla r_0(x)$. So $r_{*x}(X_i)=r_{a^*,b^*}(X_i,x)$, and 
$$\left\|\left( {\hat a - r_0(x)}\atop{ {\rm vec}\big(h(\hat b - \nabla_xr_0(x)})\big)\right) \right\|_{w_x,1}=\big\|\big(\hat a - a^*,{\rm vec}\big(h(\hat b - b^*)\big)\big)\big\|_{w_x,1}=\|r_{\hat{a},\hat{b}}-r_{a^*,b^*}\|_{w_x,h,1}.$$
\begin{equation}\label{eq:covariate:main:steps}
\begin{array}{rl}
&\quad\displaystyle Q^{FO}(\alpha_{*x},\hat r_{\hat a,\hat b}) - \lambda_x \|\hat r_{\hat a,\hat b}\|_{w_x,h,1}\\
& \geq_{(1)}\displaystyle \min_{\alpha \in \mathcal{A}_x(\hat\alpha)} Q^{FO}(\alpha,\hat r_{\hat a,\hat b}) - \lambda_x \|\hat r_{\hat a,\hat b}\|_{w_x,h,1} \\
& \displaystyle \geq_{(2)} \min_{\alpha \in \mathcal{A}_x(\hat\alpha)} Q^{FO}(\alpha,r_{*x}) - \lambda_x \|r_{a^*,b^*}\|_{w_x,h,1}\\
&\displaystyle \geq \min_{\alpha \in \mathcal{A}_x(\hat\alpha)} \{Q^{FO}(\alpha,r_{*x})-Q^{FO}(\alpha_{*x},r_{*x})\} \\
&\quad + Q^{FO}(\alpha_{*x},r_{*x})- \lambda_x \|r_{a^*,b^*}\|_{w_x,h,1} \\
&\displaystyle  \geq_{(3)}  - \mathcal{R}_x(\hat\alpha) + Q^{FO}(\alpha_{*x},r_{*x})- \lambda_x \|r_{a^*,b^*}\|_{w_x,h,1} \\
\end{array}
\end{equation}
where (1) holds from condition (v) $\alpha_{*x} \in \mathcal{A}_x(\hat\alpha)$, (2) holds by the optimality of $\hat a,\hat b$, (3) holds since $\min_{\alpha \in \mathcal{A}_x(\hat\alpha)} \{Q^{FO}(\alpha,r_{*x})-Q^{FO}(\alpha_{*x},r_{*x})\} \geq - \mathcal{R}_x(\hat\alpha)$ by condition (i). Hence,
$$\begin{array}{rl}
&\displaystyle\frac{c_0}{N}\sum_{i=1}^NK_h(X_i-x)(\Delta_{0,ab}^{FO})^2\leq_{(1)}\frac{1}{N}\sum_{i=1}^NK_h(X_i-x)(\Delta_{0,ab}^{FO})-\{Q^{FO}(\alpha_{*x},\hat r_{\hat a,\hat b})-Q^{FO}(\alpha_{*x}, r_0)\}\\
&\displaystyle=\frac{1}{N}\sum_{i=1}^NK_h(X_i-x)(\Delta_{0,ab}^{FO})-\{Q^{FO}(\alpha_{*x},\hat r_{\hat a,\hat b})-Q^{FO}(\alpha_{*x}, r_{*x})+Q^{FO}(\alpha_{*x},r_{*x})-Q^{FO}(\alpha_{*x}, r_0)\}\\
&\displaystyle\leq_{(2)}2Csdh^2\left[\sqrt{\frac{\log(1/\delta)}{Nh^d}} +  \frac{1}{N}\sum_{i=1}^N K_h(X_i-x)\widetilde{R}_i\right]+\frac{\lambda_x}{c}\|\hat r_{\hat a, \hat b} - r_{a^*,b^*}\|_{w_x,h,1}\\
&\quad\quad\displaystyle+\lambda_x\|r_{a^*,b^*}\|_{w_x,h,1}-\lambda_x\|\hat r_{\hat a,\hat b}\|_{w_x,h,1} + \mathcal{R}_x(\hat\alpha),\ \mbox{where (1) is due to }\eqref{eq:covariate:quadratic:Q}, \mbox{(2) from \eqref{eq:bound:I:II}, \eqref{eq:covariate:main:steps}}
\end{array}$$
and following from similar proof steps for \eqref{eq:covariate:quadratic:Q}, $-(Q^{FO}(\alpha_{*x},r_{*x})-Q^{FO}(\alpha_{*x}, r_0))\leq|\textrm{(II)}|$. The above inequality is equivalent to $\frac{c_0}{N}\sum_{i=1}^NK_h(X_i-x)\left\{\frac{W^{FO}(\alpha_{*x}(X_i),Y_i)^\prime(\hat{r}_{\hat a,\hat b}(X_i,x)-r_0(X_i))}{1+W^{FO}(\alpha_{*x}(X_i),Y_i)^\prime r_0(X_i)}\right\}^2\leq\lambda_x\|r_{a^*,b^*}\|_{w_x,h,1} -  \lambda_x\|\hat r_{\hat a,\hat b}\|_{w_x,h,1} + \mathcal{R}_x(\hat\alpha) + \frac{\lambda_x}{c}\|\hat r_{\hat a, \hat b} - r_{a^*,b^*}\|_{w_x,h,1}+2Csdh^2\left(\sqrt{\frac{\log(1/\delta)}{Nh^d}}+\frac{1}{N}\sum_{i=1}^N K_h(X_i-x)\widetilde{R}_i\right)$. Define $\Delta_{*x,ab}^{FO}:=\frac{W^{FO}(\alpha_{*x}(X_i),Y_i)^\prime(\hat{r}_{\hat a, \hat b}(X_i,x)-r_{*x}(X_i))}{1+W^{FO}(\alpha_{*x}(X_i),Y_i)^\prime r_0(X_i)}$, $\Delta_0W_i:=\frac{W(\alpha_0(X_i),Y_i)^\prime(r_0(X_i)-r_{*x}(X_i))}{1+W(\alpha_0(X_i),Y_i)^\prime r_0(X_i)}$.
Then for some universal constant $C'$,
$$\begin{array}{rl}
&\displaystyle \sqrt{\frac{c_0}{N}\sum_{i=1}^NK_h(X_i-x)\left\{\Delta_{0,ab}^{FO}\right\}^2}-\sqrt{\frac{c_0}{N} \sum_{i=1}^NK_h(X_i-x)\left\{\Delta_{*x,ab}^{FO}\right\}^2}\\
&\displaystyle \geq - \sqrt{\frac{c_0}{N} \sum_{i=1}^NK_h(X_i-x)\left\{\Delta^{FO}(X_i,Y_i)(r_0(X_i)-r_{*x}(X_i))\right\}^2} - \sqrt{\frac{c_0}{N} \sum_{i=1}^NK_h(X_i-x)\left\{\Delta_0W_i\right\}^2}   \\
&\displaystyle \geq - \sqrt{\frac{c_0}{N} \sum_{i=1}^NK_h(X_i-x)\left\|\Delta^{FO}(X_i,Y_i)\right\|_\infty^2\|r_0(X_i)-r_{*x}(X_i)\|^2_1}-\sqrt{\frac{c_0}{N} \sum_{i=1}^NK_h(X_i-x)\left\{\Delta_0W_i\right\}^2}   \\
&\displaystyle \geq -C'sdh^2\sqrt{\frac{c_0}{N} \sum_{i=1}^NK_h(X_i-x) \widetilde{R}_i^2} - \sqrt{\frac{c_0}{N} \sum_{i=1}^NK_h(X_i-x)\left\{\Delta_0W_i\right\}^2},   
\end{array}$$ 
where the last inequality follows from condition (iv) and the bound for $\left\|  r_{*x}(X_i)-r_0(X_i) \right\|_1$ above. Using the fact that that for non-negative numbers $a, b, c$ we have $a+c\geq b\Rightarrow 2a^2+2c^2\geq b^2$, 
$$\begin{array}{rl}
&\displaystyle\frac{2c_0}{N}\!\!\sum_{i=1}^N\!K_h(X_i-x)\!\!\left\{\!\frac{W^{FO}(\alpha_{*x}(X_i),Y_i)^\prime(\hat{r}_{\hat a, \hat b}(X_i,x)-r_0(X_i))}{1+W^{FO}(\alpha_{*x}(X_i),Y_i)^\prime r_0(X_i)}\right\}^2 \!\!\!\!\!+\!\! 4\!\left\{\!C'sdh^2\sqrt{\frac{c_0}{N} \sum_{i=1}^NK_h(X_i-x) \widetilde{R}_i^2}\right\}^2\\
&\displaystyle + \frac{4c_0}{N} \sum_{i=1}^NK_h(X_i-x)\left\{\frac{W(\alpha_0(X_i),Y_i)^\prime(r_0(X_i)-r_{*x}(X_i))}{1+W(\alpha_0(X_i),Y_i)^\prime r_0(X_i)}\right\}^2 \\
&\displaystyle \geq\frac{c_0}{N} \sum_{i=1}^NK_h(X_i-x)\left\{\frac{W^{FO}(\alpha_{*x}(X_i),Y_i)^\prime(\hat{r}_{\hat a, \hat b}(X_i,x)-r_{*x}(X_i))}{1+W^{FO}(\alpha_{*x}(X_i),Y_i)^\prime r_0(X_i)}\right\}^2.\ \mbox{Next, define }
\end{array}$$
$$\begin{array}{rl}
\widetilde{\mathcal{R}}(x,\hat{\alpha})&\displaystyle:= \mathcal{R}_x(\hat\alpha)  + 2Csdh^2\sqrt{\frac{\log(1/\delta)}{Nh^d}}+2(C')^2s^2d^2h^4\frac{c_0}{N} \sum_{i=1}^NK_h(X_i-x)\widetilde{R}_i^2\\
&\displaystyle \ +  \frac{2Csdh^2}{N}\sum_{i=1}^N K_h(X_i-x)\widetilde{R}_i+\frac{2c_0}{N} \sum_{i=1}^NK_h(X_i-x)\left\{\frac{W(\alpha_0(X_i),Y_i)^\prime(r_0(X_i)-r_{*x}(X_i))}{1+W(\alpha_0(X_i),Y_i)^\prime r_0(X_i)}\right\}^2.
\end{array}$$
The upper bound for $\frac{c_0}{N}\sum_{i=1}^NK_h(X_i-x)\left\{\frac{W^{FO}(\alpha_{*x}(X_i),Y_i)^\prime(\hat{r}_{\hat a,\hat b}(X_i,x)-r_0(X_i))}{1+W^{FO}(\alpha_{*x}(X_i),Y_i)^\prime r_0(X_i)}\right\}^2$ then further implies that $\frac{c_0}{2N} \sum_{i=1}^NK_h(X_i-x)\left\{\frac{W^{FO}(\alpha_{*x}(X_i),Y_i)^\prime(\hat{r}_{\hat a, \hat b}(X_i,x)-r_{*x}(X_i))}{1+W^{FO}(\alpha_{*x}(X_i),Y_i)^\prime r_0(X_i)}\right\}^2\leq \lambda_x\|r_{a^*,b^*}\|_{w_x,h,1} -  \lambda_x\|\hat r_{\hat a,\hat b}\|_{w_x,h,1} + \frac{\lambda_x}{c}\|\hat r_{\hat a,\hat b} - r_{a^*,b^*}\|_{w_x,h,1}   + \widetilde{\mathcal{R}}(x,\hat{\alpha})$. Define $\displaystyle D_{i*}^{\mathrm{FO}}:=\left\{\frac{W^{FO}(\alpha_{*x}(X_i),Y_i)^\prime(\hat{r}_{\hat a, \hat b}(X_i,x)-r_{*x}(X_i))}{1+W^{FO}(\alpha_{*x}(X_i),Y_i)^\prime r_0(X_i)}\right\}^2$. Let $T_x$ denote the support of $r_{a^*,b^*}$. Using the fact that $\|r_{a^*,b^*}\|_{w_x,h,1} - \|(\hat r_{\hat a,\hat b})_{T_x}\|_{w_x,h,1}\leq\|r_{a^*,b^*}-(\hat r_{\hat a,\hat b})_{T_x}\|_{w_x,h,1}$, 
\begin{equation}\label{ineq:*:ineq}
\begin{array}{rl}
&\displaystyle\quad\frac{c_0}{2N} \sum_{i=1}^NK_h(X_i-x)D_{i*}^{\mathrm{FO}}-\widetilde{\mathcal{R}}(x,\hat{\alpha})\\ 
&\displaystyle \leq \lambda_x(1+1/c)\|(\hat r_{\hat a,\hat b})_{T_x} - r_{a^*,b^*}\|_{w_x,h,1}\\
&\quad\displaystyle-\lambda_x(1-1/c)\|(\hat r_{\hat a,\hat b})_{T_x^c}\|_{w_x,h,1}
\end{array}
\end{equation}
Let $\mathcal{W}_x:={\rm diag}(w_x)$. Define $U_x,\hat{U}_x$ as follows: 
$$\begin{array}{rcl}
U_x&:=&\displaystyle\frac{c_0}{2N}\sum_{i=1}^N\displaystyle\frac{K_h(X_i-x)\left( {W(\alpha_0(X_i),Y_i)}\atop{ {\rm vec}( W(\alpha_0(X_i),Y_i) (X_i-x)'/h)}\right)\left( {W(\alpha_0(X_i),Y_i)}\atop{ {\rm vec}( W(\alpha_0(X_i),Y_i) (X_i-x)'/h)}\right)'}{(1+W(\alpha_0(X_i),Y_i)^\prime r_0(X_i))^2},\\
\hat{U}_x&:=&\displaystyle\frac{c_0}{2N}\sum_{i=1}^N\frac{K_h(X_i-x)\left( {W^{FO}(\alpha_{*x}(X_i),Y_i)}\atop{ {\rm vec}( W^{FO}(\alpha_{*x}(X_i),Y_i) (X_i-x)'/h)}\right)\left( {W^{FO}(\alpha_{*x}(X_i),Y_i)}\atop{ {\rm vec}( W^{FO}(\alpha_{*x}(X_i),Y_i) (X_i-x)'/h)}\right)'}{(1+W^{FO}(\alpha_{*x}(X_i),Y_i)^\prime r_0(X_i))^2}.
\end{array}$$ 
Let $\hat v_x = [ (\hat a)', h{\rm vec}(\hat b)' ]' - [ (a^{*})', h{\rm vec}(b^*)']'$. By H\"older's inequality, 
\begin{equation}\label{eq:holder}
\begin{array}{rl}
&\quad\displaystyle\hat v_x'\hat U_x \hat v_x  - \hat v_x' U_x \hat v_x\\
&\displaystyle\geq \! - \|\hat r_{\hat a,\hat b} - r_{a^*,b^*}\|_{w_x,h,1}^2 \! \|\mathcal{W}_x^{-1}\!(U_x-\hat U_x)\!\mathcal{W}_x^{-1}\|_{\infty,\infty}
\end{array}
\end{equation}
We consider two exclusive cases in the following: $\|(\hat r_{\hat a,\hat b})_{T_x^c}\|_{w_x,h,1} \geq 2\frac{c+1}{c-1}\|(\hat r_{\hat a,\hat b})_{T_x} - r_{a^*,b^*}\|_{w_x,h,1}$ and $\|(\hat r_{\hat a,\hat b})_{T_x^c}\|_{w_x,h,1} < 2\frac{c+1}{c-1}\|(\hat r_{\hat a,\hat b})_{T_x} - r_{a^*,b^*}\|_{w_x,h,1}$. 

Firstly suppose \textbf{Case 1}: $\|(\hat r_{\hat a,\hat b})_{T_x^c}\|_{w_x,h,1} \geq 2\frac{c+1}{c-1}\|(\hat r_{\hat a,\hat b})_{T_x} - r_{a^*,b^*}\|_{w_x,h,1}$.
then \eqref{ineq:*:ineq} implies 
$$\frac{c_0}{2N}\sum_{i=1}^NK_h(X_i-x)D_{i*}^{\mathrm{FO}} + \frac{1}{2}\lambda_x\left(1-\frac{1}{c}\right)\|(\hat r_{\hat a,\hat b})_{T^c_x}\|_{w_x,h,1}\leq\widetilde{\mathcal{R}}(x,\hat\alpha),$$ 
which immediately implies $\lambda_x\left(1-\frac{1}{c}\right)\|(\hat r_{\hat a,\hat b})_{T^c_x}\|_{w_x,h,1}\leq2\widetilde{\mathcal{R}}(x,\hat\alpha)$. Then 
$$\begin{array}{rl}
&\quad\displaystyle\hat v_x'\hat U_x \hat v_x  - \hat v_x' U_x \hat v_x \geq  - \|\hat r_{\hat a,\hat b} - r_{a^*,b^*}\|_{w_x,h,1}^2  \|\mathcal{W}_x^{-1}(U_x-\hat U_x)\mathcal{W}_x^{-1}\|_{\infty,\infty}\\
&\displaystyle\geq  - \left(1+\frac{(c-1)}{2(c+1)}\right)^2\| (\hat r_{\hat a, \hat b})_{T_x^c}\|_{w_x,h,1}^2 \|\mathcal{W}_x^{-1}(U_x-\hat U_x)\mathcal{W}_x^{-1}\|_{\infty,\infty}\\
&\displaystyle\geq-\left(1+\frac{(c-1)}{2(c+1)}\right)^2\ \frac{4\widetilde{\mathcal{R}}^2(x,\hat\alpha)}{(1-\frac{1}{c})^2\lambda^2_x}  \|\mathcal{W}_x^{-1}(U_x-\hat U_x)\mathcal{W}_x^{-1}\|_{\infty,\infty}\\
&\displaystyle\geq-\frac{9\widetilde{\mathcal{R}}^2(x,\hat\alpha)}{(1-1/c)^2\lambda_x^2}\|\mathcal{W}_x^{-1}(U_x-\hat U_x)\mathcal{W}_x^{-1}\|_{\infty,\infty}.
\end{array}$$
Rearranging the terms implies $\hat v_x' U_x \hat v_x \leq  \frac{9\widetilde{\mathcal{R}}^2(x,\hat\alpha)}{(1-1/c)^2\lambda_x^2} \|\mathcal{W}_x^{-1}(\hat U_x - U_x)\mathcal{W}_x^{-1} \|_{\infty,\infty} + \hat v_x'\hat U_x \hat v_x$, which implies $\hat v_x' U_x \hat v_x\leq\frac{9\widetilde{\mathcal{R}}^2(x,\hat\alpha)}{(1-1/c)^2\lambda_x^2} \|\mathcal{W}_x^{-1}(\hat U_x - U_x)\mathcal{W}_x^{-1} \|_{\infty,\infty} +\widetilde{\mathcal{R}}(x,\hat\alpha)$. This inequality holds since \eqref{ineq:*:ineq} implies that $\frac{c_0}{2N}\sum_{i=1}^NK_h(X_i-x)D_{i*}^{\mathrm{FO}} + \frac{1}{2}\lambda_x\left(1-\frac{1}{c}\right)\|(\hat r_{\hat a,\hat b})_{T^c_x}\|_{w_x,h,1}\leq\widetilde{\mathcal{R}}(x,\hat\alpha)$, which immediately implies $\hat v_x'\hat U_x \hat v_x=\frac{c_0}{2N}\sum_{i=1}^NK_h(X_i-x)D_{i*}^{\mathrm{FO}}\leq\widetilde{\mathcal{R}}(x,\hat\alpha)$. Thus in Case 1 we have $\sqrt{\hat v_x' U_x \hat v_x}\leq\frac{3\widetilde{\mathcal{R}}(x,\hat\alpha)}{(1-1/c)\lambda_x} \|\mathcal{W}_x^{-1}(\hat U_x - U_x)\mathcal{W}_x^{-1} \|_{\infty,\infty}^{1/2}+\widetilde{\mathcal{R}}(x,\hat\alpha)^{1/2}$. By Taylor expansion, $\|\mathcal{W}_x^{-1}(\hat U_x - U_x)\mathcal{W}_x^{-1} \|_{\infty,\infty}\leq C_1(\max_{i\in[N]}\|\hat{\alpha}(X_i)-\alpha_0(X_i)\|^2+h^2)$ for some universal constant $C_1$ from the conditions (vi) -- (viii). 
Condition (x) further implies that $\|\mathcal{W}_x^{-1}(\hat U_x - U_x)\mathcal{W}_x^{-1} \|_{\infty,\infty}^{1/2}$ is negligible for $N$ sufficiently large. Hence for some universal constant $C_2$, under case 1 we have 
\begin{equation}\label{eq:case1:covariate}
\sqrt{\hat v_x' U_x \hat v_x}\leq C_2\widetilde{\mathcal{R}}(x,\hat\alpha)^{1/2}.
\end{equation}
Next suppose \textbf{Case 2} holds: $\displaystyle\|(\hat r_{\hat a, \hat b})_{T_x^c}\|_{w_x,h,1} \leq 2\frac{c+1}{c-1}\|(\hat r_{\hat a,\hat b} - r_{a^*,b^*})_{T_x}\|_{w_x,h,1}$. Define $$\displaystyle\hat{\Delta}_{ab}^i:=\left\{\frac{W(\alpha_0(X_i),Y_i)^\prime(\hat{r}_{\hat a,\hat b}(X_i,x)-r_{*x}(X_i))}{1+W(\alpha_0(X_i),Y_i)^\prime r_0(X_i)}\right\}^2.$$ 
Then we have
\begin{equation}\label{eq:thm1:main:main:well:covariate}
\begin{array}{rl}
&\displaystyle\quad \hat v_x' U_x \hat v_x=\frac{c_0}{2N} \sum_{i=1}^NK_h(X_i-x)\hat{\Delta}_{ab}^i\\
&\displaystyle\leq_{(1)}\hat v_x' \hat{U}_x \hat v_x+\|\hat r_{\hat a,\hat b} - r_{a^*,b^*}\|_{w_x,h,1}^2\! \|\mathcal{W}_x^{-1}(U_x-\hat U_x)\mathcal{W}_x^{-1}\|_{\infty,\infty}\\
&\displaystyle \leq_{(2)} \lambda_x(1+\frac{1}{c})\|(\hat r_{\hat a,\hat b} - r_{a^*,b^*})_{T_x}\|_{w_x,h,1} + \|\hat r_{\hat a,\hat b} - r_{a^*,b^*}\|_{w_x,h,1}^2\! \|\mathcal{W}_x^{-1}(U_x-\hat U_x)\mathcal{W}_x^{-1}\|_{\infty,\infty}+ \widetilde{\mathcal{R}}(x,\hat\alpha)\\
&\displaystyle\leq_{(3)}\lambda_x(1+\frac{1}{c})\|(\hat r_{\hat a,\hat b} - r_{a^*,b^*})_{T_x}\|_{w_x,h,1}\\
&\displaystyle\quad+2\left[4(\frac{c+1}{c-1})^2+1\right]\!\!\|(\hat{r}_{\hat{a},\hat{b}}-r_{a^*,b^*})_{T_x}\|_{w_x,h,1}^2\|\mathcal{W}_x^{-1}(U_x-\hat U_x)\mathcal{W}_x^{-1}\|_{\infty,\infty}+ \widetilde{\mathcal{R}}(x,\hat\alpha)\\
&\displaystyle\leq_{(4)}\frac{\lambda_x}{\sqrt{c_0}}\left(1+\frac{1}{c}\right) \frac{\sqrt{s}}{\kappa_{x,\bar{c}}}\sqrt{\frac{c_0}{N} \sum_{i=1}^NK_h(X_i-x)\hat{\Delta}_{ab}^i} + \widetilde{\mathcal{R}}(x,\hat\alpha)\\
&\displaystyle\quad+\frac{s\eta_0\|\!\mathcal{W}_x^{-1}\!(U_x-\hat U_x)\!\mathcal{W}_x^{-1}\|_{\infty,\infty}}{\kappa_{x,\bar{c}}^2}\sum_{i=1}^N\frac{K_h(X_i-x)\hat{\Delta}_{ab}^i}{N}
\end{array}
\end{equation}
where (1) of \eqref{eq:thm1:main:main:well:covariate} follows from \eqref{eq:holder}, (2) of \eqref{eq:thm1:main:main:well:covariate} follows from \eqref{ineq:*:ineq} and that $\hat{v}_x\hat{U}_x\hat{v}_x=\frac{c_0}{2N} \sum_{i=1}^NK_h(X_i-x)D_{i*}^{\mathrm{FO}}$, (3) follows since 
$$\begin{array}{rl}
\|\hat{r}_{\hat{a},\hat{b}}-r_{a^*,b^*}\|_{w_x,h,1}^2&\displaystyle\leq2\|(\hat{r}_{\hat{a},\hat{b}}-r_{a^*,b^*})_{T_x}\|_{w_x,h,1}^2+2\|(\hat{r}_{\hat{a},\hat{b}})_{T_x^c}\|_{w_x,h,1}^2\\
&\displaystyle\leq2\left[4(\frac{c+1}{c-1})^2+1\right]\|(\hat{r}_{\hat{a},\hat{b}}-r_{a^*,b^*})_{T_x}\|_{w_x,h,1}^2,
\end{array}$$
(4) follows from the definition of restricted eigenvalue and the fact that 
$$\|(\hat{r}_{\hat{a},\hat{b}}-r_{a^*,b^*})_{T_x}\|_{w_x,h,1}\leq\sqrt{s}\|(\hat{r}_{\hat{a},\hat{b}}-r_{a^*,b^*})_{T_x}\|_{w_x,h,2},$$ 
with $\eta_0=2\left[4\left(\frac{c+1}{c-1}\right)^2+1\right]$.
Note that by condition (x), $$\frac{s\eta_0\|\mathcal{W}_x^{-1}(U_x-\hat U_x)\mathcal{W}_x^{-1}\|_{\infty,\infty}}{\kappa_{x,\bar{c}}^2} \leq c_0/4$$ 
when $N$ is sufficiently large. Further note that $\frac{1}{4}t^2 \leq \frac{\lambda_x}{\sqrt{c_0}}(1+\frac{1}{c})\frac{\sqrt{s}}{\kappa_{x,\bar{c}}} t + \mathcal{\widetilde{R}}(x,\hat \alpha) \Rightarrow t \leq\frac{4}{\sqrt{c_0}}\lambda_x(1+\frac{1}{c})\frac{\sqrt{s}}{\kappa_{x,\bar{c}}} + 2{\mathcal{\widetilde{R}}}^{1/2}(x,\hat \alpha)$. Let $t=\sqrt{2}\sqrt{\hat{v}_x'U_x\hat{v}_x}$, then we have $\sqrt{\hat{v}_x'U_x\hat{v}_x}\leq\frac{2\sqrt{2}}{\sqrt{c_0}}\lambda_x(1+\frac{1}{c})\frac{\sqrt{s}}{\kappa_{x,\bar{c}}} + \sqrt{2}{\mathcal{\widetilde{R}}}^{1/2}(x,\hat \alpha)$ in Case 2. Recall that \eqref{eq:case1:covariate} holds in Case 1. Thus the result holds from the definition of $\widetilde{\mathcal{R}}(x,\hat{\alpha})$ and the second statement of condition (i). 
\end{proof}

Recall from Section~\ref{Sec:Rates:covariates} that $\beta_{N,x}:={\displaystyle\max_{j\in[p],k\in[d]}}\left\{\mathcal{B}_{1,xj},\mathcal{B}_{2,xjk}\right\}$, where for any $j\in[p], k\in[d]$, $\mathcal{B}_{1,xj}:=\frac{1}{w_{xj}}\frac{1}{N}\sum_{i=1}^N\frac{K_h(X_i-x)|W_j(\alpha_0(X_i),Y_i)|}{1+W(\alpha_0(X_i),Y_i)^\prime r_0(X_i)}$, $\mathcal{B}_{2,xjk}:=\frac{1}{w_{xjk}}\frac{1}{N}\sum_{i=1}^N\frac{K_h(X_i-x)|W_j(\alpha_0(X_i),Y_i)(X_{ik}-x_k)/h|}{1+W(\alpha_0(X_i),Y_i)^\prime r_0(X_i)}$. 
\begin{lemma}\label{lemma:choice of lambda:covariate}
Under Assumptions \ref{ass:iid-min-den}, \ref{ass:moment-assumptions}, \ref{ass:smoothness-assumptions}, \ref{ass:local-linear}, 
$$\begin{array}{rl}
&\displaystyle\quad\left\|\mathrm{diag}^{-1}(w_x)\frac{1}{N}\sum_{i=1}^N\frac{  K_h(X_i-x) }{1+W(\hat{\alpha}(X_i),Y_i)^\prime r_0(X_i)} \left( { W(\hat{\alpha}(X_i),Y_i)}\atop{ {\rm vec}( W(\hat{\alpha}(X_i),Y_i)\frac{(X_i-x)'}{h} )}\right)^\prime \right\|_\infty \\
&\displaystyle\leq\left\|\mathrm{diag}^{-1}(w_x)\frac{1}{N}\sum_{i=1}^N\frac{  K_h(X_i-x) }{1+W(\alpha_0(X_i),Y_i)^\prime r_0(X_i)} \left( { W(\alpha_0(X_i),Y_i)}\atop{ {\rm vec}( W(\alpha_0(X_i),Y_i)\frac{(X_i-x)'}{h} )}\right)^\prime \right\|_\infty\\
&\displaystyle\quad+C(1+\beta_{N,x})\max_{\|X_i-x\|\leq h}\norm{\hat{\alpha}(X_i)-\alpha_0(X_i)}\frac{M^{1/2}K_{\alpha,1}}{\xi},
\end{array}$$
$$\begin{array}{rl}
&\displaystyle\left\|\mathrm{diag}^{-1}(w_x)\frac{1}{N}\sum_{i=1}^N\frac{  K_h(X_i-x) }{1+W^{FO}(\alpha_{*x}(X_i),Y_i)^\prime r_0(X_i)} \left( { W^{FO}(\alpha_{*x}(X_i),Y_i)}\atop{ {\rm vec}( W^{FO}(\alpha_{*x}(X_i),Y_i)\frac{(X_i-x)'}{h} )}\right)^\prime \right\|_\infty \\
&\displaystyle\leq\left\|\mathrm{diag}^{-1}(w_x)\frac{1}{N}\sum_{i=1}^N\frac{ K_h(X_i-x) }{1+W(\alpha_0(X_i),Y_i)^\prime r_0(X_i)} \left( { W(\alpha_{*x}(X_i),Y_i)}\atop{ {\rm vec}( W(\alpha_{*x}(X_i),Y_i)\frac{(X_i-x)'}{h} )}\right)^\prime \right\|_\infty\\
&\displaystyle\quad+C(1+\beta_{N,x})\frac{K_{\alpha,2}}{\xi}\max_{\alpha \in \mathcal{A}_x(\hat\alpha), \|X_i-x\|\leq h}\norm{\alpha(X_i)-\alpha_0(X_i)}^2+Ch^2.
\end{array}$$
Moreover, with probability $1-C\delta_{\alpha}-\delta(1+\mathrm{o}(1))$ for some absolute constant $C$,
$$\begin{array}{rl}
&\displaystyle\quad\left\|\mathrm{diag}^{-1}(w_x)\frac{1}{N}\sum_{i=1}^N\frac{  K_h(X_i-x) }{1+W(\alpha_0(X_i),Y_i)^\prime r_0(X_i)} \left( { W(\alpha_0(X_i),Y_i)}\atop{ {\rm vec}( W(\alpha_0(X_i),Y_i)\frac{(X_i-x)'}{h} )}\right)^\prime \right\|_\infty \\
&\displaystyle\leq\frac{\Phi^{-1}(1-\delta/2p(1+d))}{\sqrt{Nh^d}}\max_{j\in[p],k\in\{0\}\cup[d]}\frac{1}{w_{xjk}}\sqrt{\frac{1}{N}\sum_{i=1}^N\frac{K_h(X_i-x)W_j(\alpha_0(X_i),Y_i)^2}{(1+W(\alpha_0(X_i),Y_i)^\prime  r_0(X_i))^2}}.
\end{array}$$
\end{lemma}
\begin{proof}[Proof of Lemma \ref{lemma:choice of lambda:covariate}]
The proof is similar to the no-covariate case, i.e. Lemma \ref{lemma:Score:nocovariate}. 
By Taylor expansion, for each $j\in[p],i\in[N]$, there exists $\tilde{\alpha}^{ji}$, $\tilde{\alpha}$ in the line segment between $\hat{\alpha}(X_i)$ and $\alpha_0(X_i)$, such that 
$W_j(\hat{\alpha}(X_i),Y_i)=W_j(\alpha_0(X_i),Y_i)+(\hat{\alpha}(X_i)-\alpha_0(X_i))^\prime\nabla_{\alpha}W_j(\tilde{\alpha}^{ji},Y_i)$. Thus 
$$
\begin{array}{rl}
&\displaystyle\quad\frac{W_j(\hat\alpha(X_i),Y_i)}{1+W(\hat\alpha(X_i),Y_i)'r_0(X_i)}\\
&\displaystyle= \frac{W_j(\alpha_0(X_i),Y_i) +   (\hat\alpha(X_i)-\alpha_0(X_i))'\nabla_\alpha W_j(\tilde \alpha^{ji},Y_i)}{1+W(\alpha_0(X_i),Y_i)'r_0(X_i) + (\hat\alpha(X_i)-\alpha_0(X_i))'\nabla_\alpha W(\tilde{\alpha},Y_i)'r_0(X_i)} \\
&\displaystyle= \frac{W_j(\alpha_0(X_i),Y_i) }{1+W(\alpha_0(X_i),Y_i)'r_0(X_i)}\!\!\left(\!1\! - \!\frac{(\hat\alpha(X_i)-\alpha_0(X_i))'\nabla_\alpha W(\tilde \alpha,Y_i)'r_0(X_i)}{1+W(\alpha_0(X_i),Y_i)'r_0(X_i) + (\hat\alpha(X_i)-\alpha_0(X_i))'\nabla_\alpha W(\tilde{\alpha},Y_i)'r_0(X_i)}\!\right)\\
&\displaystyle\quad+\frac{(\hat\alpha(X_i)-\alpha_0(X_i))'\nabla_\alpha W_j(\tilde \alpha^{ji},Y_i)}{1+W(\alpha_0(X_i),Y_i)'r_0(X_i) + (\hat\alpha(X_i)-\alpha_0(X_i))'\nabla_\alpha W(\tilde{\alpha},Y_i)'r_0(X_i)}\\
&\displaystyle = \frac{W_j(\alpha_0(X_i),Y_i) }{1+W(\alpha_0(X_i),Y_i)'r_0(X_i)} + \widetilde R(\hat\alpha(X_i),Y_i),
\end{array}$$ 
where $\displaystyle|\widetilde R(\hat\alpha(X_i),Y_i)|\leq C\left(1+\frac{|W_j(\alpha_0(X_i),Y_i)|}{1+W(\alpha_0(X_i),Y_i)^\prime r_0(X_i)}\right)\frac{\norm{\hat{\alpha}(X_i)-\alpha_0(X_i)}M^{1/2}K_{\alpha,1}}{\xi}$ with probability $1-\delta_{\alpha}$ from Assumption~\ref{ass:smoothness-assumptions}. Note that $K_h(X_i-x)>0$ if and only if $\norm{(X_i-x)/h}\leq1$. So the first inequality holds. Now we prove the second inequality. By Taylor expansion, for each $j\in[p],i\in[N]$, there exist $\tilde{\alpha}_{ji},\tilde{\alpha}(X_i)$ on the line segment between $\hat{\alpha}(X_i)$ and $\alpha_0(X_i)$ such that 
$$\begin{array}{rl}
&\quad\displaystyle W^{FO}_j(\alpha_{*x}(X_i),Y_i)= W_j(\hat \alpha(X_i),Y_i) + (\alpha_{*x}(X_i)-\hat\alpha(X_i))'\nabla_\alpha W_j(\hat\alpha(X_i),Y_i)\\
&\displaystyle= W_j(\hat \alpha(X_i),Y_i) + (\alpha_0(X_i)-\hat\alpha(X_i))'\nabla_\alpha W_j(\hat\alpha(X_i),Y_i)+(\alpha_{*x}(X_i)-\alpha_0(X_i))'\nabla_\alpha W_j(\hat\alpha(X_i),Y_i)\\
&\displaystyle = W_j(\alpha_0(X_i),Y_i) - \frac{1}{2}(\alpha_0(X_i)-\hat\alpha(X_i))'\nabla_\alpha^2 W_j(\tilde\alpha_{ji},Y_i)(\alpha_0(X_i)-\hat\alpha(X_i))\\
&\quad\displaystyle +(\alpha_{*x}(X_i)-\alpha_0(X_i))'\nabla_\alpha W_j(\hat\alpha(X_i),Y_i),\ \mbox{and}
\end{array}$$
$\frac{W^{FO}_j(\alpha_{*x}(X_i),Y_i) }{1+W^{FO}(\alpha_{*x}(X_i),Y_i)'r_0(X_i)} =  \frac{W_j(\alpha_0(X_i),Y_i)}{1+W(\alpha_0(X_i),Y_i)'r_0(X_i)} + \widetilde R(\hat\alpha(X_i),\alpha_0(X_i),Y_i)$, where $|\widetilde R(\hat{\alpha}(X_i),\alpha_0(X_i),Y_i)|\leq C\left(1+\frac{|W_j(\alpha_0(X_i),Y_i)| }{1+W(\alpha_0(X_i),Y_i)'r_0(X_i)}\right) \max_{\alpha \in \mathcal{A}_x(\hat\alpha), \|X_i-x\|\leq h}\norm{\alpha(X_i)-\alpha_0(X_i)}^2+Ch^2$ with probability $1-\delta_{\alpha}$ from Assumption~\ref{ass:smoothness-assumptions}, for some universal constant $C$. 

Next, define 
$\tilde X_{ij} := \left(\frac{K_h(X_i-x)W_{j}(\alpha_0(X_i),Y_i)}{1+W(\alpha_0(X_i),Y_i)'r_0(X_i)},\frac{K_h(X_i-x) W_j(\alpha_0(X_i),Y_i)(X_i-x)'/h}{1+W(\alpha_0(X_i),Y_i)'r_0(X_i)}\right)\in\mathbb{R}^{d+1}$, where we denote $\tilde{X}_{ij0}:=\frac{K_h(X_i-x)W_{j}(\alpha_0(X_i),Y_i)}{1+W(\alpha_0(X_i),Y_i)'r_0(X_i)}$, and $\tilde{X}_{ijl}:=\frac{K_h(X_i-x) W_j(\alpha_0(X_i),Y_i)(X_{il}-x_{l})/h}{1+W(\alpha_0(X_i),Y_i)'r_0(X_i)}$ for any $l\in[d]$. So under Assumptions \ref{ass:iid-min-den} and \ref{ass:moment-assumptions}, we have 
$$\begin{array}{rl}
&\displaystyle\quad\min_{j\in[p],l\in\{0\}\cup[d]}\mathbb{E}[\tilde X_{ijl}^2 ]\\
&\displaystyle\geq_{(1)} \min_{j\in[p],k\in[d]}\mathbb{E}\Bigg[\frac{K(\norm{X_i-x}/h)^2\mathbf{1}\{\norm{X_i-x}\leq h\}}{h^{2d}}\frac{(X_{ik}-x_k)^2}{h^2}\left(\frac{W_j(\alpha_0(X_i),Y_i)}{1+W(\alpha_0(X_i),Y_i)'r_0(X_i)}\right)^2\Bigg]\\
&\displaystyle\geq_{(2)}\underline{C}_1\min_{j\in[p],k\in[d]}\mathbb{E}\left[\frac{\mathbf{1}\{\norm{X_i-x}\leq h\}}{h^{2d}}\frac{(X_{ik}-x_k)^2}{h^2}\left(\frac{W_j(\alpha_0(X_i),Y_i)}{1+W(\alpha_0(X_i),Y_i)'r_0(X_i)}\right)^2\right]\\
&\displaystyle\geq_{(3)}\underline{C}_1\bar\psi_2^2\int_{\norm{u-x}\leq h}\frac{(u_k-x_k)^2}{h^{2d}h^2}\rho(u)du\geq_{(4)}\underline{C}\bar\psi_2^2/h^d,
\end{array}$$
where (1) follows because $|X_{ik}-x_k|\leq h$, $K_h(X_i-x)>0\iff\norm{X_i-x}\leq h$ and $K_h(X_i-x)=\frac{1}{h^d}K(\norm{X_i-x}/h)$. (2) follows because $K(\norm{X_i-x}/h)$ is bounded away from zero on its support, (3) follows by definition of $\bar{\psi}_2$ in Assumption \ref{ass:moment-assumptions}, and $\rho(\cdot)$ is the density of covariate $X$, and (4) holds because the volume of the Euclidean ball $\{u:\norm{u-x}\leq h\}$ is equal to $Ch^d$, where $C$ is some constant independent of $h,d$, and that $\rho(\cdot)\geq\rho_0>0$ on its support by Assumption \ref{ass:local-linear}. Also note that  
$$\begin{array}{rl}
&\displaystyle\quad\max_{j\in[p],k\in\{0\}\cup[d]}\left(N^{-1}\sum_{i=1}^N\mathbb{E}[|\tilde X_{ijk}|^3 ]\right)^{1/3}\\ 
&\displaystyle\leq\max_{j\in[p]}\left(N^{-1}\sum_{i=1}^N\mathbb{E}\left[\frac{K(\norm{X_i-x}/h)^3\mathbf{1}\{\norm{X_i-x}\leq h\}}{h^{3d}}\left|\frac{W_j(\alpha_0(X_i),Y_i)}{1+W(\alpha_0(X_i),Y_i)'r_0(X_i)}\right|^3\right]\right)^{1/3}\\
&\displaystyle\leq\bar{C}\bar{K}_3h^{-2d/3},
\end{array}$$
where $\underline{C}$, $\bar{C}$ are absolute constants. So 
$$\begin{array}{rl}
&\displaystyle\mathbb{P} \left(\left| \frac{1}{N}\sum_{i=1}^N \frac{\tilde X_{ijk}}{w_{xjk}} \right| > \frac{\Phi^{-1}(1-\delta/2p(d+1))}{w_{xjk} N^{1/2}} \sqrt{\frac{1}{N}\sum_{i=1}^N \tilde X_{ijk}^2}, \mbox{for some } j\in [p], k\in\{0\}\cup[d] \right)  \\
&\displaystyle \leq  \sum_{j\in [p],k\in\{0\}\cup[d]} \mathbb{P} \left( \left|\frac{1}{N}\sum_{i=1}^N\frac{\tilde X_{ijk}}{w_{xjk}} \right| > \frac{\Phi^{-1}(1-\delta/2p(d+1))}{w_{xjk}N^{1/2}} \sqrt{\frac{1}{N}\sum_{i=1}^N \tilde X_{ijk}^2 } \right)\\
&\displaystyle\leq\sum_{j\in [p],k\in\{0\}\cup[d]}\!\!\!\!\!\!\!\!\!\!\!\!\left[1+A(\bar{K}_3/\bar{\psi}_2)^3N^{-1/2}h^{-d/2}(1+\Phi^{-1}(1-\delta/2p(d+1)))^{3}\right]\{\delta/2p(d+1)\}\\
&\displaystyle\leq \delta \left(1+\mathrm{o}(1)\right),
\end{array}$$
where the first inequality above follows from the union bound, the second inequality follows from Theorem 2.3 in \citet{jing2003self}, where $A$ is an absolute constant, and the last inequality follows since $p=2^M-M-1$ and $\sqrt{\frac{M^3}{Nh^d}}\leq\delta_N$ from Assumption~\ref{ass:moment-assumptions}, so $(1+\Phi^{-1}(1-\delta/(2p(d+1)))^3/\sqrt{Nh^d}\leq C_1\sqrt{\frac{\log^3(p(d+1))}{Nh^d}}\leq C_1\delta_N$ for some absolute constant $C_1$.
Further, for any $j\in[p],k\in\{0\}\cup[d]$, for some absolute constant $C$,
$$\begin{array}{rl}
\displaystyle\sqrt{\frac{1}{N}\sum_{i=1}^N \tilde X_{ijk}^2}&\displaystyle\leq\sqrt{\frac{1}{N}\sum_{i=1}^NK_h(X_i-x)^2\left(\frac{W_j(\alpha_0(X_i),Y_i)}{1+W(\alpha_0(X_i),Y_i)'r_0(X_i)}\right)^2}\\
&\displaystyle=\sqrt{\frac{1}{Nh^{d}}\sum_{i=1}^NK(\norm{X_i-x}/h)K_h(X_i-x)\left(\frac{W_j(\alpha_0(X_i),Y_i)}{1+W(\alpha_0(X_i),Y_i)'r_0(X_i)}\right)^2}\\
&\displaystyle\leq \frac{C}{\sqrt{h^d}}\sqrt{\frac{1}{N}\sum_{i=1}^NK_h(X_i-x)\left(\frac{W_j(\alpha_0(X_i),Y_i)}{1+W(\alpha_0(X_i),Y_i)'r_0(X_i)}\right)^2},
\end{array}$$
Thus the result follows. 
\end{proof}

\begin{lemma}\label{lemma:covariate:score*approx}
Suppose that $Z_i(X_i), i\in [N],$ are independent random variables such that \\
$\mathbb{E}[Z_i(X_i)\mid X_i]=0$,  $\sigma^2 \geq \max_{i\in[N]}\mathbb{E}[Z_i^2(X_i) \mid X_i]$. Then for a given $x\in \mathcal{X}$ we have that with some absolute constants $\bar{C},C_1$, with probability $1-\delta-C_1(\bar{K}/\sigma)^2/(Nh^d)$, where $\bar{K}=\mathbb{E}[\max_{i\in[N]}Z_i(X_i)^2]^{1/2}$, we have $\left|\frac{1}{N}\sum_{i=1}^N K_h(X_i-x) Z_i(X_i) \right|\leq\bar{C}\sigma\sqrt{\frac{\log(1/\delta)}{Nh^d}}$. In particular, under Assumption \ref{ass:iid-min-den}, \ref{ass:moment-assumptions}, with probability $\displaystyle1-\delta-C(\bar{K}_{\max}/\bar{K}_2)^2/(Nh^d)$, for some universal constant $C$ independent of $s,N,h,M,d$, we have $\displaystyle\left|\frac{1}{N}\sum_{i=1}^N K_h(X_i-x)\frac{ W(\alpha_0(X_i),Y_i)^\prime(r_{*x}(X_i)-r_0(X_i))}{1+W(\alpha_0(X_i),Y_i)^\prime r_0(X_i)}\right| \leq C \bar{K}_2sdh^2\sqrt{\frac{\log(1/\delta)}{Nh^d}}$.
\end{lemma}
\begin{proof}[Proof of Lemma \ref{lemma:covariate:score*approx}]
Let $X_1^\prime,\ldots,X_N^\prime$ be i.i.d. copies of $X_1,\ldots,X_N$, and $Z_i(X_i)^\prime$ be i.i.d. copies of $Z_i(X_i)$ for $i\in[N]$. Let $\mathbf{X}(N):=\{X_1,\ldots,X_N\}$ and $\mathbf{Z}(N):=\{Z_1(X_1),\ldots,Z_N(X_N)\}$, define $\displaystyle W = \mathbb{E}\left[\sum_{i=1}^N\left(K_h(X_i-x)Z_i(X_i)-K_h(X_i^\prime-x)Z_i(X_i)^\prime\right)^2\big| \mathbf{X}(N),\mathbf{Z}(N)\right]$.
Denote $Z:=\sum_{i=1}^NK_h(X_i-x)Z_i(X_i)$. 
Then by Theorem 12.3 of \citet{boucheron2003concentration}, for all $t\geq0$, we have $\displaystyle\mathbb{P}\left\{\frac{|Z|}{\sqrt{W}}\geq2\sqrt{t}\right\}\leq4e^{-t/4}$. So for any $\delta>0$, for some absolute constant $\bar{C}_0$, with probability $1-\delta$, we have 
$$\begin{array}{rl}
(*)=\displaystyle\frac{\left|\sum_{i=1}^N K_h(X_i-x) Z_i(X_i) \right|}{\sqrt{\mathbb{E}\left[\sum_{i=1}^N\left(K_h(X_i-x)Z_i(X_i)-K_h(X_i^\prime-x)Z_i(X_i)^\prime\right)^2\big| \mathbf{X}(N),\mathbf{Z}(N)\right]}}\leq \bar{C}_0\sqrt{\log(1/\delta)}.
\end{array}$$
$$\begin{array}{rl}
&\displaystyle\quad\mathbb{E}\left[\sum_{i=1}^N\left(K_h(X_i-x)Z_i(X_i)-K_h(X_i^\prime-x)Z_i(X_i)^\prime\right)^2\big| X_1,\ldots,X_N\right]\\
&\displaystyle\leq\mathbb{E}\left[2\sum_{i=1}^N\left(K_h(X_i-x)Z_i(X_i)\right)^2+2\sum_{i=1}^N\left(K_h(X_i^\prime-x)Z_i(X_i)^\prime\right)^2\big| X_1,\ldots,X_N\right]\\
&\displaystyle\leq_{(1)}2\sum_{i=1}^N\mathbb{E}\left[\left(K_h(X_i'-x)Z_i(X_i)'\right)^2\right]+2\sigma^2\sum_{i=1}^N\mathbb{E}\left[\frac{K((X_i-x)/h)^2\mathbf{1}\{\norm{X_i-x}\leq h\}}{h^{2d}}\big|X_i\right]\\
&\displaystyle\leq_{(2)}2\sigma^2\sum_{i=1}^N\left\{\mathbb{E}\left[\frac{K((X_i-x)/h)^2\mathbf{1}\{\|X_i-x\|\leq h\}}{h^{2d}}\right]+\mathbb{E}\left[\frac{K((X_i-x)/h)^2\mathbf{1}\{\norm{X_i-x}\leq h\}}{h^{2d}}\big|X_i\right]\right\}
\end{array}$$
where (1) follows since $X_i', Z_i(X_i)'$ are independent of $\{X_1,\ldots,X_N\}$ and $\sigma^2 \geq \max_{i\in[N]}\mathbb{E}[Z_i^2(X_i) \mid X_i]$. (2) follows since 
$$\begin{array}{rl}
\mathbb{E}\left[\left(K_h(X_i'-x)Z_i(X_i)'\right)^2\right]&\displaystyle=_{(a)}\mathbb{E}\left[\left(K_h(X_i-x)Z_i(X_i)\right)^2\right]\\
&\displaystyle=_{(b)}\mathbb{E}\{K_h(X_i-x)^2\mathbb{E}\left[Z_i(X_i)^2|X_i\right]\}\\
&\displaystyle\leq_{(c)}\sigma^2\mathbb{E}\left[\frac{K((X_i-x)/h)^2\mathbf{1}\{\|X_i-x\|\leq h\}}{h^{2d}}\right],
\end{array}$$ 
where (a) follows since $X_i', Z_i(X_i)'$ are identically distributed copies of $X_i,Z_i(X_i)$, (b) follows by tower property, and (c) follows $\sigma^2 \geq \max_{i\in[N]}\mathbb{E}[Z_i^2(X_i) \mid X_i]$. Further note that $0\leq K(\cdot)\leq\bar{C}$, and for some absolute constant $C_0$, $\mathbb{P}(\norm{X_i-x}\leq h)=\int_{B_h(x)}\rho(du)=C_0h^d$, hence
{$$\begin{array}{rl}
&\displaystyle\quad\mathbb{E}\left\{\mathbb{E}\left[\sum_{i=1}^N\left(K_h(X_i-x)Z_i(X_i)-K_h(X_i^\prime-x)Z_i(X_i)^\prime\right)^2\big| \mathbf{X}(N),\mathbf{Z}(N)\right]\right\}\\
&\displaystyle\leq\mathbb{E}\left[\frac{4\bar{C}^2\sigma^2}{h^{2d}}\sum_{i=1}^N\mathbf{1}\{\norm{X_i-x}\leq h\}\right]\leq\frac{4NC_0\bar{C}^2\sigma^2}{h^d}.
\end{array}$$}
Let $\displaystyle\bar{M}=\max_{i\in[N]}\mathbb{E}\left[K_h(X_i-x)^2Z_i(X_i)^2|X_1,\ldots,X_N,Z_1(X_1),\ldots,Z_N(X_N)\right]$ and $$\bar{M}^\prime=\max_{i\in[N]}\mathbb{E}\left[K_h(X_i^\prime-x)^2Z_i(X_i)^{\prime 2}|X_1,\ldots,X_N,Z_1(X_1),\ldots,Z_N(X_N)\right].$$ 
Then $\max_{i\in[N]}\mathbb{E}\left[(K_h(X_i-x)Z_i(X_i)-K_h(X_i^\prime-x)Z_i(X_i)^\prime)^2\big|\mathbf{X}(N),\mathbf{Z}(N)\right]\leq2(\bar{M}+\bar{M}^\prime)$.
Thus 
$$\mathbb{E}\left\{\max_{i\in[N]}\mathbb{E}\left[(K_h(X_i-x)Z_i(X_i)-K_h(X_i^\prime-x)Z_i(X_i)^\prime)^2\big|\mathbf{X}(N),\mathbf{Z}(N)\right]\right\}\leq2\mathbb{E}[\bar{M}+\bar{M}^\prime]\leq_{(a)}4\mathbb{E}[\bar{M}],$$
where (a) holds since $\mathbb{E}[\bar{M}']\leq\mathbb{E}[\bar{M}]$. Further note that $\mathbb{E}[\bar{M}]=\mathbb{E}[\max_{i\in[N]}K_h(X_i-x)^2Z_i(X_i)^2]\leq\frac{\bar{C}^2}{h^{2d}}\bar{K}^2$, where $\bar{K}=\mathbb{E}[\max_{i\in[N]}Z_i(X_i)^2]^{1/2}$. Thus using Lemma E.4 of \citet{chernozhukov2017central}, for any $t>0$, for some absolute constant $C_1$, with probability at least $\displaystyle1-C_1\frac{\bar{C}^2\bar{K}^2}{th^{2d}}$, we have 
$$\mathbb{E}\left[\sum_{i=1}^N\left(K_h(X_i-x)Z_i(X_i)-K_h(X_i^\prime-x)Z_i(X_i)^\prime\right)^2\big| \mathbf{X}(N),\mathbf{Z}(N)\right]\leq\frac{8NC_0\bar{C}^2\sigma^2}{h^d}+t,$$
indicating that with some absolute constants $\bar{C}_0$ and $\bar{C}_1$, with probability at least $\displaystyle1-\frac{C_1(\bar{K}/\sigma)^2}{Nh^d}$, $\displaystyle\mathbb{E}\left[\sum_{i=1}^N\left(K_h(X_i-x)Z_i(X_i)-K_h(X_i^\prime-x)Z_i(X_i)^\prime\right)^2\big|\mathbf{X}(N),\mathbf{Z}(N)\right]\leq\bar{C}_0\sigma^2\frac{N}{h^d}$. Thus further using the inequality for $(*)$ derived above, with some absolute constants $\bar{C},C_1$, with probability $1-\delta-\frac{C_1(\bar{K}/\sigma)^2}{Nh^d}$, $\displaystyle\left|\frac{1}{N}\sum_{i=1}^N K_h(X_i-x) Z_i(X_i) \right|\leq\bar{C}\sigma\sqrt{\frac{\log(1/\delta)}{Nh^d}}$. So the first statement holds. The second statement in the Lemma follows immediately as a corollary by letting $\displaystyle Z_i(X_i)=\frac{ W(\alpha_0(X_i),Y_i)^\prime(r_{*x}(X_i)-r_0(X_i))}{1+W(\alpha_0(X_i),Y_i)^\prime r_0(X_i)}$, so $\mathbb{E}[Z_i(X_i)|X_i]=0$. Recall that $\left\|  r_{*xj}(X_i)-r_{0j}(X_i) \right\|_1 \leq \bar{K}_{x,2}dh^2$ by component-wise Taylor expansion. Here we denote $\displaystyle\widetilde{W}_i=\frac{W(\alpha_0(X_i),Y_i)}{1+W(\alpha_0(X_i),Y_i)'r_0(X_i)}$ and let $T_i$ be the support of $r_{*x}(X_i)-r_0(X_i)$, we have $|T_i|\leq 2s$, so with some absolute constant $C$ we have $\mathbb{E}[\max_{i\in[N]}Z_i(X_i)^2]^{1/2}=\mathbb{E}\left[\max_{i\in[N]}\left(\sum_{j\in T_i}\widetilde{W}_{ij}(r_{*xj}(X_i)-r_{0j}(X_i))\right)^2\right]^{1/2}\!\!\!\!\!\!\!\!\leq C\bar{K}_{\max}sdh^2$. Also note that $\max_{i\in[N]}\mathbb{E}[Z_i^2(X_i)|X_i]\leq \max_{i\in[N]}\mathbb{E}\left[\left(\sum_{j\in T_i}\widetilde{W}_{ij}(r_{*xj}(X_i)-r_{0j}(X_i))\right)^2\big|X_i\right]\leq C\bar{K}_{2}^2s^2d^2h^4$. Since $\bar{K}=\mathbb{E}[{\displaystyle\max_{i\in[N]}}Z_i(X_i)^2]^{1/2}\leq C\bar{K}_{\max}sdh^2$, $\sigma^2\leq C\bar{K}_2^2s^2d^2h^4$. So the second statement holds.
\end{proof}

\begin{lemma}\label{lemma:MinMax:Claim1:covariate}
Suppose Assumptions~\ref{ass:iid-min-den}, \ref{ass:moment-assumptions}, \ref{ass:smoothness-assumptions} hold. Then for some universal constant $C$ independent of $N,M,s,d,h$, with probability $1-C\delta-C\epsilon-C\delta_{\alpha}$, we have 
$$\begin{array}{rl}
&\quad\displaystyle \min_{\alpha \in \mathcal{A}_x(\hat\alpha)} Q^{FO}(\alpha,r_*)-Q^{FO}(\alpha_{*x},r_*)\\
& \displaystyle 
 \geq-C\max_{\alpha \in \mathcal{A}_x(\hat\alpha), \|X_i-x\|\leq h}\|\alpha(X_i)-\alpha_0(X_i)\|_1 \sqrt{\frac{\log (M/\delta)}{Nh^d}}\\
&\quad \displaystyle {-} C\max_{\alpha \in \mathcal{A}_x(\hat\alpha), \|X_i-x\|\leq h}\!\!\|\alpha(X_i)\! - \!\alpha_0(X_i)\|^2 \!\!-Csdh^2\left(\frac{M}{\sqrt{Nh^{d}}}+Mdh^2\right),
\end{array}$$
where $\delta\geq\max\{\exp(-\bar{C}_0\{\sqrt{\log M}\sqrt{Nh^d}+Nh^d\}),\log(M)/(\delta_NNh^d),1/2^M\}$. 
\end{lemma}
\begin{proof}[Proof of Lemma~\ref{lemma:MinMax:Claim1:covariate}]
Take $\alpha(\cdot) \in \mathcal{A}_x(\hat\alpha)$, where $\alpha(X_i):=\tilde{\alpha}+\tilde{\beta}(X_i-x)$ for some $\tilde{\alpha}\in\mathbb{R}^M, \tilde{\beta}\in\mathbb{R}^{M\times d}$. Let $t=\frac{(\alpha(X_i) - \alpha_{*x}(X_i))'\nabla_\alpha W(\hat\alpha(X_i),Y_i)'r_{*x}(X_i)}{1+W^{FO}(\alpha_{*x}(X_i),Y_i)'r_{*x}(X_i)}$. Then $Q^{FO}(\alpha,r_*)-Q^{FO}(\alpha_{*x},r_*)=\frac{1}{N}\sum_{i=1}^NK_h(X_i-x) \log\left( 1 + t\right)$. By (1) applying Taylor expansions with respect to $W$, $\nabla_{\alpha}W$, (2) note that $\hat{\alpha}(X_i)\in\mathcal{A}_x(\hat{\alpha})$, (3) $\|r_{*x}(X_i)-r_0(X_i)\|_1\leq\bar{K}sdh^2$ for some universal constant $\bar{K}$, then following similar proof steps in Lemma~\ref{lemma:MinMax:Claim1}, it is straightforward to check that when $N$ is sufficiently large, 
$$\begin{array}{rl}
&\displaystyle\quad\frac{(\alpha(X_i) - \alpha_{*x}(X_i))'\nabla_\alpha W(\hat\alpha(X_i),Y_i)'r_{*x}(X_i)}{1+W^{FO}(\alpha_{*x}(X_i),Y_i)'r_{*x}(X_i)}\\
&\displaystyle=\frac{(\alpha(X_i) - \alpha_{*x}(X_i))'\nabla_\alpha W(\alpha_0(X_i),Y_i)'r_{*x}(X_i)}{1+W(\alpha_0(X_i),Y_i)'r_{*x}(X_i)} +\widetilde R(\alpha,\hat\alpha,Y_i),
\end{array}$$ 
where $\left|\frac{(\alpha(X_i) - \alpha_{*x}(X_i))'\nabla_\alpha W(\alpha_0(X_i),Y_i)'r_{*x}(X_i)}{1+W(\alpha_0(X_i),Y_i)'r_{*x}(X_i)} +\widetilde R(\alpha,\hat\alpha,Y_i)\right|\leq c_0$ for an absolute constant $c_0>0$. Furthermore note that $\log(1+t)\geq t - \frac{t^2}{1+t}$, so 
$Q^{FO}(\alpha,r_*)-Q^{FO}(\alpha_{*x},r_*)\geq \Delta q_1+\Delta q_2-\Delta q_3$, where $\Delta q_1=\frac{1}{N}\sum_{i=1}^NK_h(X_i-x)\frac{(\alpha(X_i) - \alpha_{*x}(X_i))'\nabla_\alpha W(\alpha_0(X_i),Y_i)'r_{*x}(X_i)}{1+W(\alpha_0(X_i),Y_i)'r_{*x}(X_i)}$, $\Delta q_2=\frac{1}{N}\sum_{i=1}^N K_h(X_i-x)\widetilde R(\alpha,\hat\alpha,Y_i)$, $\Delta q_3=\frac{1/c_0}{N}\sum_{i=1}^N K_h(X_i-x)\left|\frac{(\alpha(X_i) - \alpha_{*x}(X_i))'\nabla_\alpha W(\alpha_0(X_i),Y_i)'r_{*x}(X_i)}{1+W(\alpha_0(X_i),Y_i)'r_{*x}(X_i)} +\widetilde R(\alpha,\hat\alpha,Y_i))\right|^2$. In the following we focus on bounding $|\Delta q_1|$, $|\Delta q_2|$ and $|-\Delta q_3|$ respectively. Note that $\Delta q_1=\Delta q_{11}+\Delta q_{12}$, where $\displaystyle\Delta q_{11}:=\frac{1}{N}\sum_{i=1}^N K_h(X_i-x)\frac{(\alpha(X_i) - \alpha_{*x}(X_i))'\nabla_\alpha W(\alpha_0(X_i),Y_i)'r_0(X_i)}{1+W(\alpha_0(X_i),Y_i)'r_0(X_i)}$ and 
$$\begin{array}{rl}
\displaystyle\Delta q_{12}\!:=\!\frac{1}{N}\!\!\sum_{i=1}^N\!\!K_h(X_i-x)(\alpha(X_i) - \alpha_{*x}(X_i))'\!\!\left\{\!\frac{\nabla_\alpha W(\alpha_0(X_i),Y_i)'r_{*x}(X_i)}{1+W(\alpha_0(X_i),Y_i)'r_{*x}(X_i)}\! -\! \frac{ \nabla_\alpha W(\alpha_0(X_i),Y_i)'r_0(X_i)}{1+W(\alpha_0(X_i),Y_i)'r_0(X_i)}\!\right\}
\end{array}$$
So with probability $1-C\delta-C(\bar{K}_{\max}/\bar{K}_2)^2/(Nh^d)$,
$$\begin{array}{rl}
&\quad\displaystyle|\Delta q_{11}|:= \left|\frac{1}{N}\sum_{i=1}^N K_h(X_i-x)\frac{(\alpha(X_i) - \alpha_{*x}(X_i))'\nabla_\alpha W(\alpha_0(X_i),Y_i)'r_0(X_i)}{1+W(\alpha_0(X_i),Y_i)'r_0(X_i)} \right|\\
&\displaystyle\leq_{(1)}\| (\tilde{\alpha}-\alpha_0(x),h\cdot\mathrm{vec}(\tilde{\beta}-\nabla \alpha_0(x))\|_1\left\| \frac{1}{N}\sum_{i=1}^N\frac{K_h(X_i-x)\left( { \nabla_\alpha W(\alpha_{0}(X_i),Y_i)'r_0(X_i)}\atop{ {\rm vec}( \nabla_\alpha W(\alpha_{0}(X_i),Y_i)'r_0(X_i)\frac{(X_i-x)'}{h} )}\right) }{1+W(\alpha_{0}(X_i),Y_i)^\prime r_0(X_i)}\right\|_\infty \\
&\displaystyle\leq_{(2)}C'\left(\norm{\tilde{\alpha}-\alpha_0(x)}_1+h\norm{\tilde{\beta}-\nabla\alpha_0(x)}_{1,1}\right)\sqrt{\frac{\log (M/\delta)}{N{h^d}}}\leq_{(3)}C\left(\frac{Md}{\sqrt{Nh^{d}}}+Mdh^2\right)\sqrt{\frac{\log (M/\delta)}{N{h^d}}}.
\end{array}$$
Next, we show that (1) and (2) hold above. 
To show that (1) holds above, note that $\alpha(X_i)=\tilde{\alpha}+\tilde{\beta}(X_i-x)$ where $\alpha(X_i)\in\mathcal{A}_x(\hat{\alpha})$, and recall that $\alpha_{*x}(X_i)=\alpha_0(x)+\nabla\alpha_0(x)(X_i-x)$, thus 
$$\begin{array}{rl}
&\displaystyle\quad\left|\frac{1}{N}\sum_{i=1}^N K_h(X_i-x)\frac{(\alpha(X_i) - \alpha_{*x}(X_i))'\nabla_\alpha W(\alpha_0(X_i),Y_i)'r_0(X_i)}{1+W(\alpha_0(X_i),Y_i)'r_0(X_i)}\right|\\
&\displaystyle=\left|\frac{1}{N}\sum_{i=1}^N K_h(X_i-x)\frac{[(\tilde{\alpha}-\alpha_0(x))+(\tilde{\beta}-\nabla\alpha_0(x))(X_i-x)]'\nabla_\alpha W(\alpha_0(X_i),Y_i)'r_0(X_i)}{1+W(\alpha_0(X_i),Y_i)'r_0(X_i)}\right|\\
&\displaystyle\leq_{(a)}\| (\tilde{\alpha},h\mathrm{vec}(\tilde{\beta}))-(\alpha_0(x),h\mathrm{vec}(\nabla \alpha_0(x)))\|_1\\ 
&\displaystyle\quad\times\left\| \frac{1}{N}\sum_{i=1}^N\frac{  K_h(X_i-x) }{1+W(\alpha_{0}(X_i),Y_i)^\prime r_0(X_i)} \left( { \nabla_\alpha W(\alpha_{0}(X_i),Y_i)'r_0(X_i)}\atop{ {\rm vec}( \nabla_\alpha W(\alpha_{0}(X_i),Y_i)'r_0(X_i)\frac{(X_i-x)'}{h} )}\right) \right\|_\infty, 
\end{array}$$
where inequality (a) here holds by applying H\"older's inequality. Additionally, inequality (2) holds with probability $1-C\delta-C(\bar{K}_{\max}/\bar{K}_2)^2/(Nh^d)$ by the self-normalized deviation theory and Lemma \ref{lemma:covariate:score*approx}. The constant $C'$ in (2) is a universal constant independent of $N,M,s,d$. And inequality (3) holds by Proposition \ref{prop:concentration:covariate} for $\delta\geq1/2^M$. Denote $\displaystyle\Delta^{*}(X_i,Y_i,x):=\frac{ \nabla_\alpha W(\alpha_0(X_i),Y_i)'r_{*x}(X_i)}{1+W(\alpha_0(X_i),Y_i)'r_{*x}(X_i)} - \frac{ \nabla_\alpha W(\alpha_0(X_i),Y_i)'r_0(X_i)}{1+W(\alpha_0(X_i),Y_i)'r_0(X_i)}$. Then for universal constants $\bar{C},\bar{C}'$, $C$,
$$\begin{array}{rl}
&\displaystyle\quad|\Delta q_{12}|= \left|\frac{1}{N}\sum_{i=1}^N K_h(X_i-x)(\alpha(X_i) - \alpha_{*x}(X_i))'\Delta^{*}(X_i,Y_i,x)\right|\\
&\displaystyle\leq_{(1)}\norm{(\tilde{\alpha},h\mathrm{vec}(\tilde{\beta}))-(\alpha_0(x),h\mathrm{vec}(\nabla\alpha_0(x)))}_1\bigg\| \frac{1}{N}\sum_{i=1}^NK_h(X_i-x)\left({\Delta^{*}(X_i,Y_i,x)}\atop{ {\rm vec}\left(\Delta^{*}(X_i,Y_i,x)\frac{(X_i-x)'}{h} \right)}\right)\bigg\|_\infty\\
&\displaystyle\leq_{(2)}\bar{C}'\left(\frac{Md}{\sqrt{Nh^{d}}}+Mdh^2\right)sdh^2\bigg(1 \!+\!\frac{\sqrt{1/\delta}\log M}{Nh^d}\!\bigg)
\leq_{(3)} Csdh^2\left(\frac{Md}{\sqrt{Nh^{d}}}+Mdh^2\right),
\end{array}$$
where inequality (1) for $|\Delta q_{12}|$ holds by the fact that $|(X_{ik}-x_k)/h|\leq1$ for any $k\in[d]$ and then using similar argument as for inequality (1) of $|\Delta q_{11}|$. Inequality (2) for $|\Delta q_{12}|$ holds by Proposition \ref{prop:concentration:covariate} for $\delta\geq1/2^M$ and Lemma \ref{Lemma:xi-bound}, and (3) holds since $\delta\geq\log(M)/Nh^d$. $|\Delta q_2|$ is bounded via $\frac{1}{N}\sum_{i=1}^N K_h(X_i-x)\widetilde R(\alpha,\hat\alpha,Y_i)\leq\frac{C}{N}\sum_{i=1}^N K_h(X_i-x)\max_{\alpha \in \mathcal{A}_x(\hat\alpha), \|X_i-x\|\leq h}\| \alpha(X_i)-\alpha_{*x}(X_i)\|^2$. Finally, following similar argument as before, $|-\Delta q_3|$ is bounded by $\frac{C'}{N}\sum_{i=1}^N K_h(X_i-x)[\max_{\alpha \in \mathcal{A}_x(\hat\alpha), \|X_i-x\|\leq h}\|\alpha(X_i)-\alpha_0(X_i)\|^2  + |\widetilde R(\alpha,\hat\alpha,Y_i)|^2]\leq C_2\max_{\alpha \in \mathcal{A}_x(\hat\alpha), \|X_i-x\|\leq h}\|\alpha(X_i)-\alpha_0(X_i)\|^2$ for some universal constant $C_2$ according to Assumption~\ref{ass:smoothness-assumptions}. Thus the lemma holds by combining all the results above.
\end{proof}

\begin{lemma}\label{Lemma:xi-bound}
Suppose Assumptions \ref{ass:iid-min-den}, \ref{ass:smoothness-assumptions} hold, and $sdh^2 = \delta_N\xi$ for a sequence $\delta_N$ that goes zero. Then for $N$ sufficiently large, for some universal constant $C$ independent of $s,N,M,d$, with probability $1-C\epsilon-C\delta-C\delta_{\alpha}$, 
$$\begin{array}{rl}
&\displaystyle\quad\left\|\frac{1}{N}\sum_{i=1}^NK_h(X_i-x)\left[\frac{\nabla_{\alpha}W(\alpha_0(X_i),Y_i)'r_{*x}(X_i)}{1+W(\alpha_0(X_i),Y_i)'r_{*x}(X_i)}-\frac{\nabla_{\alpha}W(\alpha_0(X_i),Y_i)'r_0(X_i)}{1+W(\alpha_0(X_i),Y_i)'r_0(X_i)}\right]\right\|_\infty\\
&\displaystyle\leq Csdh^2\bigg(1 \!+\!\frac{\sqrt{1/\delta}\log M}{Nh^d}\!\bigg).
\end{array}$$
\end{lemma}
\begin{proof}[Proof of Lemma \ref{Lemma:xi-bound}]
Firstly note that $K(\cdot)$ is supported on the unit ball, so we only need to consider the covariates $X_i$ such that $\norm{X_i-x}_2\leq h$. The proof of Theorem \ref{thm:main:covariate} implies that  
$$\norm{r_0(X_i)-r_0(x)-\nabla r_0(x)(X_i-x)}_1\leq \bar{K}_{x,2}sdh^2.$$ 
Furthermore, using H\"older's inequality, for any $k\in[M]$ we have 
$$\left|\nabla_{\alpha_k} W(\alpha_0(X_i),Y_i)^\prime\left(r_0(X_i)-r_{*x}(X_i)\right)\right|\leq \bar{K}_{x,2}\bar{K}_wsdh^2,$$
\begin{equation}\label{eq:feasible-r*}
\begin{array}{rl}
&\quad1+W(\alpha_0(X_i),Y_i)^\prime r_{*x}(X_i)\\
&\geq1+W(\alpha_0(X_i),Y_i)^\prime r_0(X_i)\\
&\quad\displaystyle-|W(\alpha_0(X_i),Y_i)^\prime\left(r_0(X_i)-r_{*x}(X_i)\right)|\\
&\geq_{(a)}\xi-\bar{K}_{x,2}\bar{K}_wsdh^2\\
&\geq_{(b)} C_1^\prime\xi,
\end{array}
\end{equation}
where (a) holds with probability $1-\epsilon$ by Assumption \ref{ass:iid-min-den} and (b) holds with $C_1'\in(0,1)$ for $N$ sufficiently large since $sdh^2=\delta_N\xi$ where $\delta_N\rightarrow0$. Given $i\in[N], k\in[M]$, define $\Delta^{*}(X_i,Y_i,x):=\frac{\nabla_{\alpha} W(\alpha_0(X_i),Y_i)' r_{*x}(X_i)}{1+W(\alpha_0(X_i),Y_i)^\prime r_{*x}(X_i)}-\frac{\nabla_{\alpha} W(\alpha_0(X_i),Y_i)'r_0(X_i)}{1+W(\alpha_0(X_i),Y_i)^\prime r_0(X_i)}$, $\Delta_0r:=\frac{r_{*x}(X_i)}{1+W(\alpha_0(X_i),Y_i)^\prime r_{*x}(X_i)}-\frac{r_0(X_i)}{1+W(\alpha_0(X_i),Y_i)^\prime r_0(X_i)}$. Then 
\begin{equation}\label{eq:target:bound}
\begin{array}{rl}
    &\displaystyle\quad\left\|\frac{1}{N}\sum_{i=1}^NK_h(X_i-x)\Delta^{*}(X_i,Y_i,x)\right\|_\infty\\
    &\displaystyle=\left\|\frac{1}{N}\sum_{i=1}^NK_h(X_i-x)\nabla_{\alpha} W(\alpha_0(X_i),Y_i)'\Delta_0r\right\|_\infty.
\end{array}
\end{equation}
Note that $\Delta_0r=\frac{r_{*x}(X_i)-r_0(X_i)}{(1+W(\alpha_0(X_i),Y_i)'\tilde{r}_i)}-\frac{\tilde{r}_iW(\alpha_0(X_i),Y_i)'(r_{*x}(X_i)-r_0(X_i))}{(1+W(\alpha_0(X_i),Y_i)'\tilde{r}_i)^2}$ for some $\tilde{r}_i$ on the line segment between $r_{*x}(X_i), r_0(X_i)$. Following similarly from \eqref{eq:feasible-r*}, $1+W(\alpha_0(X_i),Y_i)'\tilde{r}_i>C_1'\xi$ for some absolute constant $C_1'$ for $N$ sufficiently large. The condition of this lemma implies $\|\nabla_{\alpha}W(\alpha_0(X_i),Y_i)'(r_{*x}(X_i)-r_0(X_i))\|_{\infty}\leq \bar{K}_w\|r_{*x}(X_i)-r_0(X_i)\|_1\leq\bar{K}_{x,2}\bar{K}_wsdh^2$. Define $\Delta_kr:=\nabla_{\alpha_k}W(\alpha_0(X_i),Y_i)'\Delta_0r$. Assumption~\ref{ass:iid-min-den} and Assumption~\ref{ass:smoothness-assumptions} imply that $|\Delta_kr|\leq\bar{K}sdh^2$ for some universal constant $\bar{K}$. Define $\bar \sigma^2:=\max_{k\in[M]}\sum_{i=1}^N\mathbb{E}\left[K_h(X_i-x)^2\{\Delta_kr\}^2\right]$. So for some universal constant $C_2$, $\bar{\sigma}^2\leq C_2Ns^2d^2h^4/h^d$. Define $\bar{M}:=  \max_{i\in[N],k\in[M]}K_h(X_i-x)|\nabla_{\alpha_k}W(\alpha_0(X_i),Y_i)'\Delta_0r|$. Then $\displaystyle\bar{M}\leq C_2sd/h^{d-2}$.
Lemma 8 of \citet{chernozhukov2015comparison} implies $\mathbb{E}\left[\max_{k\in[M]}\left|\sum_{i=1}^NK_h(X_i-x)\Delta_kr-\mathbb{E}[K_h(X_i-x)\Delta_kr]\right|\right]\leq\bar{C}_1\left(\bar{\sigma}\sqrt{\log M}+\mathbb{E}[\bar{M}^2]^{1/2}\log M\right)$ for some universal constant $\bar{C}_1$. Since $|\Delta_kr|\leq\bar{K}sdh^2$, 
$$\mathbb{E}[\max_{k\in[M]}|\sum_{i=1}^NK_h(X_i-x)\Delta_kr|]\leq\bar{C}_2\left(\bar{\sigma}\sqrt{\log M}+\mathbb{E}[\bar{M}^2]^{1/2}\log M+Nsdh^2\right)$$ 
for some universal constant $\bar{C}_2$. Furthermore using Lemma E.4 of \citet{chernozhukov2017central}, for some universal constant $\bar{C}$, with probability $1-\bar{C}\mathbb{E}[\bar M^2]/t^2$, 
\begin{equation}\label{eq:bound:term-1}
\begin{array}{rl}
&\displaystyle\max_{k\in[M]}\left|\frac{1}{N}\sum_{i=1}^NK_h(X_i-x)\Delta_kr\right|\\
&\displaystyle\leq\frac{2\bar{C}_2\bar\sigma \sqrt{\log M}}{N} + \frac{2\bar{C}_2\mathbb{E}[\bar{M}^2]^{1/2}\log(M)}{N}\\
&\displaystyle\quad+2\bar{C}_2sdh^2 + \frac{t}{N}\\
&\displaystyle \leq\bar{C}sdh^2 + \frac{t}{N}.
\end{array}
\end{equation}
Now we take $t$ in (\ref{eq:bound:term-1}) such that $\bar{C}\mathbb{E}[\bar M^2]/t^2=\delta$, i.e. $t=\sqrt{\bar{C}/\delta}\mathbb{E}[\bar{M}^2]^{1/2}$. So by (\ref{eq:bound:term-1}), with probability $1-\delta$, 
\begin{equation}\label{eq:term-1:bound:new}
\begin{array}{rl}
&\displaystyle\max_{k\in[M]}\left|\frac{1}{N}\sum_{i=1}^NK_h(X_i-x)\Delta_kr\right|\\
&\displaystyle\leq Csdh^2\bigg(1 \!+\!\frac{\sqrt{1/\delta}\log M}{Nh^d}\!\bigg). 
\end{array}
\end{equation}
Recall $\displaystyle\Delta^{*}(X_i,Y_i,x)=\frac{ \nabla_\alpha W(\alpha_0(X_i),Y_i)'r_{*x}(X_i)}{1+W(\alpha_0(X_i),Y_i)'r_{*x}(X_i)} - \frac{ \nabla_\alpha W(\alpha_0(X_i),Y_i)'r_0(X_i)}{1+W(\alpha_0(X_i),Y_i)'r_0(X_i)}$.  (\ref{eq:target:bound}), (\ref{eq:term-1:bound:new}) imply that with probability $1-\delta-\epsilon$, $\displaystyle\left\|\frac{1}{N}\sum_{i=1}^NK_h(X_i-x)\Delta^{*}(X_i,Y_i,x)\right\|_\infty\!\!\!\!\!\leq Csdh^2\bigg(1 \!+\!\frac{\sqrt{1/\delta}\log M}{Nh^d}\!\bigg)$.
\end{proof}

\section{Rates of Convergence for Marginal Probabilities}\label{appendix:marginal}
In this section, we provide rates of convergence for the nuisance parameters or functions, i.e., marginal probabilities. Proposition \ref{prop:concentration:no covariate} and Proposition~\ref{prop:concentration:covariate} provide results the no-covariate case and with-covariate case respectively. The proofs for the following two propositions are standard by applying the concentration inequalities for $\hat{\alpha}$ defined as \eqref{eq:SAA}, $\hat{\alpha}(x)$ defined as \eqref{eq:truncate-local-linear}, $\tilde\beta(x)$ defined as in Section~\ref{sec:marginals}, following from \citep{boucheron2003concentration,einmahl2005uniform,hansen2008uniform,fan1993local}, so we omit them here. 
\begin{proposition}\label{prop:concentration:no covariate} In the case of no covariates, suppose Assumption \ref{ass:iid-min-den} holds. Let $\hat{\alpha}$ and $\mathcal{A}(\hat\alpha)$ be the estimator and confidence set defined in \eqref{eq:SAA} and $\delta \in (0,1)$. Then with probability at least $1-\delta$, we have $\norm{\hat{\alpha}-\alpha_0}_{\infty}\leq c_1\sqrt{\frac{\log(M/\delta)}{N}}$, $\norm{\hat\alpha-\alpha_0}_1\leq 
c_1\left(\frac{M}{\sqrt{N}}+\sqrt{\frac{M\log(1/\delta)}{N}}\right)$, $\norm{\hat\alpha-\alpha_0}_2\leq 
c_1\left(\sqrt{\frac{M}{N}}+\sqrt{\frac{\log(1/\delta)}{N}}\right)$, 
where $c_1$ is an absolute constant independent of $M,N,s$. 
\end{proposition}

\begin{proposition}\label{prop:concentration:covariate}
Suppose Assumption \ref{ass:local-linear} holds and let $x$ be a given data point of in the interior of $\mathcal{X}$. Let $\hat{\alpha}(x)$ be the truncated local linear estimator defined by (\ref{eq:truncate-local-linear}), and let 
$\delta\geq\exp(-\bar{C}_0\{\sqrt{\log M}\sqrt{Nh^d}+Nh^d\})$. Then for some absolute constants $\bar{C}_0$, $\bar{C}_1$ which depend on $x$ while are independent of $N,d,h,s,M$, with probability $1-\delta$, we have\\ $\norm{\hat{\alpha}(x)-\alpha_0(x)}_1\leq\bar{C}_0\left(\frac{M}{\sqrt{Nh^d}}+Mdh^2+\sqrt{M\left(\frac{1}{Nh^d}+h^4d^2\right)\log(\frac{1}{\delta})}\right)$, \\
$\norm{\hat{\alpha}(x)-\alpha_0(x)}_2\!\leq\!\bar{C}_0\left(\sqrt{\frac{M}{Nh^d}}+\sqrt{M}dh^2+\sqrt{\left(\frac{1}{Nh^d}+h^4d^2\right)\log(\frac{1}{\delta})}\right)$, \\
$\norm{\hat{\alpha}(x)-\alpha_{0}(x)}_{\infty}\leq\bar{C}_0\left(dh^2+\sqrt{\frac{\log(M/\delta)}{Nh^d}}\right)$, \\
$\norm{\tilde{\beta}(x)-\nabla\alpha_0(x)}_{1,1}\leq\bar{C}_1d\left(\frac{M}{\sqrt{Nh^{d+2}}}+\!\!Mdh^2\!\!+\!\!\sqrt{M\left(\frac{1}{Nh^{d+2}}+h^4d^2\right)\log(\frac{1}{\delta})}\right)$,\\
$\norm{\tilde{\beta}(x)-\nabla\alpha_0(x)}_{2,2}\leq\bar{C}_1\sqrt{d}\left(\sqrt{\frac{M}{Nh^{d+2}}}+\sqrt{M}dh^2+\!\!\sqrt{\left(\frac{1}{Nh^{d+2}}+h^4d^2\right)\log(\frac{1}{\delta})}\right)$.
\end{proposition}

\section{Proof of Theorem \ref{thm:causal-inference:AIPW}}\label{sec:results:ATE}
For any fold $k\in[K]$ and any two treatment levels $t,t'$, define 
$$\begin{array}{rcl}
\displaystyle\hat{\eta}_{(t)}^{\mathcal{I}_k}&:=&\displaystyle\frac{1}{|\mathcal{I}_k|}\sum_{i\in\mathcal{I}_k}\left(\hat{\mu}_{(t)}^{-k}(X_i)+\mathbf{1}\{T_i=t\}\frac{O_i-\hat{\mu}_{(t)}^{-k}(X_i)}{\hat{e}_t^{-k}(X_i)}\right),\\
\displaystyle\hat{\eta}_{(t')}^{\mathcal{I}_k}&:=&\displaystyle\frac{1}{|\mathcal{I}_k|}\sum_{i\in\mathcal{I}_k}\left(\hat{\mu}_{(t')}^{-k}(X_i)+\mathbf{1}\{T_i=t'\}\frac{O_i-\hat{\mu}_{(t')}^{-k}(X_i)}{\hat{e}_{t'}^{-k}(X_i)}\right),\\
\displaystyle\hat{\eta}_{(t)}^{\mathcal{I}_k,*}&:=&\displaystyle\frac{1}{|\mathcal{I}_k|}\sum_{i\in\mathcal{I}_k}\left({\mu}_{(t)}(X_i)+\mathbf{1}\{T_i=t\}\frac{O_i-{\mu}_{(t)}(X_i)}{{e}_t(X_i)}\right),\\
\displaystyle\hat{\eta}_{(t')}^{\mathcal{I}_k,*}&:=&\displaystyle\frac{1}{|\mathcal{I}_k|}\sum_{i\in\mathcal{I}_k}\left({\mu}_{(t')}(X_i)+\mathbf{1}\{T_i=t'\}\frac{O_i-{\mu}_{(t')}(X_i)}{{e}_{t'}(X_i)}\right).
\end{array}$$
Define $\hat{\tau}^{\mathcal{I}_k}(t,t'):=\hat{\eta}_{(t)}^{\mathcal{I}_k}-\hat{\eta}_{(t')}^{\mathcal{I}_k}$, $\hat{\tau}^{\mathcal{I}_k,*}(t,t'):=\hat{\eta}_{(t)}^{\mathcal{I}_k,*}-\hat{\eta}_{(t')}^{\mathcal{I}_k,*}$, $\hat{\tau}_{\textrm{AIPW}}^*:=\sum_{k\in[K]}\frac{|\mathcal{I}_k|}{N}\hat{\tau}^{\mathcal{I}_k,*}(t,t')$. Note that $\hat{\tau}_{\textrm{AIPW}}=\sum_{k\in[K]}\frac{|\mathcal{I}_k|}{N}\hat{\tau}^{\mathcal{I}_k}(t,t')$, so $\hat{\tau}_{\textrm{AIPW}}^*=\frac{1}{N}\sum_{i=1}^N \mu_{(t)}(X_i)-\mu_{(t')}(X_i)+ \frac{1\{T_i=t\}}{e_t(X_i)}(O_i - \mu_{(t)}(X_i)) - \frac{1\{T_i=t'\}}{e_{t'}(X_i)}(O_i - \mu_{(t')}(X_i))$. To show $\sqrt{N}(\hat{\tau}_{\textrm{AIPW}}-\hat{\tau}_{\textrm{AIPW}}^*)=\mathrm{o}_{\mathrm{P}}(1)$, it suffices to show that for any $k\in[K]$, $\sqrt{N}(\hat{\tau}^{\mathcal{I}_k}(t,t')-\hat{\tau}^{\mathcal{I}_k,*}(t,t'))=\mathrm{o}_{\mathrm{P}}(1)$. Hence it suffices to show that for any $t\in\{0,1\}^M$, we have $\sqrt{N}(\hat{\eta}_{(t)}^{\mathcal{I}_k}-\hat{\eta}_{(t)}^{\mathcal{I}_k,*})=\mathrm{o}_{\mathrm{P}}(1)$. Note that  $\hat{\eta}_{(t)}^{\mathcal{I}_k}-\hat{\eta}_{(t)}^{\mathcal{I}_k,*}=(a)+(b)+(c)$, where $(a)=\frac{1}{|\mathcal{I}_k|}\sum_{i\in\mathcal{I}_k}\big\{\big(\hat{\mu}_{(t)}^{-k}(X_i)-\mu_t(X_i)\big)\big(1-\frac{\mathbf{1}\{T_i=t\}}{e_t(X_i)}\big)\big\}$, $(b)=\frac{1}{|\mathcal{I}_k|}\sum_{i\in\mathcal{I}_k}\mathbf{1}\{T_i=t\}\left((O_i-\mu_{(t)}(X_i))\left(\frac{1}{\hat{e}_t^{-k}(X_i)}-\frac{1}{e_t(X_i)}\right)\right)$, $(c)=-\frac{1}{|\mathcal{I}_k|}\sum_{i\in\mathcal{I}_k}\mathbf{1}\{T_i=t\}\left(\left(\hat{\mu}_{(t)}^{-k}(X_i)-\mu_t(X_i)\right)\left(\frac{1}{\hat{e}_t^{-k}(X_i)}-\frac{1}{e_t(X_i)}\right)\right)$. Note that 
$$\begin{array}{rl}
&\displaystyle\mathbb{E}\bigg[\bigg(\frac{1}{|\mathcal{I}_k|}\sum_{i\in\mathcal{I}_k}\bigg\{\bigg(\hat{\mu}_{(t)}^{-k}(X_i)-\mu_t(X_i)\bigg)\bigg(1-\frac{\mathbf{1}\{T_i=t\}}{e_t(X_i)}\bigg)\bigg\}\bigg)^2\bigg]\\
&\displaystyle=\frac{1}{|\mathcal{I}_k|}\mathbb{E}\left[\mathbb{E}\left[\left(\hat{\mu}_{(t)}^{-k}(X_i)-\mu_t(X_i)\right)^2\left(\frac{1}{e_t(X_i)}-1\right)\big|\mathcal{I}^{-k}\right]\right]\\
&\displaystyle\leq\frac{\mathbb{E}\big[\big(\hat{\mu}_{(t)}^{-k}(X_i)-\mu_t(X_i)\big)^2\big]/(c|\mathcal{I}_k|)}{\inf_{x\in\mathcal{X}}\prod_{j=1}^M \alpha_{0j}^{t_j}(x)(1-\alpha_{0j}(x))^{1-t_j}}\leq\frac{\tilde{C}\epsilon_N^2N^{-3/2}}{\inf_{x\in\mathcal{X}}\prod_{j=1}^M \alpha_{0j}^{t_j}(x)(1-\alpha_{0j}(x))^{1-t_j}},
\end{array}$$
where the last inequality follows from condition (iii) of Assumption~\ref{ass:causal:extra}. Further, define 
$$\Delta E:=\mathbb{E}\left[\left(\frac{1}{|\mathcal{I}_k|}\sum_{i\in\mathcal{I}_k}\mathbf{1}\{T_i=t\}\bigg(O_i-\mu_{(t)}(X_i)\bigg)\bigg(\frac{1}{\hat{e}_t^{-k}(X_i)}-\frac{1}{e_t(X_i)}\bigg)\right)^2\right].$$ Conditional on the out-of-fold sample, the fold-k summands are independent and have mean zero since $\mathbb{E}[O_i-\mu_{(t)}(X_i)|X_i,T_i=t]=0$, so
it is straightforward to check that $\Delta E\leq\frac{C}{|\mathcal{I}_k|}\mathbb{E}\left[\left(1-\frac{e_t(X_i)}{\hat{e}_t^{-k}(X_i)}\right)^2\frac{1}{e_t(X_i)}\right]$ for some absolute constant $C$. For any $x\in\mathcal{X}$, $\hat{e}_t(x)-e_t(x)=\textrm{(I)+(II)+(III)}$, where 
$$\begin{array}{rcl}
\textrm{(I)}&=&\displaystyle\left[\Big(1+W(\hat\alpha(x),t)'\hat r(x,x)\Big)-\Big(1+W(\alpha_0(x),t)'\hat r(x,x)\Big)\right]\prod_{j=1}^M\hat\alpha_j(x)^{t_j}(1-\hat\alpha_j(x))^{1-t_j},\\
\textrm{(II)}&=&\displaystyle\left[\Big(1+W(\alpha_0(x),t)'\hat r(x,x)\Big)-\Big(1+W(\alpha_0(x),t)'r_0(x)\Big)\right]\prod_{j=1}^M\hat\alpha_j(x)^{t_j}(1-\hat\alpha_j(x))^{1-t_j},\\
\textrm{(III)}&=&\displaystyle\Big(1+W(\alpha_0(x),t)'r_0(x)\Big)\left[\prod_{j=1}^M\hat\alpha_j(x)^{t_j}(1-\hat\alpha_j(x))^{1-t_j}-\prod_{j=1}^M\alpha_{0j}(x)^{t_j}(1-\alpha_{0j}(x))^{1-t_j}\right].
\end{array}$$
For any $j\in[M]$, define $\tilde{\Delta}_j\alpha_0(x):=\alpha_{0j}(x)^{t_j}(1-\alpha_{0j}(x))^{1-t_j}$, $\tilde{\Delta}_j\hat{\alpha}(x):=\hat{\alpha}_j(x)^{t_j}(1-\hat{\alpha}_j(x))^{1-t_j}$, then $\prod_{j=1}^M\tilde{\Delta}_j\hat{\alpha}(x)-\prod_{j=1}^M\tilde{\Delta}_j\alpha_0(x)=[\tilde{\Delta}_1\hat{\alpha}(x)-\tilde{\Delta}_1\alpha_0(x)]\prod_{j=2}^M\tilde{\Delta}_j\hat{\alpha}(x)+\tilde{\Delta}_1\alpha_0(x)\bigg\{\prod_{j=2}^M\tilde{\Delta}_j\hat{\alpha}(x)-\prod_{j=2}^M\tilde{\Delta}_j\alpha_0(x)\bigg\}$. Denote $\Delta_k := \prod_{j=k}^M\hat{\alpha}_j(x)^{t_j}(1-\hat{\alpha}_j(x))^{1-t_j}-\prod_{j=k}^M\alpha_{0j}(x)^{t_j}(1-\alpha_{0j}(x))^{1-t_j}$, $\hat{\delta}_j:=\hat{\alpha}_j(x)^{t_j}(1-\hat{\alpha}_j(x))^{1-t_j}-\alpha_{0j}(x)^{t_j}(1-\alpha_{0j}(x))^{1-t_j}$, $\hat{\pi}_k:=\prod_{j=k}^M\hat{\alpha}_j(x)^{t_j}(1-\hat{\alpha}_j(x))^{1-t_j},m_j:=\alpha_{0j}(x)^{t_j}(1-\alpha_{0j}(x))^{1-t_j}$. So by expanding the recursive formula of $\prod_{j=1}^M\tilde{\Delta}_j\hat{\alpha}(x)-\prod_{j=1}^M\tilde{\Delta}_j\alpha_0(x)$, we have $\Delta_1=\prod_{j=1}^M\hat{\alpha}_j(x)^{t_j}(1-\hat{\alpha}_j(x))^{1-t_j}-\prod_{j=1}^M\alpha_{0j}(x)^{t_j}(1-\alpha_{0j}(x))^{1-t_j}=\hat{\delta}_1\hat{\pi}_2+m_1\Delta_2=\hat{\delta}_1\hat{\pi}_2+m_1\hat{\delta}_2\hat{\pi}_3+m_1m_2\Delta_3=\ldots=\hat{\delta}_1\hat{\pi}_2+\sum_{j=2}^{M-1}\hat{\delta}_j\hat{\pi}_{j+1}\prod_{k=1}^{j-1}m_k+\hat{\delta}_M\prod_{k=1}^{M-1}m_k$. By noting that $0\leq\hat{\pi}_k\leq 1$ and $0\leq m_j\leq 1$ for any $j,k\in[M]$, and recall that $t_j\in\{0,1\}$, thus $|\Delta_1|\leq\sum_{j=1}^M|\hat{\delta}_j|=\sum_{j=1}^M|\hat{\alpha}_j(x)-\alpha_{0j}(x)|=\norm{\hat{\alpha}(x)-\alpha_0(x)}_1\leq CMdh^2+C\frac{M}{\sqrt{Nh^d}}$ with probability $1-\delta$ for $\delta\geq\max\{\exp(-C\{\sqrt{\log M}\sqrt{Nh^d}+Nh^d\}),1/2^M\}$ for some absolute constant $C$ according to Proposition~\ref{prop:concentration:covariate}. The above inequality and Taylor expansion imply $|\textrm{(I)}|\leq |(\hat{\alpha}(x)-\alpha_0(x))'\nabla_\alpha W(\tilde{\alpha}(x),t)'\hat{r}(x,x)|\prod_{j=1}^M\hat\alpha_j(x)^{t_j}(1-\hat\alpha_j(x))^{1-t_j}$. Since $\prod_{j=1}^M\hat\alpha_j(x)^{t_j}(1-\hat\alpha_j(x))^{1-t_j}\leq\prod_{j=1}^M\alpha_{0j}(x)^{t_j}(1-\alpha_{0j}(x))^{1-t_j}+\norm{\hat{\alpha}(x)-\alpha_0(x)}_1$, 
$$\begin{array}{rl}
|\textrm{(I)}|&\displaystyle\leq_{(1)} C'\left(\norm{\hat{\alpha}(x)-\alpha_0(x)}_1^2+\norm{\hat{\alpha}(x)-\alpha_0(x)}_1\sup_{x\in\mathcal{X}}\prod_{j=1}^M\alpha_{0j}(x)^{t_j}(1-\alpha_{0j}(x))^{1-t_j}\right)\\
&\displaystyle\leq C'\left(\frac{M}{\sqrt{Nh^d}}+Mdh^2\right)^2\!\!\!+C'\left(\frac{M}{\sqrt{Nh^d}}+Mdh^2\right)\sup_{x\in\mathcal{X}}\prod_{j=1}^M\alpha_{0j}(x)^{t_j}(1-\alpha_{0j}(x))^{1-t_j}\leq_{(2)}C\psi_N,
\end{array}$$ 
here (1) follows from Assumption \ref{ass:smoothness-assumptions}, the fact that $\hat{r}(x,x)=\hat{r}(x,x)-r_0(x)+r_0(x)$ and $\hat{r}(x,x)-r_0(x)$ is negligible by our theoretical result for the with-covariate estimator; (2) follows from (\ref{consistency:ATE:condition}) and (\ref{consistency:ATE:2}). To control (II), note $|\textrm{(II)}|\leq |W(\alpha_0(x),t)'(\hat{r}(x,x)-r_0(x))|\prod_{j=1}^M\hat\alpha_j(x)^{t_j}(1-\hat\alpha_j(x))^{1-t_j}$, where the right hand side is equal to $|W(\alpha_0(x),t)'(\hat{r}(x,x)-r_0(x))|\left[\prod_{j=1}^M\alpha_{0j}(x)^{t_j}(1-\alpha_{0j}(x))^{1-t_j}+\Delta_1\right]$. Let $\delta\geq\max\{\exp(-C'\{\sqrt{\log(M)Nh^d}+Nh^d\}),1/2^M\}$. Then according to Proposition~\ref{prop:concentration:covariate}, with probability $1-\delta$, we have $|\Delta_1|\leq C'Mdh^2+C'\frac{M}{\sqrt{Nh^d}}$. Further using H\"older's inequality, Corollary \ref{cor:FO:covariate}, Assumption~\ref{ass:smoothness-assumptions} and $r_0(x)=r_{*x}(x)$, with probability $1-\delta$, we have 
$$|\textrm{(II)}|\leq C_1'\sqrt{s}\left(\!\!\sqrt{\frac{sM}{Nh^d}}+(\sqrt{M}+s)dh^2\right)\!\left(\prod_{j=1}^M\alpha_{0j}(x)^{t_j}(1-\alpha_{0j}(x))^{1-t_j}\!+\!Mdh^2\!+\!\frac{M}{\sqrt{Nh^d}}\right)$$ for some universal constant $C_1'$ and $\delta\geq\log(M)/(\delta_NNh^d)$. Additionally, following from (\ref{consistency:ATE:condition}) and \eqref{consistency:ATE:2}, $\textrm{(II)}\leq C_2\psi_N$ for an absolute constant $C_2$. Further by Assumption \ref{ass:iid-min-den}, we have $|\textrm{(III)}|\leq C_3\norm{\hat{\alpha}(x)-\alpha_0(x)}_1\leq C_3\psi_N$ for an absolute constant $C_3$, hence $\displaystyle\sup_{x\in\mathcal{X}}|\hat{e}_t(x)-e_t(x)|\leq C\psi_N$, where $C$ is some constant independent of $N,d,h,M,s$. So the upper bound for $\Delta E$ derived earlier implies $$\begin{array}{rl}
&\quad\displaystyle\mathbb{E}\left[\left(\frac{1}{|\mathcal{I}_k|}\sum_{i\in\mathcal{I}_k}\mathbf{1}\{T_i=t\}(O_i-\mu_{(t)}(X_i))\left(\frac{1}{\hat{e}_t^{-k}(X_i)}-\frac{1}{e_t(X_i)}\right)\right)^2\right]\\
&\displaystyle\leq\frac{CN^{-1}\psi_N^2}{[\inf_{x\in\mathcal{X}}\prod_{j=1}^M \alpha_{0j}^{t_j}(x)(1-\alpha_{0j}(x))^{1-t_j}]^3}.
\end{array}$$ 
Recall that $\hat{\eta}_{(t)}^{\mathcal{I}_k}-\hat{\eta}_{(t)}^{\mathcal{I}_k,*}=(a)+(b)+(c)$, where (a), (b), (c) are defined as before. Hence by Chebyshev's inequality, for an absolute constant $\tilde{C}$, we have 
$$\mathbb{P}(\sqrt{N}|(a)|>\tilde{C} N^{-\gamma})\leq\frac{\tilde{C}\epsilon_N^2}{\inf_{x\in\mathcal{X}}\prod_{j=1}^M \alpha_{0j}^{t_j}(x)(1-\alpha_{0j}(x))^{1-t_j}}\leq\frac{\tilde{C}\epsilon_N^2}{c^M},$$ $$\mathbb{P}(\sqrt{N}|(b)|>\tilde{C}\sqrt{\psi_N})\leq\frac{\tilde{C}\psi_N}{[\inf_{x\in\mathcal{X}}\prod_{j=1}^M \alpha_{0j}^{t_j}(x)(1-\alpha_{0j}(x))^{1-t_j}]^3}\leq\frac{\tilde{C}\psi_N}{c^{3M}}.$$ Further, 
$$\begin{array}{rl}
&\displaystyle\quad\frac{1}{|\mathcal{I}_k|}\sum_{i\in\mathcal{I}_k}\mathbf{1}\{T_i=t\}\left(\left(\hat{\mu}_{(t)}^{-k}(X_i)-\mu_t(X_i)\right)\left(\frac{1}{\hat{e}_t^{-k}(X_i)}-\frac{1}{e_t(X_i)}\right)\right)\\
&\displaystyle\leq\sqrt{\frac{1}{|\mathcal{I}_k|}\sum_{i:i\in\mathcal{I}_k,T_i=t}\left(\hat{\mu}_{(t)}^{-k}(X_i)-\mu_t(X_i)\right)^2}\sqrt{\frac{1}{|\mathcal{I}_k|}\sum_{i:i\in\mathcal{I}_k,T_i=t}\left(\frac{1}{\hat{e}_t^{-k}(X_i)}-\frac{1}{e_t(X_i)}\right)^2}\\
&\displaystyle\leq_{(1)} \frac{\tilde{C}\epsilon_NN^{-1/2+\gamma}\psi_N}{[\inf_{x\in\mathcal{X}}\prod_{j=1}^M \alpha_{0j}(x)^{t_j}(1-\alpha_{0j}(x))^{1-t_j}]^{2}},\mbox{where (1) holds with probability $1-\epsilon_N$}
\end{array}$$
by Assumption~\ref{ass:causal:extra}. Thus $\sqrt{N}|(c)|\leq\frac{\tilde{C}\epsilon_NN^{\gamma}\psi_N}{[\inf_{x\in\mathcal{X}}\prod_{j=1}^M \alpha_{0j}(x)^{t_j}(1-\alpha_{0j}(x))^{1-t_j}]^2}$ with probability $1-\epsilon_N$, where $N^{\gamma}\psi_N=\mathrm{o}(1)$. Since $M=\mathrm{o}(\min\{\log(1/\psi_N),\log(1/\epsilon_N)\})$, so $\frac{\epsilon_N^2}{c^M}=\mathrm{o}(1), \frac{\psi_N}{c^{3M}}=\mathrm{o}(1), \frac{\epsilon_N}{c^{2M}}=\mathrm{o}(1)$, hence for any $k\in[K]$ we have $\sqrt{N}\left|\hat{\eta}_{(t)}^{\mathcal{I}_k}-\hat{\eta}_{(t)}^{\mathcal{I}_k,*}\right|=\mathrm{o}_p(1)$. Thus the result holds. 
\proofend

\end{document}
